\documentclass[11pt,letterpaper]{article}
\usepackage[T1]{fontenc}
\usepackage[utf8]{inputenc}
\usepackage{lmodern}
\usepackage[margin=1in,headheight=15pt]{geometry}
\usepackage{amsmath,amssymb,amsthm,mathtools,mathrsfs}
\usepackage{microtype,booktabs,longtable,tabularx,array,enumitem}
\usepackage[dvipsnames]{xcolor}
\usepackage{listings,fancyhdr}
\usepackage[colorlinks=true,linkcolor=MidnightBlue,citecolor=MidnightBlue,urlcolor=MidnightBlue]{hyperref}
\usepackage[nameinlink,noabbrev]{cleveref}
\hypersetup{%
  pdftitle={Surreal and Surcomplex Numbers in Symbolic Computer Algebra},
  pdfauthor={Merged research report},
  pdfsubject={Exact representation tiers, effective cores, certified supports, analytic and contour layers, branch and phase conventions, and three delivered Wolfram prototypes}}
\pagestyle{fancy}\fancyhf{}
\fancyhead[L]{\small Surreal and surcomplex computer algebra}
\fancyhead[R]{\small Representation, algorithms, and semantics}
\fancyfoot[C]{\thepage}
\setlength{\parskip}{0.28em}
\setlength{\parindent}{1.2em}
\setlist{leftmargin=2em,itemsep=0.15em,topsep=0.3em}
\setcounter{tocdepth}{2}
\setlength{\emergencystretch}{2em}
\numberwithin{equation}{section}
\newtheorem{theorem}{Theorem}[section]
\newtheorem{proposition}[theorem]{Proposition}
\newtheorem{lemma}[theorem]{Lemma}
\newtheorem{corollary}[theorem]{Corollary}
\theoremstyle{definition}
\newtheorem{definition}[theorem]{Definition}
\newtheorem{example}[theorem]{Example}
\theoremstyle{remark}
\newtheorem{remark}[theorem]{Remark}
\newcommand{\No}{\mathbf{No}}
\newcommand{\SC}{\mathbf{SC}}
\newcommand{\R}{\mathbb R}
\newcommand{\C}{\mathbb C}
\newcommand{\Q}{\mathbb Q}
\newcommand{\N}{\mathbb N}
\newcommand{\Z}{\mathbb Z}
\newcommand{\RA}{\mathbb R_{\mathrm{alg}}}
\newcommand{\OO}{\mathcal O}
\newcommand{\mm}{\mathfrak m}
\newcommand{\HH}{\mathscr H}
\newcommand{\Acal}{\mathscr A}
\newcommand{\st}{\operatorname{st}}
\newcommand{\fin}{\operatorname{fin}}
\newcommand{\supp}{\operatorname{supp}}
\newcommand{\lc}{\operatorname{lc}}
\newcommand{\sgn}{\operatorname{sgn}}
\newcommand{\cis}{\operatorname{cis}}
\newcommand{\Res}{\operatorname{Res}}
\newcommand{\res}{\operatorname{res}}
\newcommand{\NF}{\operatorname{NF}}
\newcommand{\Tr}{\operatorname{Tr}}
\newcommand{\Disc}{\operatorname{Disc}}
\newcommand{\id}{\operatorname{id}}
\newcommand{\vH}{v_{\!H}}
\newcommand{\Sin}{\operatorname{Sin}}
\newcommand{\Cos}{\operatorname{Cos}}
\newcommand{\ExpNo}{\operatorname{Exp}_{\No}}
\newcommand{\LogNo}{\operatorname{Log}_{\No}}
\newcommand{\Expi}{\operatorname{Exp}_{\mathrm{inf}}}
\newcommand{\Logi}{\operatorname{Log}_{\mathrm{inf}}}
\newcommand{\Ov}[1]{O_v(#1)}
\newcommand{\sh}{^{\sharp}}
\newcommand{\HInt}{\mathop{\int^{\!H}}\nolimits}
\DeclareMathOperator{\Arg}{Arg}
\DeclareMathOperator{\Logop}{Log}
\newcolumntype{P}[1]{>{\raggedright\arraybackslash}p{#1}}
\newcolumntype{L}[1]{>{\raggedright\arraybackslash}p{#1}}
\newcolumntype{Y}{>{\raggedright\arraybackslash}X}
\lstdefinestyle{code}{basicstyle=\ttfamily\small,columns=fullflexible,keepspaces=true,
  breaklines=true,showstringspaces=false,frame=single,framerule=0.25pt,
  rulecolor=\color{gray!45},backgroundcolor=\color{gray!3},
  xleftmargin=0.3em,xrightmargin=0.3em,aboveskip=0.7em,belowskip=0.7em}
\lstset{style=code}
\makeatletter
\g@addto@macro{\UrlBreaks}{\UrlOrds}
\makeatother
\renewcommand{\UrlFont}{\small\ttfamily}

\begin{document}
\hypersetup{pageanchor=false}
\begin{titlepage}
\centering
\vspace*{0.6cm}
{\large\scshape A merged research and implementation report\par}
\vspace{0.9cm}
{\Huge\bfseries Surreal and Surcomplex Numbers\par}
\vspace{0.4cm}
{\LARGE in Symbolic Computer Algebra\par}
\vspace{0.7cm}
{\large Exact Representation Tiers, Effective Cores, Certified Supports,\par
Analytic and Contour Layers, Branch and Phase Conventions,\par
and Three Delivered Wolfram Prototypes\par}
\vspace{0.9cm}
{\large September 21, 2026\par}
\vspace{0.7cm}
\begin{minipage}{0.93\textwidth}
\begin{abstract}
A symbolic computer algebra system can support large and useful parts of the
surreal field $\No$ and its complexification $\SC=\No[i]$ without attempting to
store an arbitrary surreal number or to decide an arbitrary identity. The right
product is not one universal numeric type but a layered family of exact domains
with separately advertised, independently queryable capability contracts: an
aggressively decidable algebraic core surrounded by explicitly partial
extensions. Throughout we separate exact denotation, canonical representation,
effective coefficient access, decidable equality and order, certified
approximation, and proof-carrying evidence; these are six different promises,
and a system that conflates them will overstate what it can do.

We fix an effective rational-monomial core and prove that it is an effective
ordered field with terminating arithmetic, equality, valuation, leading
coefficient and real sign, and we record the three different specializations of
that core that were actually implemented. We prove three negative results ---
no total zero test for arbitrary computable coefficient programs (by two
genuinely different reductions, whose morals differ), no total support-admissibility
validator, and no universal leading-term or sign algorithm for high-rank
program-described grid series even under a nonzero promise --- together with one
positive counterweight beyond the algebraic core, the 2026 D-algebraic
transseries zero-test, carried with its hypotheses intact. A constructive
finite-grid dependency lemma, given here in two independent constructions of its
auxiliary positive functional, separates effective coefficient access from finite
truncation at a requested valuation; a positive rank-one criterion and three
counterexamples of two different kinds delimit that separation.

On the analytic side we develop exact functions in four separate categories:
formal germs, ordinary holomorphic lifts, Hahn-coherent data on a fixed common
ordinary domain, and the strictly weaker radius-free Hahn-germ algebra, in which
coefficients share no ordinary radius. Polynomial algebra receives initial
polynomials, Newton profiles, Hensel cluster lifting and two separate sharp
root-perturbation theorems. Finite quotient algebras receive residue and trace
pairings, multiplication matrices, a perturbed residue formula and a worked
two-scale coupled system. Contour operations receive the coefficientwise Hahn
functional, algebraic rational periods, microscopic rescaling, a separating
torus and a certificate-based Stokes rewrite. Surcomplex angles receive the
finite-angle circle, an exact cut-free phase, the Cayley coordinate, the
canonical strips, two logarithm branches, and infinite circular phases under
two distinct and separately attributed accounts of nonuniqueness.

Wolfram Language facilities are assessed facility by facility against official
documentation consulted on September 21, 2026, and a set of host-integration
prohibitions is stated. A proposed object architecture is given in
proposal voice only; nothing named in it exists in any delivered file.
\textbf{Three} separately delimited Wolfram Language prototypes are delivered:
\texttt{RationalHahn} over $\Q$, \texttt{SurrealCASCore} over $\Q(i)$ with
rational exponents, and \texttt{HahnRational} over $\Q(i)$ with integer
exponents. Their envelopes differ on every axis that matters, two of them export
six identical public symbol names with incompatible signatures, and no sentence
in this report asserts the union of their capabilities.
\end{abstract}
\end{minipage}
\vfill
\begin{minipage}{0.93\textwidth}\small
\textbf{Status and provenance.} This is an implementation-oriented mathematical
exposition merged from three independently written source articles. Six distinct
supplied research manuscripts stand behind it, one of them in two different
revisions; they are treated as sources, not as independently refereed or
machine-verified publications. Established theorems, deductions proved in the
sources, results imported from the supplied manuscripts, proposed interfaces,
and executed example tests are distinguished throughout. Nothing here is
machine-verified: no Lean file, formalization blueprint or formalization audit
exists in any of the three source packages, and no git commit, Lean toolchain or
mathlib revision is pinned anywhere. No universal surreal simplifier, complete
surcomplex analytic system, or complete surreal CAS is claimed.
\end{minipage}
\end{titlepage}
\hypersetup{pageanchor=true}
\pagenumbering{roman}
\begingroup\footnotesize\setlength{\parskip}{0pt}\tableofcontents\endgroup
\clearpage
\pagenumbering{arabic}

\section{The answer: a layered exact system with queryable capabilities}
\label{cas:sec-answer}

\subsection{What should be supported first}
\label{cas:sub-first}
A useful implementation should begin with an exact ordered field of rational
expressions in explicitly ordered surreal monomials, extend it by algebraic root
objects, and attach certified series representations when expansions are wanted.
This already permits genuine infinite and infinitesimal quantities rather than a
symbolic placeholder for a limit. Complexifying the real field gives an exact
surcomplex coefficient domain, after which polynomial division, greatest common
divisors, resultants, finite linear algebra and polynomial root specifications
become natural operations.

The following calculations should all be straightforward in an initial system:
\begin{align*}
 (\omega+1)^2-\omega^2-2\omega&=1,\\
 (\omega+1)^{-1}&=\frac{t}{1+t},\qquad t=\omega^{-1},\\
 \sqrt{\omega^2+1}&=\omega\sqrt{1+t^2},\\
 \arctan\omega&=\frac\pi2-\arctan t.
\end{align*}
The second equality is already an exact answer; expanding it into infinitely
many monomials is optional. The third requires an algebraic extension, and the
fourth requires the canonical finite-angle analytic layer. A system should be
able to explain these changes of domain rather than perform them silently.

A second layer should represent exact infinite series by finite rules with
proved support conditions: evaluate Taylor series at infinitesimals, compute
requested coefficients, and, where justified, return finite truncations with
valuation error bounds. A third layer can add selected exponential--logarithmic
and transserial constructors, and ordinary holomorphic functions lifted to
finite surcomplex arguments. These are extensions with additional semantics, not
consequences of adding one arithmetic head.

The design has a deliberately asymmetric policy. Exact algebra is aggressive
inside a declared decidable domain; more general analytic expressions remain
symbolic until their domain, support and branch obligations have been
discharged. An unevaluated but well-defined expression is a legitimate exact
result. An expression with an unproved summability obligation is not yet an
admitted exact number. The system should never silently replace an unsupported
value by zero, by an ordinary infinity, by a large real approximation, or by a
convention-dependent phase.

Three questions guide every interface below. What mathematical object does the
expression denote? What operations can the representation carry out with
guaranteed termination? What information is retained when the system returns a
finite approximation? Answering the first does not settle the other two.

\subsection{Six promises, and the four-way version two of the sources state}
\label{cas:sub-promises}
The organizing distinction of this report is that ``exact'' names several
independent capabilities that must be queried separately. We use the six-way
list, which is a strict refinement of the four-way lists given independently by
two of the three source articles. The refinement splits out
\emph{canonicality} --- equal values have identical normalized data --- and
\emph{proof-carrying representation}, neither of which the four-way lists name.
Canonicality is not a cosmetic addition: it is exactly the property that
distinguishes the canonical sparse layer of one delivered prototype from that
same prototype's deliberately unreduced fraction wrapper
(\cref{cas:sec-prototypes}).

\begingroup\small
\begin{longtable}{@{}P{0.22\textwidth}P{0.35\textwidth}P{0.36\textwidth}@{}}
\toprule
Promise & What it promises & What it does not promise\\
\midrule\endhead
Exact denotation & A finite description singles out one mathematical value,
without rounding. & A normal form, a numerical approximation, or a terminating
equality test.\\[0.45em]
Canonical representation & Equal values have identical normalized data. &
That the normalization is cheap, or that it exists for every admitted
representation.\\[0.45em]
Effective coefficient access & A requested coefficient can be calculated by a
terminating procedure. & That the first nonzero coefficient can be found, or
that all coefficients below a cutoff can be listed.\\[0.45em]
Decidable algebra & Equality and the required field operations terminate on all
valid inputs in a declared class. & Closure under every elementary or special
function; acceptable complexity.\\[0.45em]
Certified approximation & The answer includes a proved error specification in a
named valuation or topology. & Equality to the truncated expression, or
approximation in the full surreal fine topology.\\[0.45em]
Proof-carrying representation & Checkable evidence is supplied for the
conditions that make the expression meaningful. & A machine-checked
formalization; the evidence is checked against trusted constructors.\\
\bottomrule
\end{longtable}
\endgroup

Two of the three source articles present a four-way version of this taxonomy:
denotationally exact / coefficient-effective relative to an indexed support
description / decision-effective for a named operation set / certified
truncation in one case, and exact denotation / effective coefficient access /
decidable algebra / certified approximation in the other. Both insist, as we do,
that the levels are independent and must be queried independently. We do not
print ``there are four levels of exactness'' as the merged claim, because that
would silently drop two capabilities.

\subsection{The running illustration}
\label{cas:sub-abc}
Consider three outputs:
\[
 A=\frac1{1-t},\qquad
 B=1+t+\cdots+t^9+\Ov{t^{10}},\qquad
 C=\operatorname{StrongSum}(n\mapsto t^n).
\]
Here $A$ names one exact element; it is a finite exact field expression whose
Hahn normal form has infinitely many terms. $B$ names a \emph{precision class}
unless it carries an additional exact-tail specification; the displayed
polynomial $1+t+\cdots+t^9$ is not the same exact number as $A$. And $C$ names
the same element as $A$ only after the system has certified strong summability
and the intended monomial interpretation. Printing an infinite summation symbol,
or leaving a host expression inactive, supplies no certificate. Conversely, the
program description of a series can be exact while equality of two such
programs' values is undecidable (\cref{cas:thm-zerotest}). These are structural
distinctions, not implementation inconveniences.

\subsection{An initial capability map}
\label{cas:sub-capmap}
The following map is the union of the two presentation taxonomies given by the
sources. Its rows are not a claim that every object in an abstract Hahn field or
every analytic germ is finitely describable; they identify workable
\emph{presentation families} inside the mathematical objects.

\begingroup\small
\begin{longtable}{@{}L{0.22\textwidth}L{0.34\textwidth}L{0.35\textwidth}@{}}
\toprule
Representation & Exact content available & Principal computational boundary\\
\midrule\endhead
Hereditarily finite cuts and sign data & Dyadic rationals; birthday and
simplicity data & Not closed under division; recursive trees grow; $1/3$ is
already outside.\\[0.4em]
Finite Hahn sums / normal forms & Finite sums of admitted monomials; nested
finite scales & Closed under $+$ and $\times$; reciprocals generally need
another representation.\\[0.4em]
Ordered rational monomial expressions & A non-Archimedean fraction field &
Total arithmetic, equality and real order for effective parents; algebraic
independence and the monomial order must be fixed; not algebraically
closed.\\[0.4em]
Ordered algebraic data & Real closure, complex polynomial roots & Needs
non-Archimedean root selectors, not floating-point isolation.\\[0.4em]
Certified series grammars / lazy Hahn series & Exact infinite sums, local
Taylor values, positive grids & Support validity, effective zero and
leading-term access are separate services.\\[0.4em]
Exp--log and transserial graphs & Infinite exponential scales, asymptotic
composition & Partial normalization; the input class of an actual transseries
algorithm is narrower than ``all transseries''.\\[0.4em]
Analytic germs and Hahn families & Local functions, finite jets, coherent
periods & Domain, support, branch and ordinary coefficient algorithms must all
be supplied.\\[0.4em]
General oracular or program-described Hahn data & A broad countable collection
of exact values & No universal equality, support validation or finite cutoff
materialization.\\
\bottomrule
\end{longtable}
\endgroup

\subsection{Provenance: six supplied manuscripts, one of them in two revisions}
\label{cas:sub-provenance}
This report merges three independently written source articles. They were given
different supplied research manuscripts, so the phrase ``the supplied
manuscript'' is ambiguous across the group and is never used in the singular
here. The six distinct manuscripts are:
\begin{itemize}
\item \emph{Surcomplex Analysis: Infinitesimal Calculus, Hahn-Coherent
Holomorphy, and Contour Theory}, in \textbf{two revisions}. Revision~(2)
\cite{cas:SourceAnalysisTwo} was supplied to the capability-tiers article;
revision~(3) \cite{cas:SourceAnalysisThree} was supplied to the other two. We
write \textbf{SA} when the distinction does not matter and name the revision
when it does.
\item \emph{Trigonometry on the Surcomplex Plane: Canonical Angles, Infinite
Triangles, Infinitesimal Contact, and Hahn-Analytic Oscillation}
\cite{cas:SourceTrig}, supplied only to the capability-tiers article. We write
\textbf{TR}.
\item \emph{Surcomplex Polynomial Algebra: Order Geometry, Newton Profiles, and
Scale-Resolved Stability} \cite{cas:SourcePolynomial}, supplied only to the
layered-representations article. We write \textbf{PA}.
\item \emph{Hartogs Coherence, Finite Maps, and Residue Duality over the
Surcomplex Numbers} \cite{cas:SourceResearch}, supplied only to the
layered-representations article. We write \textbf{HC}.
\item \emph{Infinitesimal Analytic Geometry over the Surcomplex Numbers:
Radius-Free Hahn Germs, a Nullstellensatz, and Conservation of Singular Zeros}
\cite{cas:SourceAG}, supplied only to the Gaussian-rational-prototype article.
We write \textbf{AG}.
\item \emph{Jordan Separation, Cauchy Theory, and Stokes Formulae on the
Surcomplex Plane: A Category-Sensitive Contour Calculus} \cite{cas:SourceCT},
supplied only to the Gaussian-rational-prototype article. We write \textbf{CT}.
\end{itemize}
Every result imported below is attributed to the specific manuscript, and where
relevant to the specific revision, that supplies it.

Only one of the three delivered archives bundled its sources: the
capability-tiers archive shipped unchanged copies of SA~(2) and TR under its own
\texttt{sources/} directory, so only its provenance is independently inspectable
from the delivery itself. Those two files are already committed in this
repository and are not re-shipped here: the trigonometry manuscript is
byte-identical to
\begin{sloppypar}\noindent
\path{docs/surcomplex/trigonometry/sources/18-canonical-angles-oscillation.tex},
and SA~(2) is identical to
\path{docs/surcomplex/analysis/sources/07-infinitesimal-calculus-contour-theory.tex}
after CRLF normalization.
\end{sloppypar} The other two archives cite their manuscripts rather
than bundling them.

Two distinct source-gap disclaimers survive from the sources and are both
honest limitations. First, two further archives mentioned \emph{inside} TR were
not supplied, are not included, are not treated as independently reviewed
sources, and their own referenced verification scripts were not rerun. Second, a
fourth, global-divisor manuscript cited \emph{inside} PA is not among the three
files supplied to that article and is not used here as an independently
inspected source.

Finally, and once for the whole report: nothing here is machine-verified. None
of the three source packages contains a Lean file, a formalization blueprint, or
a formalization audit of any kind, and no git commit hash, Lean toolchain or
mathlib revision is pinned anywhere in any of them. The finite symbolic checks
recorded in \cref{cas:sec-executed} are exactly that.

\section{Objects the software must distinguish}
\label{cas:sec-objects}

\subsection{Normal forms, workspaces, and set-sized parents}
\label{cas:sub-workspaces}
Following all three sources and all six supplied manuscripts we write
\begin{equation}\label{cas:eq-normal}
 t^\gamma:=\omega^{-\gamma},\qquad
 A=\sum_{\gamma\in S}a_\gamma t^\gamma,\qquad
 S\subseteq\Gamma\text{ well ordered},
\end{equation}
where the displayed power of $\omega$ is the \emph{Conway monomial map} and
$\Gamma$ is a set-sized ordered additive subgroup of $\No$. Put
$F_\Gamma=\R((t^\Gamma))$ and $K_\Gamma=\C((t^\Gamma))$. For nonzero $A$,
\[
 v(A)=\min\supp(A),\qquad \lc(A)=a_{v(A)},\qquad v(0)=+\infty,
\]
so $v(AB)=v(A)+v(B)$ and $v(A+B)\ge\min\{v(A),v(B)\}$. For a real series the
sign is the sign of the leading coefficient. When $\Gamma$ is divisible the real
Hahn field is real closed and the complex one is algebraically closed; this
algebraic-closure input is a theorem about the complete mathematical field, not
a terminating algorithm for arbitrary coded series
\cite[Corollary 4]{cas:Poonen}.

This is a semantic workspace, not a finite data structure. Even $\Q((t))$
contains uncountably many arbitrary rational coefficient sequences. A practical
implementation replaces $\R$ or $\C$ by a declared coefficient representation,
replaces arbitrary exponents by an effective ordered-group representation, and
admits only selected finitely described supports and coefficient rules. Such a
replacement preserves some closure properties and loses others; every closure
claim must refer to the replacement, not to the surrounding Hahn field. Every
collection of coefficients, every support and every family summed in this report
is a set: $\No$ is a proper class, but every individual expression is finite
data interpreted in suitable set-sized workspaces. A class-sized theory may
justify those parents without ever allocating a class-sized list, ring or
support.

\subsection{Conway monomials are not generic powers}
\label{cas:sub-conway}
The symbol $\omega^a$ in a Conway normal form denotes the omega-map. It satisfies
$\omega^{a+b}=\omega^a\omega^b$, but it is \emph{not} defined by the generic
prescription $\ExpNo(a\LogNo\omega)$: the canonical real surreal exponential and
logarithm are an additional structure \cite{cas:Gonshor,cas:vdDE}. The
distinction matters as soon as the exponent is surreal but not an ordinary real.
Concretely,
\[
 \omega^{1/\omega}\ \text{(Conway notation) is \emph{infinite}},
 \qquad\text{whereas}
\]
\[
 \ExpNo\!\left(\frac{\LogNo\omega}{\omega}\right)
 =1+\text{a positive infinitesimal}.
\]
The first has positive exponent, hence an Archimedean class strictly above that
of $1$; in the second, $\LogNo\omega/\omega$ is infinitesimal, by the growth
properties of the surreal exponential. Rational powers of $\omega$ cause no such
conflict: both constructions give the unique appropriate positive algebraic
root. The dangerous step is extending that rational-power simplification to
arbitrary surreal exponents.

An interface should therefore distinguish at least three constructors,
\[
 \texttt{ConwayMonomial[a]},\qquad \texttt{SurrealPower[x,y]},\qquad
 \texttt{OrdinalPower[a,b]},
\]
where for positive surreal $x$ the second may mean $\ExpNo(y\LogNo x)$ and the
third is ordinal exponentiation. Identities relating them must be proved on
their particular domains, and the two constructions need an actual compatibility
theorem on the chosen argument class rather than an inherited notation. A
printed superscript is not sufficient semantic metadata. General complex powers
need still more: a chosen logarithm or a specified algebraic-root branch
(\cref{cas:sub-powers}).

\subsection{Ordinal arithmetic, field arithmetic, and the host's infinity}
\label{cas:sub-ordinals}
Ordinary ordinal addition is not surreal field addition:
$1+_{\mathrm{ord}}\omega=\omega$, whereas $1+\omega=\omega+1>\omega$ in the
field. On embedded ordinals, surreal field arithmetic agrees with natural
(Hessenberg) arithmetic, not with an unrestricted replacement of ordinal
operations by field operations \cite{cas:Conway}. A package supporting both
needs distinct constructors, and birthday or simplicity queries should be a
separate capability from field equality.

A genuine infinite field element is also not a projective point at infinity and
not the host's \texttt{Infinity}:
\[
 \omega-\omega=0,\qquad \omega/\omega=1,\qquad \omega^{-1}>0.
\]
A head used for a pole or an extended limit cannot supply those field laws;
\texttt{Infinity} participates in extended-limit arithmetic and is not a named
element of a non-Archimedean ordered field \cite{cas:WLInfinity}. In all three
delivered prototypes \texttt{Infinity} appears only as the conventional metadata
value $v(0)=+\infty$, never as a stored field element.

Nor is a surreal infinitesimal a nilpotent: $t^n\ne0$ for every ordinary $n$.
See \cref{cas:sub-fourinf}.

\subsection{Order, valuation, standard part, and the finite part}
\label{cas:sub-st}
With ordinary Archimedean coefficients, $v(z)\ge0$ means $z$ is finite and
$v(z)>0$ means it is infinitesimal. Let $\OO$ denote the finite elements and
$\mm$ the infinitesimals. Every $z\in\OO$ decomposes uniquely as $z=c+\eta$ with
$c$ ordinary and $\eta\in\mm$; the standard part $\st(z)=c$ is a ring
homomorphism \emph{on this finite subring}. The zero coefficient of an infinite
element can also be extracted, but calling it a standard part would misstate the
domain of the residue homomorphism.

The \emph{finite part} $\fin$ used by TR is a different projection: it removes
all negative-exponent terms from any normal form. It is additive but not
multiplicative, since
\[
 \fin(\omega\cdot\omega^{-1})=1\ne\fin(\omega)\fin(\omega^{-1})=0 .
\]
These operations require different names and different rewrite rules. Even when
a standard part exists mathematically, extracting it from an arbitrary
program-defined series may not be computable.

For a surcomplex $z=x+iy$, conjugation and modulus are $\bar z=x-iy$ and
$|z|=\sqrt{x^2+y^2}$. The modulus satisfies the ordinary triangle inequality
while the valuation satisfies an ultrametric inequality: an order disk
$|z-a|<r$ is not a valuation ball $v(z-a)>v(r)$ \cite{cas:SourcePolynomial}.

Software must distinguish an exact valuation from a lower bound. If the leading
displayed coefficient is not known to be nonzero, the valuation has not been
established. Cancellation in a sum is not cosmetic: it can change sign, scale,
and the number of roots in a selected ball. The delivered kernels make this
concrete; see the cancellation trap in \cref{cas:sub-cancellationtrap}.

\subsection{Four incompatible uses of ``infinitesimal''}
\label{cas:sub-fourinf}
A nonzero surreal infinitesimal has an inverse and is not nilpotent. A dual
number $a+b\epsilon$ with $\epsilon^2=0$ belongs to a different algebra. A
truncation modulo $t^N$ belongs to a quotient ring with nilpotents. These rings
are useful for jets and automatic differentiation, but neither is a subfield of
the surreals: storing a surreal $t$ in a dual-number object and applying its
multiplication law asserts the false equality $t^2=0$.

An \emph{asymptotic variable tending to zero} is a fourth kind of object. It is
a variable in a function or germ, not yet a particular surreal constant.
Converting its formal expansion into an exact Hahn value requires a selected
interpretation and a summability proof. We retain all four distinctions even
where one displayed expression is convenient for all four purposes.

\subsection{A type is more than a coefficient container}
\label{cas:sub-type}
An exact object should carry or reference its coefficient domain, its exponent
domain and order, its embedding into the surreal numbers, its analytic
interpretation, and any branch or derivation convention. A formal generator
called $s$ is not enough: $s=t^2$ and $s\prec t^n$ for every ordinary $n>0$ give
different fields and different equality tests. Domain identifiers should be
structural and serializable rather than dependent on session-local assumptions
that can later change. Ordinary substitution can destroy the intended
relationships: replacing $u=t^\omega$ by $t^{100}$ collapses a rank, and
\textbf{a different variable order is a different parent even when the printed
rational expression is identical}. That last point is not hypothetical here ---
two of the three delivered prototypes use opposite coordinate significance
conventions (\cref{cas:sub-orderconventions}).

\section{Limits of finite representation and universal simplification}
\label{cas:sec-limits}

\subsection{There is no finite notation for every surreal}
\label{cas:sub-countable}
\begin{proposition}[The finite-description bound]\label{cas:prop-countable}
For any fixed finite or countable alphabet, each symbol itself specified by
finite data, a language of finite strings can denote at most countably many
individually specified numbers. It cannot denote every real number, and
therefore cannot denote every surreal number.
\end{proposition}
\begin{proof}
There are countably many strings of each finite length and countably many
lengths, so countably many strings in all. The image of this countable set under
a partial denotation map cannot have larger cardinality. The reals are
uncountable and are contained in $\No$.
\end{proof}

All three sources prove or state this independently, and draw the same three
consequences, which we record as remarks rather than as three propositions.

\begin{remark}[Oracles move the information, they do not encode it]
\label{cas:rem-oracle}
A symbol for an arbitrary unspecified real parameter gives a symbolic family,
not a finite encoding of every possible value of that parameter. An oracle or an
opaque atom can stand for an arbitrary real or surreal, but the missing
information has then been moved outside the finite input rather than encoded in
ordinary finite memory; the resulting object is exact \emph{relative to the
oracle}. A symbol $c$ carrying assumptions names a parameter or a family until a
particular value is specified.
\end{remark}

\begin{remark}[Set-sized support is not a computability statement]
\label{cas:rem-setsized}
``Every normal form has set-sized support'' is a statement about cardinality and
order type, not about computation. A set may be uncountable, noncomputable, or
lack any finite structural description. Likewise, a proof assistant can quantify
over all surreals without making them machine-representable data, and a computer
algebra system can state universal identities in abstract variables without
materializing their values.
\end{remark}

\begin{remark}[No largest representable subset]\label{cas:rem-nolargest}
There is no uniquely largest useful finitely representable subset independent of
a chosen language. Adding one new constant, one new support notation, or one new
certified constructor may strictly extend an existing language, and each
enlargement changes the available decision procedures. Even two representations
of the \emph{same} mathematical subset can have very different computational
properties --- which is exactly the content of \cref{cas:thm-zerotest}. The
practical question is which representations offer a useful and honest collection
of operations.
\end{remark}

\subsection{The zero-test obstruction: one theorem, two reductions}
\label{cas:sub-zerotest}
\begin{theorem}[Total coefficient access does not decide zero]
\label{cas:thm-zerotest}
There is no algorithm which, for every total program specifying rational
coefficients $a_n$ of a valid Hahn series $A=\sum_{n\ge0}a_nt^n$, decides
whether $A=0$. Consequently no universal terminating equality or three-way sign
procedure exists for this unrestricted representation class.
\end{theorem}

Two of the three sources reduce from halting in one way and the third in
another. The two reductions have different punchlines, both load-bearing, so we
print both.

\begin{proof}[Reduction A: values in a decidable rational field]
For a machine $M$ and its input, let $a_n=1$ if $M$ has halted \emph{within} $n$
steps, and $a_n=0$ otherwise. Computing $a_n$ needs only bounded simulation, so
the sequence is total computable, and its support lies in $\N$, making the
displayed series a valid Hahn sum at $t=\omega^{-1}$. It is zero exactly when
$M$ never halts. A total zero test would therefore decide halting. All
coefficients are nonnegative, so a sign procedure distinguishing zero from
positive would do the same; equality reduces to zero testing by subtraction.
\end{proof}

\begin{remark}[Moral of Reduction A]\label{cas:rem-moralA}
In this reduction every individual value is either $0$ or the rational function
$t^N/(1-t)$ for some halting time $N$. Thus even values lying in a
\emph{decidable rational field} become undecidable when presented by
unrestricted coefficient programs. \emph{Representation, not merely membership
in an abstract field, determines the algorithmic guarantee.} The whole
capability-contract thesis of this report rests on that sentence.
\end{remark}

\begin{proof}[Reduction B: a promise of at most one nonzero coefficient]
For a machine $M$ and its input, let $b_M(n)=1$ if $M$ \emph{first} halts at
step $n$ (for $n\ge1$) and $b_M(n)=0$ otherwise; this is computable by
simulating exactly $n$ steps. Then $a_M=\sum_{n\ge1}b_M(n)t^n$ has support
either empty or a singleton, and is zero precisely when $M$ never halts.
\end{proof}

\begin{remark}[Moral of Reduction B]\label{cas:rem-moralB}
Here every input in the reduction actually denotes either zero or a single
monomial --- indeed a polynomial. The obstruction therefore survives even when
\emph{every represented series is promised to have at most one nonzero
coefficient}. This is the sharper form of the same point: an explicit finite
polynomial list is different information from a program that might reveal one
nonzero coefficient arbitrarily late. There is no tension with
\cref{cas:thm-core}: opaque programs promised merely to denote finite sparse
polynomials need not have a computable support bound, and can describe exactly
the same set of values with a much weaker computational contract.
\end{remark}

Neither reduction is ``the'' proof; both survive.

\begin{corollary}[No universal three-way sign; equality reduces to zero testing]
\label{cas:cor-sign}
There is no universal terminating three-way sign procedure for this
representation class, and equality of two such presentations reduces to zero
testing of their difference.
\end{corollary}

\begin{corollary}[No universal leading-term algorithm]\label{cas:cor-leading}
A leading-term algorithm that always either reports zero or returns the least
nonzero exponent would also decide zero. No such algorithm exists for arbitrary
total coefficient programs.
\end{corollary}

All three sources insist independently on the same operational consequence.
\textbf{Returning \texttt{False} after a fixed unsuccessful search is unsound.}
The correct result is an explicit unknown status or a resource-limited
inconclusive result. A lazy series implementation remains useful, but must admit
nontermination, resource-bounded \texttt{Unknown}, or restricted input grammars
carrying extra zero certificates. Numerical sampling does not repair the
problem. Sage's lazy-series documentation makes the related practical
distinction between coefficient generation and equality decidable only in
special cases \cite{cas:SageLazy}.

The obstruction requires neither arbitrary real coefficients nor complicated
surreal exponents: a certified support container $\N$ is available from the
outset, and the difficulty is locating the first nonzero coefficient or proving
none exists. It also does not prevent a zero test for rational generating
functions, algebraic series with finite defining data, or other specifically
certified families (\cref{cas:sub-positive}). The promise must be attached to
the representation class rather than to the word ``exact''.

\subsection{Support admissibility is not a universal automatic test}
\label{cas:sub-validation}
\begin{proposition}[Unrestricted support validation is undecidable]
\label{cas:prop-validation}
No algorithm decides Hahn admissibility for every program which outputs a
candidate support using total bounded-time coefficient tests.
\end{proposition}
\begin{proof}
For a machine $M$, assign at exponent $-n$ (for $n\ge1$) the coefficient $1$ if
$M$ has halted by step $n$ and $0$ otherwise. Every individual test terminates.
If $M$ never halts the exact support is empty and admissible. If $M$ halts, the
support contains an infinite strictly decreasing tail $-n,-(n+1),\dots$ and is
not well ordered. Deciding admissibility would decide halting.
\end{proof}
The remedy is not to omit validation. It is to use constructors whose support
properties follow from their syntax or from checked certificates, and to treat
unrestricted user generators as untrusted until obligations are supplied. A
validator can be total and complete for a restricted grammar; it cannot be both
for all programs of the preceding proposition.

\subsection{Known nonzero does not make the leading term computable}
\label{cas:sub-highrank}
\begin{proposition}[A high-rank leading-term and sign obstruction]
\label{cas:prop-leading}
For total program-described two-scale grid series, even a promise that the value
is nonzero does not give a universal algorithm for its exact valuation or its
sign.
\end{proposition}
\begin{proof}
Use exponents in $\Q^2_{\mathrm{lex}}$ and let $a_n$ be the bounded-simulation
coefficients of Reduction~A, with halting steps indexed from $1$. Put
\[
 B_M=t^{(1,0)}-\sum_{n\ge1}a_nt^{(0,n)} .
\]
This is \emph{always nonzero} and actually has one or two terms. If $M$ never
halts it is positive with valuation $(1,0)$; if $M$ halts at step $n$ it is
negative with valuation $(0,n)$. Either a sign algorithm or an exact-valuation
algorithm would decide halting.
\end{proof}

\begin{remark}[The rank-one contrast, which must not be generalized]
\label{cas:rem-rankone}
There is a genuine contrast with a one-variable power series carrying a known
integer lower exponent bound. Under a nonzero promise, sequentially computing
its coefficients eventually finds the first nonzero one, because only finitely
many exponents precede any candidate. In higher rank, infinitely many potential
smaller exponents can precede a known nonzero coefficient. An implementation
must not generalize the rank-one search argument without checking this.
\end{remark}

\subsection{Positive counterweights}
\label{cas:sub-positive}
The negative results above delimit \emph{unrestricted} representations. Rich
structured languages escape them, and three positive routes should be recorded
explicitly so that the picture is not read as uniformly negative.

First, rational and algebraic representations already avoid the obstruction:
equality in the effective core is decided by the zero test on a single
polynomial (\cref{cas:thm-core}), and algebraic elements with finite defining
data and valid root selectors have decidable equality and order
(\cref{cas:thm-realclosure}).

Second, for an ordinary-point solution of a scalar linear differential equation
with analytic coefficients and nonvanishing highest-order coefficient, finitely
many initial derivatives determine the solution. With effective coefficient
arithmetic and the equation recorded, vanishing of the requisite initial data is
a finite zero certificate. Within a representation that can construct an
annihilating equation for a difference, identity reduces to a finite
initial-data test. Products and derivatives are handled in the usual D-finite
setting; reciprocals and arbitrary composition need not remain D-finite.
Singular points require a certified recurrence start and treatment of
exceptional indices: the phrase ``satisfies a differential equation'' alone does
not guarantee uniqueness. Wolfram's \texttt{DifferentialRoot} is a relevant
ordinary-function precedent, not an automatic Hahn-series implementation
\cite{cas:WLDifferentialRoot}.

Third, and most significant, is a published positive decidability result beyond
the algebraic core. Chen, Fang and van der Hoeven's \emph{A Zero-Test for
D-Algebraic Transseries} appeared in the ISSAC 2026 proceedings,
pp.~105--113, published July 12, 2026 \cite{cas:CFH}. Their framework uses
effective differential ambient fields and structured grid-based transseries;
their Theorem 18 extends the ambient field, and its zero-test, by a
distinguished solution of an appropriate quasi-linear differential equation,
with Ritt reduction and separation bounds central to the argument.

\begin{remark}[Scope of the 2026 zero-test]\label{cas:rem-cfh}
This is a concrete route beyond elementary expressions for structured
differential-algebraic extensions. It is \textbf{not} a zero-test for arbitrary
series-producing programs. A CAS implementing it must preserve its
effective-field, transbasis and solution-selection assumptions rather than
replacing them with an unrestricted \texttt{DifferentialRoot} label.
\end{remark}

\subsection{Richardson's undecidability and its restrictions}
\label{cas:sub-richardson}
Richardson supplies further undecidability results for sufficiently rich
expression languages of real functions \cite{cas:Richardson}. The restrictions
of that particular theorem matter, and it must not be inflated into a blanket
claim that every restricted exponential-constant identity problem is known to be
undecidable. The elementary coefficient-program reductions of
\cref{cas:sub-zerotest} are enough for the unrestricted Hahn-generator
interface considered here, and they are independent of difficult identities
among elementary constants.

\section{A taxonomy of exact representations}
\label{cas:sec-representations}

\subsection{Cuts, sign expansions, and the simplest-number semantics}
\label{cas:sub-cuts}
The original cut presentation and the sign-sequence presentation retain
simplicity and birthday information. Hereditarily finite legal cut trees,
equivalently finite sign expansions, represent exactly the dyadic rationals.
They are excellent for teaching the recursive construction and for checking
small identities, but a poor universal arithmetic storage format: $1/3$ has a
short rational expression and is not a finite-birthday surreal, so division
immediately leaves this representation.

A legal number cut has every left option smaller than every right option, and
its value is the \emph{simplest} number between the options, not an arbitrary
point in an interval. For instance $\{0\mid4\}=1$, not the midpoint $2$. A cut
constructor therefore cannot be implemented by choosing a convenient solution of
the inequalities \cite{cas:Gonshor}.

Compressed infinite cuts and ordinal-indexed sign blocks can be exact --- the
familiar description of $\omega$ by all natural numbers on the left is the
standard example --- and can name real numbers with effective digit
descriptions and much more. Their effectiveness depends on the language used for
the options, not on the braces. The compression language must say how a
transfinite block is indexed, how its legality is certified, and which
comparisons terminate; an ordinal notation appearing in an index is not a
universal algorithm for all ordinal or surreal questions. Computing the simplest
number between arbitrary program-defined classes of options is not a generic
finite operation. A cut constructor should record an ordinal or
well-foundedness certificate and an inequality certificate rather than accept an
arbitrary predicate and promise a canonical normal form.

Neither birthday nor simplicity is determined by the field structure alone. A
Hahn-field backend that knows addition and order should not pretend to supply
shortest sign expansions or the simplest number between arbitrary cuts.
Birthdays and simplicity comparisons belong in an optional capability, not among
the prerequisites for ordinary field arithmetic.

\subsection{Finite normal forms}
\label{cas:sub-finitenormal}
A sparse finite sum
\[
 A=c_1t^{\gamma_1}+\cdots+c_mt^{\gamma_m},\qquad \gamma_1<\cdots<\gamma_m,
\]
is a good exact representation once coefficient and exponent equality are
decidable. Addition merges supports; multiplication forms pairwise exponent sums
and collects equal terms; conjugation acts coefficientwise; for real
coefficients, ordering is a leading-term calculation. Recursive finite normal
forms for the exponents extend the hierarchy: a coefficient/exponent tree can
describe nested scales without materializing a sign sequence.

This representation is generally a ring, not a field:
$(1-t)^{-1}=\sum_{n\ge0}t^n$ is not a finite-support normal form. An
implementation must either move to rational-expression storage, create a lazy
infinite sum with a support certificate, or return an explicit
unsupported-domain result. Declaring the inverse to have a finite normal form
would be incorrect. If the coefficient zero test is only partial, a sparse map
may contain unresolved terms and cannot honestly advertise its first stored key
as the valuation.

\subsection{Rational expressions}
\label{cas:sub-rationalrep}
Fractions of finite monomial polynomials are especially effective: their
arithmetic can often be performed with no series expansion at all. If numerator
and denominator are sparse polynomials in independent monomial generators,
equality is decided by cross multiplication, and the result is an exact field
element even when no finite normal form exists. A common-denominator normal form
and polynomial gcd cancellation can improve performance but are not necessary
for the correctness of that equality test. This is the representation the
effective core of \cref{cas:sec-core} formalizes and all three delivered
prototypes implement, in three different specializations.

\subsection{Algebraic specifications and branch selectors}
\label{cas:sub-algebraic}
An algebraic object should contain a polynomial over the declared coefficient
field together with enough data to select one root. A bare equation $P(z)=0$
denotes a finite set or a relation, not one number. Valid selectors include an
ordered real root index, a Thom encoding, a certified simple residual root after
a scale change, a certified infinitesimal expansion distinguishing a branch, or
a tower of ordered algebraic extensions. A valuation ball containing a multiple
cluster is not a unique-root selector, and a numeric approximation to an
ordinary real root is not a sufficient selector for roots separated only by
infinitesimals. An arbitrary integer root index does not inherit the meaning of
an ordinary complex root ordering at non-Archimedean coefficients
\cite{cas:WLRoot}.

An exact algebraic number need not have a finite radical expression, and the
same is true over surreal coefficient fields. Implicit root objects are
therefore a central representation, not an unfortunate fallback: they give
finite descriptions of infinitely supported roots and algebraic closure of a
well-chosen effective core, at the cost of root-branch management and
potentially large algebraic degrees. Thus $\sqrt{1+t}$ is most naturally stored
as the positive root of $X^2-1-t$, with its binomial expansion an auxiliary view
rather than its definition. The output of an algebraic operation need not belong
to the input's finite-support grammar, yet can still have a finite exact
implicit representation.

For a simple binomial root the construction is explicit. If $a=c\,t^\gamma(1+\epsilon)$
with $v(\epsilon)>0$, a chosen $n$th root is
\begin{equation}\label{cas:eq-binomialroot}
 c^{1/n}\,t^{\gamma/n}\sum_{j\ge0}\binom{1/n}{j}\epsilon^j ,
\end{equation}
which requires a coefficient-root choice and the exponent $\gamma/n$; the series
is strongly summable, and multiplication by the $n$th roots of unity supplies
the other roots. Formula \eqref{cas:eq-binomialroot} is an exact lazy
description, usually not a finite normal form. Adjoining $\gamma/n$ may enlarge
the exponent group and adjoining $c^{1/n}$ may enlarge the coefficient field;
neither enlargement should be concealed. For complex $a$ a branch selector
remains necessary even though algebraic closedness supplies all roots, and
formulas such as $(ab)^{1/q}=a^{1/q}b^{1/q}$ must not be used with unrelated
complex root choices.

\subsection{Puiseux data and ramification}
\label{cas:sub-puiseux}
In rank one and characteristic zero, algebraic branches can be described by
Puiseux expansions after ramification: the union of the fields $k((t^{1/n}))$
over an algebraically closed coefficient field is algebraically closed
\cite[Corollary 7]{cas:Poonen}. For polynomial equations over an effective
rational-function base, Newton--Puiseux expansion with residual factorization
and branch selection is an appropriate implementation strategy.

The \emph{field of all Puiseux series} still permits arbitrary infinite
coefficient sequences, including noncomputable ones; its algebraic closedness
does not encode those sequences, and it should not be identified with the
finitely described algebraic functions. Finite defining equations plus branch
data give an effective class; an arbitrary coefficient stream leads back to
\cref{cas:thm-zerotest}. Higher-rank algebraic algorithms should use nested
valued fields, valuation extensions and explicit selected-root objects rather
than flattening every scale into one real exponent. Merely passing a custom Hahn
head into a general-purpose \texttt{Root} or \texttt{Solve} call does not
install those algorithms.

\subsection{Grids, recursive towers, and unrestricted Hahn data}
\label{cas:sub-grids}
We use the more general of the two grid definitions given by the sources.

\begin{definition}[Effective positive grid]\label{cas:def-grid}
Fix vectors $\gamma_0,e_1,\ldots,e_m\in\Q^d$, each $e_j$ lexicographically
positive. A grid presentation has support contained in
\begin{equation}\label{cas:eq-grid}
 \gamma_0+\N e_1+\cdots+\N e_m ,
\end{equation}
and supplies its coefficients either by a direct exponent procedure or by a
coefficient procedure on the multi-indices, with collisions added. \textbf{A
finite union of such presentations is permitted.}
\end{definition}

Two of the three sources use only the single-translate form; the finite-union
form is strictly more general and costs nothing. The coefficient rule on the
grid is still part of the input. Finite grids are useful because the support
condition follows from positivity and the Hahn--Neumann lemma, and because
coefficient extraction can be made effective under explicit assumptions
(\cref{cas:thm-grid}). Different tuples may represent the same exponent, so
coefficients must be combined by finite convolution rather than treated as
distinct monomials. Grid-based methods have a substantial existing computational
theory \cite{cas:vdH} and should not be presented as newly invented. They do not
make every coefficient sequence decidable, nor do they imply that every
valuation cutoff contains only finitely many terms
(\cref{cas:sub-finiteprefix}).

Recursive exponent towers can represent finitely many nested Conway monomials
whose exponents themselves have finite normal forms; a well-founded finite
dependency graph makes arithmetic and comparison recursive. A monomial with
exponent $\omega$ is finitely describable even though it is not in a rank-one
rational-exponent lattice. Such a language captures far more than rational
powers, but every extension must specify normalization of exponents, cross-level
coercions, and which operations preserve the syntax. A recursive exponent
backend should export its own equality, order and normalization procedures with
a well-founded construction rank; an arbitrary reference back to the same object
is not a legitimate normal form. A general expression in \texttt{SurrealExp},
\texttt{SurrealLog} and \texttt{ConwayMonomial} need not possess a known
complete normalizer merely because its expression graph is finite.

General Hahn supports can be well ordered without lying in a finite grid. Even a
countable support of order type $\omega+1$ cannot be traversed completely by a
monotonically increasing ordinary stream: after its first $\omega$ terms, no
finite step reaches the final term. Such supports require nested blocks,
symbolic support sets, indexed access, or exact structural expressions.
``Lazy'' must not be used as a synonym for ``an ordinary iterator suffices''.
A useful support language should include finite supports, arithmetic
progressions, positive finitely generated grids, nested lexicographic blocks and
selected well-ordered families with supplied dependency routines; more exotic
computable ordinal-indexed supports may be added as optional extensions.

At the broadest level one can store an abstract Hahn datum with a proof of
well-ordered support and a coefficient specification. This supports theorem
statements and some symbolic transformations. It should advertise only the
computational capabilities actually supplied by its certificates, rather than
pretending that a proof of existence is an enumeration algorithm.

\subsection{Expression graphs, transseries, and hyperseries}
\label{cas:sub-graphs}
A directed acyclic graph built from exact constants, field operations, algebraic
roots, selected exponentials and logarithms, and certified sums is a flexible
denotational language. Its finite size does not imply a finite normal form.
Common-subexpression sharing is particularly valuable for algebraic towers and
repeated local expansions. Descriptions using a rational generating function, a
total coefficient recurrence with specified initial data, or a formal implicit
equation with a uniqueness certificate are likewise finite and exact; these
mechanisms have different equality algorithms, and arbitrary coefficient
programs have none.

Transseries offer a structured enlargement for exponential and logarithmic
scales and for composition at infinity, with a substantial algebraic and
differential apparatus \cite{cas:vdH}; the 2026 D-algebraic zero test of
\cref{cas:rem-cfh} is an important additional algorithmic direction
\cite{cas:CFH}. Published work interprets LE-, EL- and omega-series as functions
on positive infinite surreals, with composition and differentiation in specified
classes \cite{cas:BMFunctions}. For implementation, select finite or effectively
generated fragments with explicit coefficient and support algorithms; the full
mathematical classes are not all finitely described.

The hyperseries representation theorem of Bagayoko and van der Hoeven provides a
much broader mathematical correspondence between surreal values and growth-rate
expressions \cite{cas:Hyperseries}. It is cited here for the mathematical result
only. It is not a finite-string encoding of every surreal; the countability
argument of \cref{cas:prop-countable} still applies, and a representation
theorem with infinitary data does not by itself supply normalization, equality,
or a storage format.

\subsection{A representation and closure map}
\label{cas:sub-closuremap}
\begingroup\small
\begin{longtable}{@{}P{0.24\textwidth}P{0.31\textwidth}P{0.37\textwidth}@{}}
\toprule
Representation & Exact values included & Main closure or decision boundary\\
\midrule\endhead
Finite cut trees / sign data & Dyadic rationals; finite-birthday and simplicity
data & Not closed under division; simplicity retained but not supplied by a
Hahn backend.\\[0.45em]
Compressed cuts and ordinal-indexed blocks & $\omega$, effectively described
reals, selected transfinite objects & Effectiveness lies in the option
language; simplest-number queries are not generic.\\[0.45em]
Finite sparse normal forms & Finite sums of admitted monomials & Closed under
addition and multiplication; inverse usually needs another representation.\\[0.45em]
Monomial rational functions & Fractions of those sums & Exact field algebra and
equality with effective independent generators; not algebraically closed.\\[0.45em]
Effective real closure and its complexification & All algebraic elements over
the chosen ordered monomial field & Real order and algebraic equality decidable;
complexity may be severe.\\[0.45em]
Rational or algebraic series & Infinite supports defined by finite algebraic
data & Strong exact structure; root selection remains necessary.\\[0.45em]
Certified recurrence or D-finite series & Many special-function expansions,
including divergent formal series & Coefficients effective with adequate initial
data; closure and identity are representation-specific.\\[0.45em]
Certified positive grids (finite unions) & Multiscale series with finite
generators & Fixed coefficients can be effective; a bounded valuation prefix can
still be infinite.\\[0.45em]
Selected transserial syntax & Exponential and logarithmic scale hierarchies &
Requires a specified ordered differential and composition framework.\\[0.45em]
Analytic germs and Hahn families & Local functions, finite jets, coherent
periods & Domain, support, branch and ordinary coefficient algorithms must be
supplied.\\[0.45em]
General program- or proof-described Hahn data & A broad countable collection of
exact values & No universal equality, support validation or truncation
algorithm.\\
\bottomrule
\end{longtable}
\endgroup

\section{The effective core}
\label{cas:sec-core}

\subsection{One core theorem at the generality all three sources support}
\label{cas:sub-coretheorem}
All three source articles prove the same effective-core theorem, over three
different coefficient fields and with three different treatments of rational
exponents and of the real closure. Those are differences of scope, not of
mathematics. We state the theorem once at the generality all three support, and
then print the three specializations.

\begin{theorem}[Effective rational monomial core]\label{cas:thm-core}
Let $k$ be an effective field, with an effective order when real comparisons are
requested. Let $\Gamma$ be an effectively ordered additive group and let
$X_1,\dots,X_d$ be monomials whose exponents $\beta_1,\dots,\beta_d\in\Gamma$ are
$\Q$-linearly independent. Then $k(X_1,\dots,X_d)$ embeds in $k((t^{\Gamma}))$,
and finite rational-function descriptions give an \emph{effective field}:
arithmetic terminates, and equality of $P/Q$ and $P'/Q'$ is decided by the zero
test on the single polynomial $PQ'-P'Q$. For nonzero $f=P/Q$, writing $\alpha(P)$
for the least exponent of $P$ in the induced order and $c(P)$ for that
coefficient,
\begin{equation}\label{cas:eq-leadingrational}
 v(f)=\alpha(P)-\alpha(Q),\qquad
 \lc(f)=\frac{c(P)}{c(Q)},\qquad
 \sgn(f)=\sgn\bigl(c(P)c(Q)\bigr),
\end{equation}
the last for real inputs, where $\sgn$ is the sign in $k$. If $v(f)>0$ then
$\st(f)=0$; if $v(f)=0$ then $\st(f)=\lc(f)$; if $v(f)<0$ then $f$ is infinite
and the finite standard-part operation is inapplicable. The zero case is handled
separately, with $v(0)=+\infty$ as a distinguished marker.
\end{theorem}
\begin{proof}
Distinct integer exponent vectors give distinct values of $\sum n_j\beta_j$ by
$\Q$-linear independence, so a nonzero polynomial in the $X_j$ maps to a
nonempty finite normal form and cannot vanish; the generators are algebraically
independent over $k$ and the embedding extends to fractions. Concretely, write a
nonzero denominator as $Q=c\,t^\beta(1+\epsilon)$ with all exponents of the
finite $\epsilon$ strictly positive; the strong geometric series
$\sum_{n\ge0}(-\epsilon)^n$ is a Hahn inverse, justified coefficientwise by the
positive-support lemma of \cref{cas:sub-twoconditions}.

Polynomial arithmetic over $k$ is effective because coefficient arithmetic and
equality are. Equality of fractions reduces to the zero test on $PQ'-P'Q$, a
finite coefficient computation. The leading term of a product cannot cancel
because the least exponent of each factor is unique, which gives
\eqref{cas:eq-leadingrational}; the Hahn ordering gives the sign statement. No
infinite expansion is required by any of these algorithms.
\end{proof}

Polynomial gcd reduction, when supplied by the coefficient backend, yields
reduced fractions, and normalizing the denominator's leading coefficient to one
removes scalar ambiguity; neither is required for the equality test.

\subsection{The three specializations, and which kernel implements which}
\label{cas:sub-threecores}
\begingroup\small
\begin{longtable}{@{}P{0.155\textwidth}P{0.15\textwidth}P{0.175\textwidth}P{0.15\textwidth}P{0.115\textwidth}P{0.115\textwidth}@{}}
\toprule
Source article & Coefficient field & Exponent group and generators &
Coordinate significance & Rational powers inside the theorem? & Real closure
proved?\\
\midrule\endhead
Capability tiers (02) & $\RA$ in the theorem; $\Q$ in its kernel &
$\Z$-spans of $\gamma_j=\omega^{j-1}$, i.e.\ $t_j=\omega^{-\omega^{j-1}}$ &
\textbf{Reversed}: highest-index coordinate compared first & No &
Yes, separately, as $R_d$ with $C_d=R_d[i]$\\[0.5em]
Layered representations (04) & $\RA$ & $\Q^r_{\mathrm{lex}}$, embedded by
$(q_1,\dots,q_r)\mapsto q_1\omega^{r-1}+\cdots+q_r$ & First coordinate most
significant & \textbf{Yes}, by common-denominator lattice refinement, inside the
same theorem & Yes, as $F_r$ with $C_r=F_r[i]$ algebraically closed\\[0.5em]
Gaussian prototype (05) & \textbf{An arbitrary effective field} $k$ &
$\Z^d_{\mathrm{lex}}$, standard basis & First coordinate most significant &
No; a remark on lattice refinement only & Discussed as a design target, not
proved\\
\bottomrule
\end{longtable}
\endgroup

Two strengthenings survive from this table and must not be flattened. The
\emph{rational-power extension} belongs to the layered-representations article:
choose a common positive denominator in each coordinate for all exponents
appearing in the finitely many operands, replace $T_j$ by $U_j^{N_j}$ where
$U_j=t^{\beta_j/N_j}$, and the same proof applies to the algebraically
independent $U_j$. The \emph{arbitrary-$k$ generality} belongs to the
Gaussian-prototype article, and is what licenses the statement of
\cref{cas:thm-core} above. The \emph{reversed coordinate convention} belongs to
the capability-tiers article, and is the convention its kernel actually matches.

The concrete embedding shared by the second and third specializations sends
$(q_1,\dots,q_r)\mapsto q_1\omega^{r-1}+\cdots+q_r$, so that every earlier
generator is smaller than every positive integer power of a later one; for $r=2$
the generators are $\omega^{-\omega}$ and $\omega^{-1}$. The first
specialization's generators are $\gamma_j=\omega^{j-1}$, giving
$v(t_1)=1,\ v(t_2)=\omega,\ \dots$, so that $0<t_{j+1}<t_j^{\,n}$ for every
ordinary positive integer $n$, with the exponent vector $(q_1,\dots,q_d)$
ordered by its \emph{highest-index} differing coordinate first.

\begin{remark}[State the convention; never infer it from names]
\label{cas:rem-convention}
The order convention is not inferred from names such as \texttt{eps1} and
\texttt{eps2}; it belongs to the domain descriptor. Different conventions give
different ordered embeddings even when the same polynomial syntax is used. This
report therefore fixes one global rule --- \emph{lexicographic order on a
finitely generated ordered exponent group, with coordinate significance declared
as part of the domain descriptor} --- and then states each kernel's own
convention in that kernel's own subsection, quoting each kernel's outputs in its
own convention and never rewriting them into a common form
(\cref{cas:sub-orderconventions}).
\end{remark}

\subsection{The cancellation trap}
\label{cas:sub-cancellationtrap}
Cancellation before leading-term extraction is essential. In the capability-tiers
convention, the numerator
\[
 (1+t_1)t_2-t_2
\]
has leading exponent $(1,1)$, \emph{not} $(0,1)$. A heuristic inspection of the
uncollected expression is therefore wrong. This is not an academic point: the
shipped kernel of that article carries a check for exactly this case, and
re-executing it for this report returns the exponent vector
$\{1,1\}$ (\cref{cas:sec-executed}).

\subsection{Effective real closure and complexification}
\label{cas:sub-realclosure}
Let $E$ denote an effective ordered rational-monomial field as in
\cref{cas:thm-core} over a real coefficient field, let $F$ be its relative real
closure realized inside $\No$ --- the elements of $\No$ algebraic over $E$,
which is a real closure because $\No$ is real closed --- and put $C=F[i]$.

\begin{theorem}[An implementable real-closed and algebraically closed layer]
\label{cas:thm-realclosure}
$F$ admits finite effective algebraic descriptions with decidable real order and
equality, and $C$ admits finite effective descriptions with decidable equality
and field operations. Every finite-degree univariate polynomial over $C$ has all
its roots in $C$. Semialgebraic sentences over these effective coefficient
fields can be decided, and finite polynomial systems can be treated by effective
algebraic methods with the output representation appropriate to their dimension.
No Hahn expansion of each algebraic element is required for these guarantees.
\end{theorem}
\begin{proof}[Algorithmic justification]
Represent a real algebraic element by a finite polynomial equation over $E$
together with a unique ordered-root selector, for example a realizable Thom
encoding: the signs of successive derivatives at the selected real root. A
finite tuple of such descriptions can be handled in one real closure, and a
finite tower is an alternative to reducing every operation immediately to a
single primitive element. Addition, multiplication and inversion introduce
finite systems expressing the operation together with the selectors of its
arguments; elimination produces equations for the result and sign determination
chooses the intended real root. Equality and order are likewise finite formulas
in the defining equations and selectors. Real-closed-field quantifier
elimination reduces every such formula to finitely many polynomial sign tests
over $E$, and those terminate by \cref{cas:thm-core}. Purely algebraic
quantifier-elimination and sign-determination algorithms based on Thom encodings
are available independently of Archimedean topology and are a published
foundation for this step \cite{cas:PerrucciRoy}.

Pairs of real elements represent $C$, with complex equality the equality of both
coordinates. Since $F$ is real closed, $F[i]$ is algebraically closed;
consequently a polynomial over it factors there and no additional roots appear
in a larger extension. Root specifications can also be obtained by writing the
polynomial equation in real and imaginary coordinates and selecting among its
finitely many solutions. For a polynomial system the same real-algebraic
encoding decides existential solvability and supplies descriptions of its
solution conditions; a positive-dimensional solution set is represented by
equations or a suitable parametric or semialgebraic description, not by a finite
list of points.
\end{proof}

\begin{remark}[Implementability, not complexity]\label{cas:rem-complexity}
\cref{cas:thm-realclosure} establishes \emph{implementability}, not acceptable
complexity in every case. Elimination can be expensive and degree growth severe.
In practice, maintain extension towers and specialized univariate routines and
use general elimination as a completeness fallback. Quantifier elimination for a
mathematical field and executable elimination over a chosen input encoding are
separate assertions; an arbitrary surreal constant represented only by a symbol
does not supply the comparison oracle these algorithms require. Nor does the
theorem assert that a host's native \texttt{Root} or \texttt{Reduce} already
accepts an arbitrary custom surreal head as its coefficient domain: an
implementation must supply the reductions, the ordered-field adapter, and
root-description management.
\end{remark}

\subsection{The non-density trap: why ordinary interval isolation is insufficient}
\label{cas:sub-nondensity}
A subtle implementation trap is to assume that the ground ordered field is dense
in its real closure. It need not be, and all three sources give a witness ---
two infinitesimal and one infinite.

\begin{itemize}
\item In $E=\Q(t)$ with $t>0$ infinitesimal there is no $q\in E$ with
$\sqrt t<q<2\sqrt t$. Any positive $q$ has leading form $c\,t^m$ with $m\in\Z$
and $c>0$; if $m\le0$ it exceeds both endpoints, and if $m\ge1$ it is below
both. Rational-function endpoints therefore cannot separate the roots $\sqrt t$
and $2\sqrt t$ by a point between them.
\item In $F=\RA(t)$ with $t>0$ infinitesimal, $\alpha=\sqrt t$ and
$\beta=\sqrt t+t$ are distinct, yet no element of $F$ lies strictly between
them, by the same integer-valuation argument. Root isolation by refining
endpoints inside $F$ cannot separate this pair.
\item If $T$ is positive \emph{infinite}, no element of $\Q(T)$ lies between
$\sqrt T$ and $\sqrt T+1$: a nonzero rational function of $T$ has an integer
leading power, whereas every point of that interval has leading power $T^{1/2}$.
\end{itemize}

The obstruction is the same in all three cases and is an algorithm-design
obstruction, not a failure of ordered algebraic root representation. Real-root
indices, Thom encodings, algebraic endpoints, or enlargement of the monomial
field by fractional powers all resolve it; naively porting an
interval-bisection algorithm that assumes rational density does not give a
complete method. In particular, \emph{increasing the decimal precision of a real
number cannot distinguish two roots with identical real standard part and
infinitesimal separation}, at any precision.

\subsection{What the core deliberately omits}
\label{cas:sub-coreomits}
The countable union of finite-rank algebraic cores is useful but is not the
entire surreal field. The core admits $\omega$, $\omega^{-1}$,
$\omega^{-\omega}$, their rational powers and rational combinations at
\emph{finitely many} hierarchy levels; it does not contain arbitrary surreal
exponents such as every possible infinite normal form. Algebraic closure does
not imply exponential closure, support closure under every countable sum, or
access to birthdays. Even an exactly named transcendental such as $\pi$ is not
in the minimal algebraic coefficient domain; it can be added through a larger
coefficient backend, but that backend must declare what it can decide about the
new constants. Higher finite rank is added by an explicit order-preserving
embedding of the old coordinates into the new domain --- in the
first-coordinate-most-significant convention, passing from rank $r$ to rank
$r+1$ while preserving the old scales means prepending a zero coordinate.

A symbolic parameter may instead remain algebraically independent over the base.
This permits rational-function algebra, but is not an assertion that its
intended real value is independent of other named constants. The difference
between an indeterminate and a particular named number must survive
serialization, substitution and simplification.

\subsection{What transfer does and does not buy}
\label{cas:sub-transfer}
The ordered-field laws and the real-closed algebraic theorems transfer to this
core, so finite matrix identities, polynomial identities and semialgebraic
geometry can be reused with a coefficient-domain adapter. The statement that $t$
is smaller than $1/n$ for every \emph{ordinary integer} $n$ is not a single
ordinary ordered-field formula defining an infinitesimal in $\R$, and passing a
few inequalities to a real solver is not an implementation of that condition.

What should therefore \textbf{not} be done is to ask a native real-domain
\texttt{Reduce} to find a positive real $t$ smaller than $1/n$ for every
positive standard integer $n$. No such real exists. Replacing the condition by
finitely many inequalities only specifies a small ordinary parameter, not an
infinitesimal. The mathematical transfer principle concerns finite first-order
field formulas in an appropriate real-closed model; it does not turn unrestricted
quantification over the integers into a field-language assumption
\cite{cas:SourcePolynomial,cas:PerrucciRoy}. Likewise, an elementary-extension
theorem for a surreal exponential structure is a semantic theorem, not a
universal decision algorithm for exponential expressions \cite{cas:vdDE}; a CAS
must not infer that its ordinary real simplifier now knows every surreal
exponential identity.

One route that does work: perform real-closed-field quantifier elimination with
symbolic parameters first, then evaluate the resulting finitely many sign
conditions using the custom ordered-field backend.

\section{Certified supports and finite dependency}
\label{cas:sec-support}

\subsection{The two conditions, and four standard non-examples}
\label{cas:sub-twoconditions}
A family $(A_j)_{j\in J}$ of Hahn series is \emph{strongly summable} when the
union of supports is well ordered \emph{and} every exponent receives
contributions from only finitely many members; the sum is then formed by finite
coefficient addition at each exponent. All six supplied manuscripts use exactly
this meaning. Both conditions are necessary, and the sources give the same four
standard non-examples, with $t=\omega^{-1}$:
\begin{align*}
 \sum_{n\ge0}t^n&=\frac1{1-t}\quad\text{(valid)}, &
 \sum_{n\ge0}n!\,t^n&\quad\text{is a valid Hahn sum},\\
 \sum_{n\ge0}t^{-n}&\quad\text{is \emph{not} a Hahn sum}, &
 \sum_{n\ge0}\frac1{n!}&\quad\text{is \emph{not} a strongly summable family}.
\end{align*}
The third support decreases indefinitely. In the fourth family infinitely many
summands use exponent zero; its ordinary real sum is of course $e$, and a CAS
may evaluate that \emph{ordinary coefficient sum} first and embed the result ---
a different operation, not a relaxation of finite coefficient multiplicity. The
same remark applies to $\sum_{n\ge0}2^{-n}=2$. Two further standard failures:
$\{1/n:n\ge1\}$ is a set of positive exponents that is not well ordered, so
positivity of every summand's valuation is insufficient; and $\sum_{n\ge1}t^{1/n}$
must be rejected as an unordered-support Hahn sum. The CAS must always specify
which summation operation is being used.

For well-ordered supports $S,T$, the support of a product lies in $S+T$ with
finitely many contributing pairs at each exponent. If $E$ is well ordered and
strictly positive, the monoid $E^*$ of finite sums from $E$ is well ordered and
each exponent has finitely many word decompositions up to the usual finite
permutations. For a source support $S$, expressions indexed by one member of $S$
and a finite word in $E$ have support in $S+E^*$, with finitely many
contributions at each exponent. This Hahn--Neumann calculus justifies reciprocal
series, infinitesimal Taylor substitution and compatible coefficient regrouping,
and is exactly the mechanism used repeatedly by AG and CT
\cite{cas:SourceAG,cas:SourceCT}.

The factorial series deserves separate mention: it has zero ordinary radius of
convergence but positive, well-ordered infinitesimal support. Its exact surreal
value does not automatically identify it with a particular analytic continuation
or summation of an ordinary divergent function; a special-function or
resurgent-extension provider must supply that extra interpretation
(\cref{cas:sub-factorial}).

\subsection{Compositional certificates and the five support services}
\label{cas:sub-services}
An implementation should construct support certificates compositionally: a
finite support supplies a base certificate; a product uses a Minkowski-sum
constructor; inversion of a normalized infinitesimal perturbation uses a
positive-monoid constructor; analytic substitution combines the outer support
with such a monoid; finite unions and translates are available throughout. A
bare statement that a user-supplied program ``probably produces increasing
exponents'' is not such a certificate.

A useful support interface distinguishes five separate services.
\begin{itemize}
\item \textbf{Membership}: is a given exponent inside a support bound?
\item \textbf{Coefficient}: return the actual coefficient, including cancellation.
\item \textbf{Decomposition}: list the finitely many ways an exponent can arise
in a product or substitution.
\item \textbf{Leading term}: additionally certify that all smaller coefficients
vanish.
\item \textbf{Prefix}: describe everything below a chosen cutoff.
\end{itemize}
\emph{No one of these may be inferred automatically from the existence of the
others.} For finite supports all five are straightforward. For infinite
supports, certificates can be built compositionally as above, but an arbitrary
well-order proof can justify a mathematical expression without furnishing any of
these operational services.

\subsection{A certificate is not an evaluation algorithm}
\label{cas:sub-notalgorithm}
To compute a requested coefficient, a backend also needs an effective way to
identify all contributing decompositions and to know when the list is complete.
\textbf{A theorem that each such list is finite does not supply a terminating
search that detects its end.} AG itself is credited by one of the sources with
issuing this warning about its own formulas. Useful constructors therefore carry
finite dependency bounds or explicit decomposer procedures in addition to
semantic well-ordering evidence.

Similarly, a transfinite coefficient recursion proves the existence of a
complete Hahn object without being an ordinary finite loop through the entire
support. A correct implementation either produces a finite implicit object,
computes specified coefficients with certified finite dependencies, or works
inside a restricted effective support grammar. It must not report that it has
``finished the transfinite recursion'' merely because it computed many initial
coefficients.

The capability-tiers article identifies precisely this missing service, and
demands finite dependency bounds or explicit decomposer procedures, without
supplying a construction. The other two articles supply one. We state the gap
first because it is the motivation, rather than letting the lemma below silently
absorb the caution.

\subsection{The finite-grid dependency lemma, with both constructions}
\label{cas:sub-gridlemma}
\begin{theorem}[Finite extraction from a positive grid]\label{cas:thm-grid}
Let $e_1,\dots,e_m\in\Q^d$ be lexicographically positive, let
$\gamma_0\in\Q^d$, and let $c:\N^m\to k$ be a total computable coefficient rule
in an effective field $k$. The monomial family
\begin{equation}\label{cas:eq-gridseries}
 \bigl(c(\boldsymbol n)\,t^{\gamma_0+\sum_jn_je_j}\bigr)_{\boldsymbol n\in\N^m}
\end{equation}
is strongly summable, and at any specified $\lambda\in\Q^d$ its coefficient is
computable by a finite search --- even when the map
$\boldsymbol n\mapsto\gamma_0+\sum n_je_j$ is not injective.
\end{theorem}
\begin{proof}
The positive-monoid form of the Hahn--Neumann lemma makes the grid well ordered
and gives finite fibres. For an effective bound, it suffices to produce a
rational linear functional $\ell:\Q^d\to\Q$ with $\ell(e_j)>0$ for every $j$.
Two constructions are given below. Granting one, put $b=\ell(\lambda-\gamma_0)$.
Every nonnegative integer solution of $\sum_jn_je_j=\lambda-\gamma_0$ satisfies
$\sum_jn_j\ell(e_j)=b$. If $b<0$ there are no solutions; otherwise
\[
 0\le n_j\le\left\lfloor\frac{b}{\ell(e_j)}\right\rfloor
 \qquad\text{for each }j.
\]
Enumerate this finite box and test the \emph{original vector equation}, not
merely its image under $\ell$. The desired coefficient is the finite sum of
$c(\boldsymbol n)$ over the solutions found. The same bound proves the
finite-contribution requirement for strong summability directly; zero
coefficients can only shrink the support. For a fixed multi-index there are
finitely many ordered words and finite combinatorial multiplicities, so products
and substitutions also give finite coefficient computations whenever the
coefficient operations themselves terminate.
\end{proof}

\paragraph{Construction I: increasing weights, forwards.}
Try
\[
 \ell_M(x_1,\dots,x_d)=M^{d-1}x_1+M^{d-2}x_2+\cdots+Mx_{d-1}+x_d
\]
for successively larger positive integers $M$. For each lexicographically
positive $e_j$, its first nonzero coordinate eventually dominates the later
coordinates in this expression. Because there are finitely many $e_j$, the
search terminates. \emph{This is the construction that is actually exercised in
code}: the layered-representations article's Python script implements it with
$M$ doubling (\cref{cas:sec-executed}).

\paragraph{Construction II: choosing weights backwards.}
Start with weight $1$ on the last coordinate and work backwards. When choosing
the weight on coordinate $k$, make it large enough that every $e_j$ whose first
nonzero coordinate is $k$ has positive weighted sum, despite its already
weighted later coordinates; such a first nonzero coordinate is positive by
lexicographic positivity. There are only finitely many inequalities. This
construction is not implemented in any delivered code.

The two constructions are genuinely different algorithms for the same object and
both survive here.

\begin{remark}[The functional is not an order embedding]\label{cas:rem-notembedding}
Both sources that supply a construction issue the same warning independently,
and it is load-bearing. The functional $\ell$ is an \textbf{auxiliary device on
the finite generator set only}. It is \emph{not} an order-preserving embedding
of the whole lexicographic group into $\Q$ or $\R$ --- no such embedding exists
for rank greater than one --- so it must never be reused as a substitute
valuation for arbitrary cutoff comparisons. Its finite box proves termination
but can be large; integer linear-equation methods, dynamic programming, sparse
convolution and memoized dependencies can improve it, and must preserve exact
treatment of coincident exponent sums by \emph{adding} their coefficients rather
than retaining multiple conflicting entries. A total coefficient rule still
needs a supplied termination argument, or membership in a language where
termination is built into the syntax.
\end{remark}

\subsection{The small-scale theorem and its effectivity gap}
\label{cas:sub-smallscale}
SA proves that an arbitrary formal series $\sum a_nX^n$ with set-many surcomplex
coefficients can be evaluated after a sufficiently small rescaling; its proof
chooses a surreal exponent dominating a set of exponents drawn from \emph{all}
the coefficients and shifts their supports into separated bands
\cite{cas:SourceAnalysisTwo,cas:SourceAnalysisThree}. For arbitrary surreal
coefficients, positivity of $v(\epsilon)$ alone is not enough: the coefficient
valuations may decrease too rapidly.

This is a valuable semantic existence theorem and \textbf{not} a generic finite
procedure for discovering that bound from an arbitrary coefficient program. An
effective interface should ask for a support-growth bound, a dominating-scale
certificate, or a restricted coefficient grammar from which one can be derived.
The constructed scale may also lie outside the user's current workspace; domain
enlargement must be explicit rather than hidden inside a purported fixed-field
operation.

\subsection{Inverses, positive recursions, and the reciprocal recurrence}
\label{cas:sub-inverses}
For a known nonzero series write
\begin{equation}\label{cas:eq-unit}
 A=c\,t^a(1+E),\qquad c\ne0,\quad \supp(E)\subset\Gamma_{>0}.
\end{equation}
Then
\begin{equation}\label{cas:eq-inverse}
 A^{-1}=c^{-1}t^{-a}\sum_{n\ge0}(-E)^n
\end{equation}
is an exact Hahn identity. For finite $E$, or for a certified positive-grid
representation of $E$, the support and finite decompositions are explicit. If a
general input has no computable leading normalization, \eqref{cas:eq-inverse} is
a \emph{mathematical identity rather than a ready algorithm}. A package can
retain an exact rational-expression inverse without expanding it when that
representation is available --- which is exactly what all three delivered
prototypes do.

The same holds for HC's positive-operator formula
$(I+T)^{-1}X=\sum_{n\ge0}(-T)^nX$ and its nonlinear preparation recursions,
whose finite coefficient dependencies are controlled by support words
\cite{cas:SourceResearch}. To make them executable an implementation needs
effective word decomposition and executable ordinary coefficient operations; for
a fixed-domain analytic division problem those operations can themselves involve
ordinary integrals or implicitly specified functions. Such a coefficient is
still representable by a well-defined symbolic operator when no elementary
closed form is available, but it is not thereby in the smallest decidable scalar
field.

For a normalized one-parameter series $A(t)=1+\sum_{n\ge1}a_nt^n$, the reciprocal
coefficients come from a recurrence needing no symbolic differentiation:
\[
 b_0=1,\qquad b_n=-\sum_{k=1}^n a_kb_{n-k}.
\]
Each requested $b_n$ has finitely many dependencies. In a general Hahn group the
corresponding recursion is indexed by actual exponents; a positive-support
certificate supplies the well-founded relation, but a computational decomposer
is still required.

\section{Precision is a mathematical type, not a display setting}
\label{cas:sec-precision}

\subsection{The \texorpdfstring{$\Ov{\cdot}$}{O\_v} convention}
\label{cas:sub-ovconvention}
Write
\[
 x=\widetilde x+\Ov{\alpha}
 \quad\Longleftrightarrow\quad
 v(x-\widetilde x)\ge\alpha,\qquad v(0)=+\infty .
\]
This is an exact enclosure statement in the chosen valued workspace: a statement
about a \emph{set} of possible errors, not an equality in the field after
deleting a formal $O$-symbol, and not ordinary real asymptotic big-$O$. The
precision parameter $\alpha$ ranges over the \emph{full ordered value group},
not over an ordinary numerical accuracy parameter. A strict bound
$v(\cdot)>\alpha$ is a separately represented thing; the convention used
throughout this report is the \textbf{weak} inequality unless a strict one is
displayed.

\subsection{Certified precision propagation}
\label{cas:sub-precisionrules}
All three sources state the following identically, with the same convention.

\begin{proposition}[Basic certified precision rules]\label{cas:prop-precision}
Suppose $A=\widehat A+e_A$ and $B=\widehat B+e_B$ with $v(e_A)\ge\alpha$ and
$v(e_B)\ge\beta$. Then
\begin{align}
 v\bigl((A+B)-(\widehat A+\widehat B)\bigr)&\ge\min\{\alpha,\beta\},
 \label{cas:eq-addprecision}\\
 v\bigl(AB-\widehat A\widehat B\bigr)&\ge
 \min\{\alpha+v(\widehat B),\ \beta+v(\widehat A),\ \alpha+\beta\}.
 \label{cas:eq-mulprecision}
\end{align}
If $\widehat A\ne0$, $\mu=v(\widehat A)$ and $\alpha>\mu$, then $A\ne0$,
$v(A)=\mu$, and
\begin{equation}\label{cas:eq-invprecision}
 v\bigl(A^{-1}-\widehat A^{-1}\bigr)\ge\alpha-2\mu .
\end{equation}
\end{proposition}
\begin{proof}
The sum error is $e_A+e_B$. The product error is
$e_A\widehat B+\widehat Ae_B+e_Ae_B$; apply the valuation inequalities. Under
the last hypothesis the leading term of $\widehat A$ cannot be cancelled by
$e_A$, so $v(A)=\mu$. Finally
$A^{-1}-\widehat A^{-1}=-e_A/(A\widehat A)$, of valuation at least
$\alpha-2\mu$.
\end{proof}

These formulas expose genuine precision loss through inversion of a small
quantity: an error of valuation $10$ in a quantity of valuation $3$ gives an
inverse error bound of only $4$. Preserving the same absolute cutoff through
division would be incorrect.

\begin{remark}[Dependency tracking]\label{cas:rem-dependency}
A centre polynomial that cancels to zero does not prove the exact answer zero.
Subtracting two identical truncations of independently represented numbers
leaves a remainder enclosure unless exact equality of the underlying objects has
been proved. Conversely, subtracting an object from itself should use dependency
information and return exact zero: exact identity tracking can prove $A-A=0$
even when only an approximate expansion of $A$ has been computed, and treating
the two appearances as unrelated error balls loses that information
unnecessarily. A zero test on a truncated series may return ``zero to the
supplied precision'', but \emph{that is not exact zero}. If a nonzero leading
coefficient is known below the error cutoff, sign and leading valuation are
certified; if all retained terms cancel, more information is required.
\end{remark}

\subsection{The positive finite-prefix criterion, and three counterexamples}
\label{cas:sub-finiteprefix}
\begin{proposition}[Finite-prefix criterion for rank-one grids]
\label{cas:prop-finiteprefix}
For $a,e_1,\dots,e_m,B\in\Q$ with every $e_j>0$, the grid $a+\sum_j\N e_j$
contains only finitely many exponents below $B$, and those exponents can be
enumerated effectively. The corresponding statement is false for general
lexicographically ordered groups, even with one positive generator.
\end{proposition}
\begin{proof}
If $a+\sum_jn_je_j<B$ then every $n_j<(B-a)/e_j$, so a finite box contains all
possibilities. For the counterexample use generator $(0,1)$ and cutoff $(1,0)$
in $\Q^2_{\mathrm{lex}}$: every $(0,n)$ lies below that cutoff.
\end{proof}

This positive criterion is the only one of its kind in the group, and it is what
makes the rank-one case tractable. It is also exactly as far as the positive
result goes. Three counterexamples, of two genuinely different kinds, delimit it,
and no one of them should be allowed to stand for the family.

\paragraph{Higher-rank kind, witness 1 (grid).}
In an ordered group containing positive $\lambda,\mu$ with $n\lambda<\mu$ for
every ordinary $n$, put
\begin{equation}\label{cas:eq-cutoff}
 A=\sum_{n\ge0}t^{n\lambda}+t^\mu=\frac1{1-t^\lambda}+t^\mu .
\end{equation}
The support $\{n\lambda:n\ge0\}\cup\{\mu\}$ is well ordered and lies in the grid
generated by $\lambda$ and $\mu$, and every first summand has exponent below
$\mu$. Hence the truncation below $\mu$ has infinitely many nonzero terms, and a
finite list cannot represent that support completely. Moreover, after $N$
geometric terms the remainder has valuation $(N+1)\lambda$ and \emph{never}
reaches $\mu$: a refinement loop that counts terms runs forever, even though
every finite step is mathematically correct.

\paragraph{Higher-rank kind, witness 2 (geometric remainder plus $u$).}
Take $\Gamma=\Q\omega+\Q$ and $u=t^\omega$, so $u<t^n$ for every ordinary
positive $n$. The exact remainder identity is
\begin{equation}\label{cas:eq-geometricremainder}
 \frac1{1-t}-\sum_{j=0}^{N-1}t^j=\frac{t^N}{1-t},
 \qquad v\!\left(\frac{t^N}{1-t}\right)=N<\omega .
\end{equation}
No finite prefix of this expansion achieves error valuation at least $\omega$,
yet the strong sum exists and the left side is a finite exact description. The
same issue affects $\sum_{n\ge0}t^n+u$: an ordinary stream that insists on
emitting all terms in increasing valuation order never reaches the $u$ term.

\paragraph{Rank-one nondiscrete kind.}
Even rank one can fail the finite-prefix property. The supplied example
\begin{equation}\label{cas:eq-accumulating}
 A=\sum_{k\ge0}t^{1-1/(k+2)}
\end{equation}
has well-ordered support of order type $\omega$ with infinitely many terms below
exponent $1$ \cite{cas:SourcePolynomial,cas:SourceResearch}; equivalently, the
set $\{1-1/n:n\ge1\}$ is increasing and well ordered but has infinitely many
elements below $1$. Hence \emph{``rank one'' alone is not enough}: discreteness
or a suitable positive-grid bound is needed. The obstruction is therefore not
merely about rank.

\subsection{Three truncation interfaces, and two worked compressions}
\label{cas:sub-truncation}
The shared conclusion of all three sources is that a single ``precision''
request conflates three different things, which must be separate interfaces:
\begin{enumerate}
\item a finite number of \textbf{leading terms};
\item a \textbf{valuation-cutoff object}, whose representation may remain
compressed;
\item a finite rectangular or total-degree \textbf{jet} in selected grid
generators.
\end{enumerate}
To these the layered-representations article adds a fourth request, a single
fixed coefficient, which the others treat as part of the coefficient service.

There is no obstruction to an \emph{exact compressed answer}. Two worked
compressions appear in the sources. In \eqref{cas:eq-cutoff} the truncation
below $\mu$ is exactly $(1-t^\lambda)^{-1}$. In
\eqref{cas:eq-geometricremainder} with $N=10$, the exact value of
$1/(1-t)-\sum_{j=0}^{9}t^j$ is the rational block $t^{10}/(1-t)$, of valuation
$(0,10)$ in the two-coordinate parent $\{u,t\}$, while the requested error scale
$u$ has valuation $(1,0)$, strictly larger than $(0,N)$ for every ordinary
integer $N$. A correct finite-prefix routine must report that the requested cut
needs an infinite block; a compressed-truncation routine returns the rational
block exactly.

\begin{remark}[Regression framing]\label{cas:rem-regression}
This example is a regression test for any future precision module:
\emph{increasing an ordinary integer precision parameter must never claim to
certify accuracy beyond $u$}. The valuation-vector comparison, not a numerical
norm chosen for convenience, decides whether a claimed bound is valid.
Correspondingly, a precision request whose result cannot be a finite term list
must return a compressed certified subseries, an explicit unsupported-request
result, or an alternative finite jet chosen by the caller. It must
\textbf{never} silently stop after an arbitrary number of terms and claim the
requested bound.
\end{remark}

\subsection{Nested descriptions and the \texorpdfstring{$\omega$th}{omega-th} position}
\label{cas:sub-nested}
For higher-rank data, nested descriptions are useful: a coefficient at one outer
scale may itself be an infinite series in a smaller-rank field, which compactly
exposes a term occurring after an order-type-$\omega$ block. A whole $t$-block
can be stored as a rational or lazy coefficient in an outer $u$-series; another
suitable representation provides direct exponent access. The nesting order must
agree with the intended ordered group, and operations must respect the resulting
support conditions. A flat iterator that only returns the next
natural-number-indexed term cannot reach an $\omega$th support position in
finitely many calls, and neither mechanism is equivalent to taking more digits
in a single real-valued precision parameter.

\subsection{Coefficient precision is a second, independent axis}
\label{cas:sub-coeffprecision}
Surreal scale precision and ordinary numerical precision are independent. The
coefficient of $t^{1/2}$ may be known to 100 binary digits while the series is
known only through exponent $3$; increasing coefficient precision does not
discover an omitted $t^4$ term, and producing more exact rational coefficients
does not yield an ordinary floating-point approximation to an infinite surreal
value. An interval coefficient enclosure containing zero does not determine a
leading sign. For nonalgebraic ordinary coefficients, validated real or complex
enclosures can sometimes certify nonvanishing and sign; if an enclosure contains
zero one must refine it or retain an unresolved status. \emph{Heuristically
small is not exactly zero}, and an ordinary approximate coefficient must not be
silently recast as a nearby rational value inside an exact-number constructor.
A serious system records both the exponent cutoff and the coefficient-error
model and refines whichever prevents the requested conclusion. Certified
interval or ball arithmetic for coefficients and exact valuation truncation
solve different problems.

\subsection{These are not fine-topological limits}
\label{cas:sub-finetopology}
The full-class fine topology of SA uses every positive surreal tolerance, and in
it \textbf{every set-indexed convergent net is eventually equal to its limit}
\cite{cas:SourceAnalysisTwo,cas:SourceAnalysisThree}. A sequence of finite
truncations must therefore not be advertised as converging in that topology.
Within a chosen rank-one valued field a sequence of truncations can converge in
its valuation topology; the same words refer to a different structure. Nor does
the ordinary sequence $1/n$ converge finely to zero in the full surreal field: a
fixed positive surreal infinitesimal is smaller than every term.

An ordinary real limit and a surreal evaluation are different commands:
\[
 \lim_{X\to+\infty}\frac1{X+1}=0\ \text{in ordinary real analysis},
 \qquad\text{whereas}\qquad
 \frac1{\omega+1}>0\ \text{in }\No .
\]
The latter retains the infinitesimal information the former discards. Limit
modes, Hahn summation, and finite-scale specialization must remain explicit
operations rather than different spellings of one implicit ``infinite
evaluation'' routine.

\subsection{Jets, asymptotics, and exact tails}
\label{cas:sub-jets}
A finite Taylor jet in $K[z]/(z)^N$ has exact quotient-ring semantics: it
determines products modulo $(z)^N$, not a unique analytic function. An
asymptotic object records a function or value modulo an explicitly stated error
class. An exact lazy object specifies all coefficients even though only finitely
many are cached, and an exact rational tail specifies infinitely many
coefficients in one node. These distinctions must survive serialization and
printing. A finite multivariate jet is a quotient by a chosen monomial ideal in
auxiliary variables; after substituting the grid generators, its omitted
monomials must be converted into an explicit valuation error bound, a conversion
that may be coarse and may not describe an entire valuation prefix.

\subsection{Rounding and the \texorpdfstring{$\Pi\oplus\Z$}{Pi + Z} floor}
\label{cas:sub-floor}
For an infinite surreal there is no ordinary integer floor in $\Z$: no ordinary
integer brackets such a value, so an ordinary $\Z$-valued floor cannot satisfy
the usual defining inequalities. A package must state what its target
``integers'' are before defining a global floor, rather than silently reusing
ordinary integer rounding.

One useful normal-form choice is the ordered subring $\Pi\oplus\Z$, where $\Pi$
is the purely infinite additive class. Products of purely infinite parts have
only negative exponents, so this is indeed a subring. For $x=p+r+\eta$ with $p$
purely infinite, $r\in\R$ and $\eta$ infinitesimal, define
\begin{equation}\label{cas:eq-floor}
 \lfloor x\rfloor_{\Pi\oplus\Z}
 =p+\begin{cases}
 \lfloor r\rfloor-1,&r\in\Z\text{ and }\eta<0,\\
 \lfloor r\rfloor,&\text{otherwise.}
 \end{cases}
\end{equation}
This is the \emph{unique} element $m\in\Pi\oplus\Z$ with $m\le x<m+1$:
noninteger $r$ has positive ordinary distance from its neighbouring integers,
while at an integer the infinitesimal sign selects the side, and any distinct
purely infinite part differs by an infinite quantity and cannot satisfy the same
unit-interval condition. Thus
\[
 \lfloor1-t\rfloor=0,\qquad
 \lfloor\omega/2\rfloor_{\Pi\oplus\Z}=\omega/2 ,
\]
the second not an ordinary integer. The operation is effective only when the
backend can split the normal form, decide the ordinary integer predicate, and
determine the infinitesimal sign; a compact rational node need not expose its
entire purely infinite part as a finite list. An omnific-integer convention is a
different target again and needs its own definition.

\section{Polynomial algebra and root-cluster algorithms}
\label{cas:sec-polynomials}

\subsection{Finite algebra over a declared coefficient domain}
\label{cas:sub-finitealgebra}
Over an effective field, polynomial degree is an ordinary finite integer. Long
division terminates by degree descent; the Euclidean algorithm gives gcds and
B\'ezout identities; squarefree decomposition uses the derivative; resultants
and discriminants are finite determinants; matrices admit exact row reduction
whenever pivot zero tests are available. These apply to the core of
\cref{cas:sec-core} without expanding any coefficient into an infinite Hahn
normal form. Gaussian elimination decides matrix rank and solves consistent
systems, characteristic polynomials and algebraic root selectors represent
eigenvalues, positive square roots give distances, and exact determinant signs
decide orientation --- so TR's triangle identities can be supported without ever
choosing an infinite circular phase, since normalized directions lie on the
finite-coordinate unit circle \cite{cas:SourceTrig}.

For multivariate polynomials, Buchberger-type ideal algorithms depend on finite
monomial orders and effective coefficient operations; a completed Gr\"obner
basis certifies ideal membership or inconsistency by a finite identity. A
zero-dimensional quotient has a finite basis and multiplication matrices, while
a positive-dimensional solution set needs a different output shape: the
mathematical existence of a root over an algebraically closed Hahn field is not
a promise to list a positive-dimensional family.

If coefficient equality is only partial, the polynomial routines become partial
or conditional. An unresolved leading coefficient can change the degree; an
unresolved pivot can change matrix rank; a determinant that fails to simplify is
not a proof of nonsingularity. \emph{Testing a coefficient with an ordinary
numerical heuristic and proceeding as though the result were an exact algebraic
fact is unsound.} A domain adapter should expose the conditions under which a
matrix result was obtained.

\subsection{The initial polynomial and its exact root counts}
\label{cas:sub-initial}
PA supplies a particularly useful computational interface. Write
\[
 P(a+Y)=\sum_{j=0}^n b_jY^j,\qquad
 \lambda=\min_{b_j\ne0}\{v(b_j)+j\rho\},
\]
and define the nonzero ordinary residual polynomial
\begin{equation}\label{cas:eq-initial}
 I_{a,\rho}(P)(X)=\st\!\left(t^{-\lambda}P(a+t^\rho X)\right).
\end{equation}
Its computation uses finite translation, exact coefficient valuations and
leading coefficients, exponent comparison and residue arithmetic. It need not
find any full root expansion.

For a complete root list $\alpha_i$ in a containing algebraically closed
workspace, PA proves
\begin{equation}\label{cas:eq-initialfactor}
 I_{a,\rho}(P)(X)=C\!\!\prod_{v(\alpha_i-a)\ge\rho}\!\!
 \left(X-\st\frac{\alpha_i-a}{t^\rho}\right),\qquad C\ne0 .
\end{equation}
Consequently its \textbf{degree} counts the roots in the closed valuation ball,
its \textbf{order at zero} counts the roots in the open ball, and the
\textbf{multiplicity of a residual root} $c$ counts the roots in the
corresponding residue-direction ball. These are exact finite conclusions even
when the represented roots have infinite supports \cite{cas:SourcePolynomial}.

The required valuations are effective for rational monomial coefficients by
\cref{cas:thm-core}. For algebraic or more general coefficients an additional
valuation and leading-coefficient service is needed; \emph{decidable field
equality alone should not be confused with an already implemented expansion
engine}. Algebraic root descriptors can remain valid while that service is still
being developed.

\subsection{Newton profiles without a real-valued plot}
\label{cas:sub-newton}
At centre zero the profile is
\[
 \phi_P(\rho)=\min_j\{v(a_j)+j\rho\},
\]
with candidate breakpoints at the finitely many values
$\rho=(v(a_j)-v(a_k))/(k-j)$ for $j\ne k$. At a breakpoint, the difference
between the largest and smallest active indices is the number of nonzero roots
at that valuation. Exact division by ordinary integers in a divisible exponent
group suffices; the exponents need \emph{not} be embedded in $\R$
\cite{cas:SourcePolynomial}.

A straightforward implementation enumerates candidates, sorts them by the
ordered-group comparator, and evaluates all affine expressions at each
candidate; a lower-hull implementation reduces repeated work. Drawing an
ordinary Euclidean Newton polygon is only a visualization for suitable rank-one
inputs. For exponents such as $1$ and $\omega$, a finite algebraic profile is
more faithful than a misleading numerical plot.

The general leading-valuation condition behind this is elementary: for
$P(X)=\sum_{j}a_jX^j$ with a nonzero root beginning $c\,t^\lambda$,
\begin{equation}\label{cas:eq-newtonmin}
 \min_j\bigl(v(a_j)+j\lambda\bigr)\quad\text{must be attained at least twice},
\end{equation}
since otherwise the unique term of least valuation could not cancel; the
coefficients attaining the minimum give the initial equation
$\sum_{j\in J}\lc(a_j)c^j=0$, whose nonzero roots are the candidate leading
coefficients. Substituting $X=t^\lambda(c+Y)$ and normalizing produces a smaller
perturbation problem; fractional $\lambda$ can force ramification. Two small
examples need no infinite search:
\[
 X^2-2X+1-t=0\ \Longrightarrow\ X=1\pm t^{1/2},
 \qquad X^2+t=0\ \Longrightarrow\ X=\pm i\,t^{1/2}.
\]
Both roots of the first have standard part $1$ but are distinct exact elements;
storing only standard parts loses multiplicity and branch information. The
second has no real roots with uncertain sign. These are the algebraic steps
behind Newton--Puiseux and generalized Newton-polygon methods, for which
grid-based treatments supply a developed framework \cite{cas:vdH}.

\begin{example}[The cubic]\label{cas:ex-cubic}
Let $P(z)=z^3-t^2z-t^5$. Its coefficient profile is $\min\{5,\,2+\rho,\,3\rho\}$.
The active indices at $\rho=1$ are $\{1,3\}$ and at $\rho=3$ are $\{0,1\}$, so
two roots have valuation $1$ and one has valuation $3$. The residual polynomials
are
\[
 I_{0,1}(P)=X^3-X,\qquad I_{0,3}(P)=-X-1,
\]
whose nonzero residual roots select the leading terms $t$, $-t$, $-t^3$. After
these simple residual branches have been selected, ordinary coefficient
recursion gives
\begin{align*}
 \alpha_+&=t+\tfrac12t^3-\tfrac38t^5+\Ov{t^7},\\
 \alpha_-&=-t+\tfrac12t^3+\tfrac38t^5+\Ov{t^7},\\
 \alpha_0&=-t^3-t^7+\Ov{t^{11}} .
\end{align*}
These finite displays do not replace the exact root descriptors; they are
certified expansions attached to them.
\end{example}

\begin{example}[The quartic and derivative-root counts]\label{cas:ex-quartic}
For $P(z)=z(z-t^3)(z-t)(z-1)$ the same machinery gives the residual weights and
polynomials at $\rho=0,1,3$, and the corresponding counts for $P'$. The general
statement from PA is that a ball containing $k\ge1$ root multiplicities contains
$k-1$ derivative-root multiplicities. The \textbf{occupation hypothesis} must be
retained: a ball containing no roots of $P$ can still contain a root of $P'$.
These routines count roots; they do not locate them.
\end{example}

\subsection{Hensel lifting, clusters, and algebraic fallback}
\label{cas:sub-hensel}
When a monic polynomial has pairwise coprime ordinary residual factors, PA
constructs unique lifted cluster factors: at each positive exponent the
correction is obtained by the same finite linear inverse, and the support lies
in the positive monoid generated by the input correction support
\cite{cas:SourcePolynomial}. With effective support decomposition this is an
implementable coefficient algorithm.

A \emph{multiple} residual root is different. For $z^m-t$ every root reduces to
zero and the first useful scale is $t^{1/m}$; a routine that only attempts
ordinary simple-root Newton iteration at zero will fail. It should split or
rescale the cluster, enlarge the exponent group when required, or return an
exact algebraic object with an unresolved expansion. The algebraic fallback is
valuable precisely because it separates exact solvability from the cost of
finding a convenient normal form. At multiple roots the correct interface may
return a finite \emph{cluster} object rather than arbitrarily numbered
individual branches: the preparation and coprime-factor lifting results retain
the support of \emph{factor coefficients}, while individual roots can require
fractional scales, and factor coefficients are often stable under perturbation
even when individual root labels are not.

Newton corrections $x_{k+1}=x_k-P(x_k)/P'(x_k)$ require a justification of the
invertibility of $P'(x_k)$ and of the precision improvement.

\begin{remark}[Cofinality is not free in higher rank]\label{cas:rem-cofinal}
There is \textbf{no unrestricted guarantee} that a natural-number sequence of
successive corrections becomes cofinal in every desired higher-rank cutoff: in a
higher-rank field, repeatedly doubling a positive error valuation need not reach
a cutoff in a strictly higher Archimedean class. Newton iteration is therefore a
coefficient-refinement method with a stated target policy, not a universal
convergence proof in the full surreal topology. PA's general lifting proof is a
well-founded support recursion. Classical Newton--Puiseux is a rank-one
specialization and should be identified as such rather than extrapolated to
arbitrary program-described Hahn coefficients
\cite{cas:Poonen,cas:SourcePolynomial}.
\end{remark}

\subsection{Root-perturbation thresholds: two different sharp theorems}
\label{cas:sub-rootperturbation}
Two of the sources give sharp perturbation results. They are \emph{not} two
versions of one theorem, they come from two different supplied manuscripts, and
they concern two different situations: univariate coefficient perturbation with
a multiplicity-driven exponent loss, and multivariate system perturbation with
an inverse-Jacobian-driven uniqueness ball. They are printed separately and must
not be merged into one ``root stability'' statement.

\subsubsection{Univariate coefficient perturbation (from PA)}
\label{cas:subsub-univariate}
For monic degree-$n$ polynomials with finite coefficients, PA proves that
coefficient errors of valuation at least $\varepsilon>0$ permit a matching of
root multisets with
\begin{equation}\label{cas:eq-rootmatch}
 v(\alpha_i-\beta_i)\ge\varepsilon/n .
\end{equation}
The exponent loss is \textbf{sharp}: compare $z^n$ with $z^n-t^\varepsilon$.
For a squarefree integral polynomial, put
\[
 d_i=v(P'(\alpha_i)),\qquad \delta_i=\max_{j\ne i}v(\alpha_i-\alpha_j).
\]
PA gives the sharper sufficient condition $\varepsilon>\max_i(d_i+\delta_i)$
with
\[
 v(\beta_i-\alpha_i)\ge\varepsilon-d_i ,
\]
and the easily stated coefficient-level sufficient condition
$\varepsilon>v(\Disc P)$ \cite{cas:SourcePolynomial}. These are certificates for
stable cluster labels, not heuristic tolerances. The strictness matters, and the
collision occurs exactly at the threshold:
\begin{equation}\label{cas:eq-collision}
 z(z-t^a)+\frac{t^{2a}}4=\left(z-\frac{t^a}2\right)^{\!2} .
\end{equation}
Before applying these estimates to infinite root locations, the polynomial must
be normalized to the stated finite-coefficient setting and the error transformed
with it.

\begin{remark}[Operational framing]\label{cas:rem-univariateops}
A software request for roots to valuation $\rho$ may require coefficient
information through approximately $n\rho$ near an unresolved multiple cluster.
\emph{Ignoring multiplicity overstates the precision.}
\end{remark}

\subsubsection{Multivariate system perturbation (from AG)}
\label{cas:subsub-multivariate}
For a simple zero $a$ of a nonnegative-support system, AG sets
\begin{equation}\label{cas:eq-kappa}
 \kappa_a=\max\!\left(0,\ -\min_{i,j}v\bigl((DF(a)^{-1})_{ij}\bigr)\right).
\end{equation}
For a perturbation $D$ with $\delta=\min_iv_H(D_i)>2\kappa_a$, AG proves a
\emph{unique} continued zero in the ball $v_{\min}(b-a)>\kappa_a$, with the
displacement bound
\begin{equation}\label{cas:eq-displacement}
 v_{\min}(b-a)\ge\delta-\kappa_a .
\end{equation}
These are exact inequalities in the full ordered value group
\cite{cas:SourceAG}.

The strict inequality is \textbf{sharp}: $z^2-t^{2\eta}$ perturbed by
$t^{2\eta}$ collapses two simple roots into a double root precisely at
$\delta=2\kappa_a$. \emph{The equality case cannot be accepted.}

\begin{remark}[Operational framing]\label{cas:rem-multivariateops}
A root tracker can use this as a certificate-producing operation: return the old
root descriptor, the perturbation bound, the inverse-Jacobian bound, and the new
implicit root with a verified uniqueness ball. It need not enumerate an infinite
Newton sequence. These guarantees remain conditional on obtaining the exact
bounds: \emph{a numerical estimate at one real specialization of the monomials
is not a proof of the higher-rank inequality}, and interval bounds for ordinary
coefficients cannot replace exponent-group comparisons.
\end{remark}

\subsection{Implicit analytic roots}
\label{cas:sub-implicit}
Suppose an ordinary analytic equation has $f(c)=0$ and $a=f'(c)\ne0$, and
consider $f\sh(z)+\tau g\sh(z)=0$ for a positive infinitesimal $\tau$. SA gives
the unique solution in the monad of $c$ as a support-controlled implicit germ,
whose first coefficients are
\begin{equation}\label{cas:eq-implicit}
 z=c-\frac{g(c)}{a}\tau
 +\left(\frac{g(c)g'(c)}{a^2}-\frac{f''(c)g(c)^2}{2a^3}\right)\tau^2+O(\tau^3).
\end{equation}
A CAS can store the implicit equation, the selected monad and a nonzero-Jacobian
certificate as the exact answer, computing coefficients on demand; the proof is
coefficient recursion, not a numerical root search in a large finite
approximation \cite{cas:SourceAnalysisTwo}.

More generally, for a formal system $\Phi(X,Y)=0$ with specified initial point
and invertible Jacobian in $Y$, each new homogeneous degree's coefficient is
determined by one fixed finite matrix inverse; when the coefficient operations
and input generators are effective, the solution's finite jets are computable.
A natural exact object is
\texttt{ImplicitSeries[equations, variables, initialData, certificate]}: the
equation and uniqueness certificate denote the whole solution, and a computed
jet is only a view of it. A returned solution should include the root centre,
the leading branch data, the Jacobian certificate, and support or precision
bounds. \emph{A singular Jacobian is not a proof that no solution exists}; it
means this simple-root procedure is inapplicable and a preparation, ramification
or finite-algebra method is needed.

\subsection{Near-tangency conditioning is algebraic information}
\label{cas:sub-tangency}
TR proves, for $0<\delta\in\mm$,
\begin{equation}\label{cas:eq-acos}
 q(\delta)=\arccos(1-\delta)
 =\sqrt{2\delta}\left(1+\frac{\delta}{12}+\frac{3\delta^2}{160}
 +\frac{5\delta^3}{896}+O(\delta^4)\right),
\end{equation}
and, for $e\ne0$ with $e\prec\delta$, the branch-stability formula
\begin{equation}\label{cas:eq-acosprecision}
 v\bigl(q(\delta+e)-q(\delta)\bigr)=v(e)-\tfrac12v(\delta).
\end{equation}
The inverse derivative therefore \emph{loses half the valuation of the tangency
parameter}. This is not a roundoff anecdote: it is an exact input--output
precision rule that a branch-aware algebraic or trigonometric evaluator can
propagate \cite{cas:SourceTrig}.

\section{Exact functions: formal germs, ordinary lifts, coherent data, radius-free germs}
\label{cas:sec-functions}

\subsection{Formal evaluation versus ordinary lift}
\label{cas:sub-formaleval}
Let $k\subseteq\C$ be an admitted coefficient field and $A(X)=\sum_{n\ge0}a_nX^n$
a specified formal power series over $k$. For an infinitesimal $\epsilon$ with a
certified well-ordered positive support, the evaluation
\begin{equation}\label{cas:eq-formaleval}
 A(\epsilon)=\sum_{n\ge0}a_n\epsilon^n
\end{equation}
is strongly summable: the support lies in $(\supp\epsilon)^*$ and finite word
multiplicity gives finitely many contributions at each exponent.
\textbf{No ordinary radius of convergence is required}
\cite{cas:SourceAnalysisThree}.

Thus both $\sum t^n/n!$ and $\sum n!\,t^n$ have exact Hahn interpretations, but
only the first is the Taylor series of an ordinary holomorphic function near
zero. The second may be stored as a formal-series value with a recurrence and a
support certificate. Calling both objects ordinary analytic germs would erase a
distinction on which the whole coherence theory depends.

\subsection{The Taylor lift, and what it does not confer}
\label{cas:sub-taylorlift}
For an ordinary holomorphic germ $f$ at $c\in\C$ --- with a branch specified
where necessary --- and $\epsilon\in\mm$, define
\begin{equation}\label{cas:eq-lift}
 f\sh(c+\epsilon)=\sum_{n\ge0}\frac{f^{(n)}(c)}{n!}\epsilon^n .
\end{equation}

\begin{proposition}[Correctness of infinitesimal substitution]\label{cas:prop-lift}
Suppose a function descriptor specifies an ordinary analytic germ at $c$, an
exact centre, and an effective Taylor-coefficient procedure, and let $\epsilon$
have certified well-ordered positive support $E$. Then \eqref{cas:eq-lift} is a
strong Hahn sum with support contained in $E^*$, and supplies an exact symbolic
value. If the support and coefficient dependencies are effective, each requested
finite coefficient calculation terminates. Sums, products and compositions of
such analytic germs are respected whenever the ordinary centres match and the
relevant germs are defined. \textbf{Equality of arbitrary resulting values is
not guaranteed by this construction alone.}
\end{proposition}
\begin{proof}
Every coefficient contribution is indexed by a finite word in $E$, and the
positive-support lemma of \cref{cas:sub-twoconditions} gives strong
summability, so the semantic existence is immediate; the descriptor fixes every
coefficient, hence a unique value. At a fixed exponent the formal Taylor
identities for sums, products and composition involve only finitely many
contributing terms and therefore remain valid after substitution. Centre values
are handled as ordinary coefficients before the infinitesimal expansion is
performed. Each specified effective coefficient request uses finitely many
certified dependencies, and \cref{cas:thm-grid} turns this into a coefficient
algorithm when the supports are given as effective grids. The final
qualification follows from \cref{cas:thm-zerotest}: a general formal-series
constructor cannot automatically acquire a complete zero test.
\end{proof}

The domain is $U\sh=\{c+\epsilon:c\in U,\ \epsilon\in\mm\}$, \emph{not} every
surcomplex number whose printed formula resembles an ordinary argument. Ordinary
values such as $f(c)$ are evaluated or named first; it is not correct to insert
a nonzero ordinary $c$ directly into a Hahn power series about zero and invoke
ordinary convergence. Local instances:
\begin{align}
 \exp\sh(c+\epsilon)&=e^c\sum_{n\ge0}\frac{\epsilon^n}{n!},
 \label{cas:eq-localexp}\\
 \log\sh(c+\epsilon)&=\log(c)+\sum_{n\ge1}\frac{(-1)^{n+1}}{n}
 \left(\frac{\epsilon}{c}\right)^{\!n},\qquad c\ne0,
 \label{cas:eq-locallog}\\
 \sin\sh(c+\epsilon)&=\sin(c)\cos\sh(\epsilon)+\cos(c)\sin\sh(\epsilon).
 \label{cas:eq-localsin}
\end{align}
The logarithm's ordinary germ must be chosen; there is no holomorphic logarithm
germ at zero. At $c=0$ the exponential coefficients lie in $\Q$, whereas at an
algebraic nonzero $c$ the coefficient $e^c$ may require a transcendental
constant descriptor: domain closure must be checked at both the coefficient and
the exponent level. For finite real $r+\eta$ one likewise evaluates the ordinary
constants first,
\begin{align}
 \sin(r+\eta)&=\sin r\sum_{n\ge0}\frac{(-1)^n\eta^{2n}}{(2n)!}
 +\cos r\sum_{n\ge0}\frac{(-1)^n\eta^{2n+1}}{(2n+1)!},
 \label{cas:eq-sinfinite}\\
 \cos(r+\eta)&=\cos r\sum_{n\ge0}\frac{(-1)^n\eta^{2n}}{(2n)!}
 -\sin r\sum_{n\ge0}\frac{(-1)^n\eta^{2n+1}}{(2n+1)!}.
 \label{cas:eq-cosfinite}
\end{align}

The construction is also more general than its simplifier. The expression
$\sin\sh(\pi+t)=-\sin\sh(t)$ is exact once the ordinary constants and functions
are admitted, but a minimal $\RA$-coefficient implementation does not thereby
gain a complete decision procedure for every ordinary constant generated by
arbitrary special-function evaluation. \emph{A value constructor may be more
general than its simplifier.}

\subsection{Which familiar functions can be included}
\label{cas:sub-admissible}
At regular ordinary centres the mechanism applies to elementary analytic
functions and to specified branches of special functions: exponential,
logarithm away from its selected cut, trigonometric functions, gamma away from
its poles, error functions, Airy functions, Bessel functions with their
parameter and branch conditions, and hypergeometric functions at regular
parameter and argument values. This is a rule about known ordinary germs, not an
assertion that every possible parameter value or singularity is regular.

Ordinary CAS differentiation and series tools are valuable coefficient engines
and supply exact special-function expressions, derivatives and local expansions
\cite{cas:WLSpecial}. A successful coefficient call does not by itself certify
an entire surreal support or a global function domain, and the extension layer
must preserve that missing information. Calling \texttt{SeriesCoefficient} on a
supported classical function can supply coefficients, but a returned unevaluated
expression is not a termination guarantee for all requested indices, and
user-supplied recurrences require their own totality and domain assumptions.

\paragraph{Coefficient-domain promotion.}
$\Gamma\sh(2+t)$ is an exact value given by the ordinary germ at $2$. Its
coefficients may involve constants such as Euler's constant and derivatives of
gamma, which lie \emph{outside} the real-algebraic coefficient core. The system
should promote to an exact ordinary-constant domain with correspondingly weaker
general equality guarantees, not silently call those constants algebraic.

Many useful special-function germs can instead be specified by a finite
differential equation or recurrence, initial values and a chosen ordinary
domain. A linear differential equation at an ordinary point provides recursive
Taylor coefficients; algebraic equations provide implicit-function recurrences
at simple branches. Parameter singularities, resonance, irregular singular
points and branch changes must be part of the input checks.

\subsection{Singular centres and the summability boundary}
\label{cas:sub-singular}
At an ordinary pole one uses a finite-principal-part Laurent germ on a punctured
infinitesimal neighbourhood, factoring the meromorphic germ into a finite
negative power and a holomorphic germ; evaluation then requires a nonzero
displacement. At an algebraic branch point one uses a ramified coordinate and a
selected Puiseux branch. Formula \eqref{cas:eq-acos} is a basic example: the
ordinary Taylor expansion of $\arccos$ at $1$ is not an analytic power series in
$\delta$, but becomes a convergent ordinary series after the positive
square-root uniformizer is introduced.

At an \emph{essential} singularity a Laurent expression can fail Hahn
summability outright. If $\epsilon\ne0$ is infinitesimal, the terms
$\epsilon^{-n}/n!$ have strictly \emph{decreasing} valuations, so
\[
 \sum_{n\ge0}\frac{\epsilon^{-n}}{n!}
\]
is \textbf{not} a strong Hahn definition of $e^{1/\epsilon}$. SA explicitly
distinguishes this failed substitution from a separately selected global
exponential evaluated at $1/\epsilon$
\cite{cas:SourceAnalysisTwo,cas:SourceAnalysisThree}. A separate global real
exponential or surcomplex extension may assign a value, but that is a different
construction.

\subsection{Formal, analytic, and asymptotic descriptors are different}
\label{cas:sub-descriptors}
A formal series with computable coefficients can be meaningful at infinitesimal
inputs even when it is not an ordinary analytic germ, and an asymptotic
expansion of an ordinary function is a third kind of data. A general smooth
function does not have a unique infinitesimal extension determined by its Taylor
coefficients: an ordinary smooth function flat at zero has the same Taylor jet
as the zero function there, and assigning only that jet to all infinitesimals
discards information about its nonzero values away from zero. CT explicitly
separates analytic target coefficients from merely smooth parameter
coefficients, and a CAS should retain the same distinction
\cite{cas:SourceCT}. See also \cref{cas:sub-asymptotic}.

\subsection{Hahn-coherent data on a fixed common ordinary domain}
\label{cas:sub-coherent}
The first of the two coherence certificates in this report is the one all three
sources define.

\begin{definition}[Common-domain Hahn-coherent datum]\label{cas:def-coherent}
A Hahn-coherent datum on a \emph{fixed} ordinary domain $U\subseteq\C$ is
\begin{equation}\label{cas:eq-coherent}
 F=\sum_{\gamma\in S}f_\gamma t^\gamma,\qquad
 f_\gamma\in\OO(U),\quad S\text{ well ordered},
\end{equation}
where $\OO(U)$ denotes the ordinary holomorphic functions on $U$, with
\emph{one} common well-ordered support and \emph{one} common domain. Its value
at $c+\epsilon\in U\sh$ is
\begin{equation}\label{cas:eq-coherenteval}
 F\sh(c+\epsilon)=\sum_{\gamma\in S}\sum_{n\ge0}
 t^\gamma\frac{f_\gamma^{(n)}(c)}{n!}\epsilon^n ,
\end{equation}
with source support bound $S+(\supp\epsilon)^*$.
\end{definition}

Coefficient functions must share the stated ordinary domain. \textbf{Merely
holding an unrelated germ at each centre does not provide coherent data}, and
pointwise knowledge that selected values are infinitesimal is no substitute for
the uniform support and domain condition
\cite{cas:SourceAnalysisThree,cas:SourceResearch}. A computational object should
therefore contain both a support description and an ordinary analytic domain
descriptor, plus a coefficient-function provider. Finite support with rational
or explicitly holomorphic coefficient functions is a particularly accessible
initial implementation; infinite coherent data need further operational
certificates, just as infinite scalar series do.

HC's Hartogs and gluing theorems are mathematically useful but \textbf{should
not be exposed as unconditional algorithms} that discover a global support or
decide arbitrary holomorphic identities. To invoke them computationally, the
input must contain the relevant function descriptions and hypothesis
certificates, or a restricted syntax in which those hypotheses can be checked.

\subsection{The radius-free Hahn-germ algebra: a strictly weaker certificate}
\label{cas:sub-radiusfree}
The second certificate is present in only one of the three sources, imported
from AG, and it is \emph{strictly weaker} than
\cref{cas:def-coherent}. The two must carry distinct names, and a system must
not silently identify them.

\begin{definition}[Radius-free Hahn germ]\label{cas:def-radiusfree}
Put
\begin{equation}\label{cas:eq-germalgebra}
 R_n=\C\{z_1,\dots,z_n\},\qquad \Acal_n=R_n((t^\Gamma)).
\end{equation}
Each coefficient is an ordinary \emph{convergent germ at zero}, but different
coefficients need not share any common ordinary radius.
\end{definition}

A common-domain family in $\OO(U)((t^\Gamma))$ is a stronger and different
certificate \cite{cas:SourceAG}. A finite computational descriptor must
therefore \textbf{separate its Hahn-support program from its ordinary
coefficient-germ program}.

The distinction has executable consequences. Let
\begin{equation}\label{cas:eq-radiusfree}
 H(z)=\sum_{n\ge1}\frac{t^n}{1-nz}.
\end{equation}
Every coefficient is a valid ordinary germ at zero, and there is \emph{no}
positive ordinary disk on which all of them are holomorphic, since their poles
are $1/n$. Nevertheless $H(a)$ is well defined for every infinitesimal $a$, by
strong Taylor substitution. \textbf{Calling an ordinary contour routine on a
fixed disk containing all coefficient domains is an error}; infinitesimal
rescaling is the correct alternative (\cref{cas:sub-rescaling}). Substituting
$a=t^2$,
\begin{equation}\label{cas:eq-radiusfreesub}
 H(t^2)=\sum_{n\ge1}\frac{t^n}{1-nt^2}
 =t+t^2+2t^3+3t^4+5t^5+9t^6+16t^7+31t^8+\cdots,
\end{equation}
and at exponent $m$ the coefficient is the finite integer sum
\begin{equation}\label{cas:eq-radiusfreecoeff}
 [t^m]H(t^2)=\sum_{\substack{n\ge1,\ k\ge0\\ n+2k=m}}n^k .
\end{equation}
The coefficient algorithm terminates with an explicit finite dependency set even
though the coefficient germs share no radius. The values for $m=1,\dots,16$ are
\[
 1,\ 1,\ 2,\ 3,\ 5,\ 9,\ 16,\ 31,\ 61,\ 125,\ 266,\ 579,\ 1305,\ 3009,\ 7120,\ 17255,
\]
independently checked by symbolic computation (\cref{cas:sec-executed}). CT
identifies the same value as a microscopic Cauchy integral at an appropriate
radius \cite{cas:SourceCT}. \emph{None of the three delivered prototypes
represents the outer infinite sum}; \eqref{cas:eq-radiusfreecoeff} describes a
straightforward next lazy-series module.

AG also distinguishes $\Acal_n$ from $K[[z]]$. The formal expression
\begin{equation}\label{cas:eq-notinA1}
 (z-t^\eta)^{-1}=-\sum_{j\ge0}t^{-(j+1)\eta}z^j
\end{equation}
is a power series over $K$ but is \textbf{not} an element of $\Acal_1$: its
combined Hahn exponents are not well ordered. In $\Acal_1$, the function
$z-t^\eta$ actually \emph{vanishes at an infinitesimal point}, so it must not be
marked as an invertible germ on the whole monad. A CAS needs separate types for
a formal completion at a point and a function algebra on the entire
infinitesimal neighbourhood \cite{cas:SourceAG}.

\subsection{Composition contracts and the blind-substitution failure}
\label{cas:sub-composition}
Formal power-series composition $A(B(X))$ is immediately defined when the inner
constant term is zero, because each resulting coefficient depends on finitely
many terms. A nonzero inner constant generally requires an analytic
re-expansion or a separate summability condition; at finite surreal inputs the
ordinary-centre Taylor lift supplies that re-expansion when the ordinary
analytic domains match.

In the coherent category, for $F=f_0+E$ with nonnegative support, $E$ of
positive support, and $f_0(U)\subseteq V$, the sources define
\begin{equation}\label{cas:eq-composition}
 H\circ F=\sum_{\delta,n}\frac{h_\delta^{(n)}\circ f_0}{n!}\,t^\delta E^n,
 \qquad H=\sum_\delta h_\delta t^\delta\in\HH(V),
\end{equation}
with support contained in $\supp H+(\supp E)^*$. This common-domain condition is
stronger than pointwise existence. For a transseries function at infinity,
admissible composition uses its large-argument theory instead. These contracts
must not be collapsed into one unchecked substitution routine. An executable
implementation needs a verified ordinary-domain inclusion, ordinary derivatives
and compositions, and effective support bookkeeping; a symbolic domain
obligation can remain in a conditional result when it cannot be resolved
automatically.

\paragraph{Blind substitution fails.}
The geometric expansion $\sum_{n\ge0}t^nz^n$ is coherent on the finite plane and
equals $(1-tz)^{-1}$ there. Substituting $z=t^{-2}$ into the expansion produces
exponents $-n$, which are \emph{not} well ordered. The rational expression still
has a value at $t^{-2}$, but its valid expansion is now
\begin{equation}\label{cas:eq-blindsub}
 \frac1{1-t^{-1}}=-\frac{t}{1-t}=-\sum_{n\ge1}t^n ,
\end{equation}
and at $z=t^{-1}$ there is a pole. A failed series substitution therefore
requires re-expansion, a different representation, or rejection, depending on
the actual argument.

\subsection{Evaluation, translation, and finite jets}
\label{cas:sub-jetsgerms}
For $F=\sum_\gamma f_\gamma t^\gamma$ and an infinitesimal tuple $a$, AG defines
\begin{align}
 F(a)&=\sum_{\gamma,\alpha}t^\gamma
 \frac{\partial^\alpha f_\gamma(0)}{\alpha!}a^\alpha,
 \label{cas:eq-germeval}\\
 (\tau_aF)(z)&=\sum_{\gamma,\alpha}t^\gamma
 \frac{\partial^\alpha f_\gamma(z)}{\alpha!}a^\alpha,
 \label{cas:eq-germtranslation}
\end{align}
with support certificate $S+E_a^*$, where $S=\supp F$ and $E_a$ is the union of
the coordinate supports of $a$. A coefficient request can be compiled into a
finite dependency graph when those supports have effective presentations and the
ordinary germ derivative requests terminate. This gives an actual algorithmic
interpretation of AG's formula, \emph{not} an assumption that every convergent
germ is already a CAS object.

Finite jets are particularly accessible: at a represented infinitesimal $a$,
\begin{equation}\label{cas:eq-finitejets}
 \Acal_n/(z-a)^N\ \simeq\ K[X_1,\dots,X_n]/(X_1,\dots,X_n)^N .
\end{equation}
Once the required finite derivatives have been computed, the right-hand side
uses ordinary sparse polynomial arithmetic with a degree cutoff, and a separate
Hahn precision cutoff may be applied to its coefficients. \textbf{The Taylor
degree and the Hahn valuation are different axes of approximation and must not
be compressed into one integer option.}

\subsection{Preparation and division as operator algorithms}
\label{cas:sub-weierstrass}
AG's one-variable preparation and its several-variable division formulas are
useful implementation specifications. For $F=f+E$ with leading ordinary germ $f$
regular in a selected variable $w$ and $\vH(E)>0$, and with $D$ and $\Pi$ the
ordinary Weierstrass quotient and remainder operators,
\begin{equation}\label{cas:eq-weierstrass}
 Q=(\id+DM_E)^{-1}DG=\sum_{r\ge0}(-DM_E)^rDG,\qquad
 R=\Pi(G-EQ),
\end{equation}
with support contained in $\supp G+(\supp E)^*$. These are exact formulas from
AG \cite{cas:SourceAG}, \textbf{not} an implemented routine in any of the three
delivered packages.

Three executability caveats survive and all three matter.
\begin{itemize}
\item The system needs a \emph{concrete ordinary division backend}. For
polynomial or algebraic coefficient germs, polynomial reduction and selected
local algebra procedures are natural choices; for a suitable differential or
analytic descriptor, additional constructive division methods may be possible.
Even a finite output coefficient can call an ordinary operation unavailable for
the given germ language.
\item \emph{An arbitrary choice-based complex-linear splitting is not
executable} merely because it preserves the space of analytic germs. AG's
arbitrary existence of a linear splitting is sufficient for its theorem but not
for a program.
\item The ordinary leading order and regularity direction must be found or
supplied, and \emph{if detecting whether a coefficient germ is zero is only
semidecidable, searching for the first nonzero homogeneous part need not
terminate on zero input}. A reliable API accepts certified leading data, tries
supported methods, and otherwise returns a structured unresolved result rather
than asserting a Weierstrass polynomial.
\end{itemize}

\subsection{A nonpolynomial algebraic cluster: two developments of \texorpdfstring{$\sin^2z=t$}{sin\textasciicircum 2 z = t}}
\label{cas:sub-sinsquared}
Two of the sources develop the same example in two compatible ways, and both
survive.

\paragraph{Local monad development.}
$\sin^2z-t$ has exactly two zeros in the monad of zero, counted with
multiplicity:
\[
 z_\pm=\pm\arcsin\sh(\sqrt t)
 =\pm\left(t^{1/2}+\frac{t^{3/2}}{6}+\frac{3t^{5/2}}{40}
 +\frac{5t^{7/2}}{112}+\cdots\right),
\]
representable locally as two analytic-root descriptors after ramification. The
exact prepared cluster polynomial is
\begin{equation}\label{cas:eq-prepared}
 P(z)=z^2-\bigl(\arcsin\sh(\sqrt t)\bigr)^2 .
\end{equation}
It is polynomial in $z$, but its coefficients lie in an \emph{analytic
extension} of the minimal rational monomial field.

\begin{remark}[Algebraicity is relative to the ambient domain]
\label{cas:rem-algebraicity}
A package must not label these roots ``algebraic over the original coefficient
field'' merely because they occur in a prepared polynomial over a newly enlarged
field. The ambient coefficient domain is part of every algebraicity claim.
\end{remark}

\paragraph{Full-strip development.}
On the canonical strip modulo $2\pi$, putting $U=E_{\HH}(iz)$ (see
\cref{cas:sub-strips}) turns the equation into
\begin{equation}\label{cas:eq-stripquartic}
 -U^4+(2-4t)U^2-1=0 ,
\end{equation}
which has four roots counted with multiplicity; the two other angular roots lie
near $\pi$ modulo $2\pi$. A local monad solver and a full-strip periodic solver
are answering different, compatible questions.

\section{Finite quotient algebras and residues}
\label{cas:sec-quotients}

\subsection{Remainders, traces, and norms without choosing roots}
\label{cas:sub-remainders}
Let $P\in K[z]$ be monic of degree $n$ over an effective surcomplex coefficient
domain $K$, and put $A_P=K[z]/(P)$. Polynomial remainders represent elements of
this $n$-dimensional algebra, and multiplication matrices give traces, norms,
characteristic polynomials and resultants \emph{without choosing individual
roots}. The residue functional of PA is
\begin{equation}\label{cas:eq-residuefunctional}
 \lambda_P(H)=[z^{n-1}](H\bmod P),
\end{equation}
with trace identity $\Tr(m_H)=\lambda_P(P'H)$. Both are finite exact coefficient
calculations \cite{cas:SourcePolynomial}.

\subsection{Residue pairing versus trace pairing}
\label{cas:sub-pairings}
The pairing $(G,H)\mapsto\lambda_P(GH)$ has determinant $(-1)^{n(n-1)/2}$ in the
power basis, \textbf{even at repeated roots}. The trace pairing, by contrast,
has determinant $\Disc P$ and can degenerate. This is an algorithmic
distinction, not a curiosity: a residue calculation by remainders stays regular
through a root collision even where formulas using the individual weights
$1/P'(\alpha_i)$ develop apparent singularities.

The design policy that follows is to \emph{compute symmetric invariants,
remainders and finite-algebra traces before requesting labelled root
expansions}.

\subsection{An analytic observable in a polynomial quotient, by two routes}
\label{cas:sub-observable}
For $P=z^2-t$, the relation $z^2=t$ gives the formal reduction
$e^z\bmod P=A(t)+zB(t)$ with
\[
 A(t)=\sum_{n\ge0}\frac{t^n}{(2n)!},\qquad
 B(t)=\sum_{n\ge0}\frac{t^n}{(2n+1)!},
\]
both strongly summable at infinitesimal $t$. Hence
\begin{equation}\label{cas:eq-residueexample}
 \lambda_P(e^z)=B(t)=1+\frac t6+\frac{t^2}{120}+\frac{t^3}{5040}+\cdots
 =\frac{\sinh\sh(\sqrt t)}{\sqrt t}.
\end{equation}

The same value arises by the contour and residue route. With
$M(z)=e^z/(z^2-t)$ and $s=t^{1/2}$, the poles $\pm s$ have residues $e^s/(2s)$
and $-e^{-s}/(2s)$, each of which is \emph{infinite}; their sum is nevertheless
the exact finite value
\begin{equation}\label{cas:eq-residue}
 \frac1{2\pi i}\HInt_{|z|=1}M(z)\,dz=\frac{\sinh s}{s}
 =\sum_{n\ge0}\frac{t^n}{(2n+1)!}
\end{equation}
\cite{cas:SourceAnalysisTwo,cas:SourceAnalysisThree}. Here the exponential on
the infinitesimal poles is the canonical lift and the contour is the ordinary
circle with the coefficientwise functional of \cref{cas:sub-coefficientwise}.

Both routes survive because they teach different lessons. The quotient
computation avoids both a branch choice for the intermediate square root and the
subtraction of two infinite leading residue weights. The contour computation
shows that a correct CAS must combine the exact residue expressions
\emph{before} applying standard part or coarse truncation: treating the
individual residues as generic infinities produces an indeterminate expression
and destroys a computable answer, whereas keeping their algebraic dependencies
makes the cancellation automatic. For more general analytic observables an exact
remainder constructor may be supported by the coherent division theorem and a
certificate rather than by a finite polynomial reduction alone.

\subsection{Normal forms, multiplication matrices, and the perturbed residue}
\label{cas:sub-normalforms}
For a leading isolated complete intersection $f=(f_1,\dots,f_n)$ with
$\mu=\dim_\C R_n/(f_1,\dots,f_n)<\infty$ and a positive-support perturbation
$F=f+E$, AG proves preservation of the finite algebra length and supplies
multiplication matrices, residues and trace identities \cite{cas:SourceAG}. This
is an attractive CAS target because the result is represented by finitely many
matrices even when their entries are exact lazy Hahn series.

Choose a monomial basis $b_1,\dots,b_\mu$ of the leading quotient and write an
ordinary splitting $G=i\,p(G)+\sum_jf_jh_j(G)$ with $pi=\id$ and $hi=0$. For an
effective implementation, $p$ and the $h_j$ must be \emph{algorithms} for the
chosen coefficient-germ presentation; AG's arbitrary existence of a linear
splitting is sufficient for its theorem but not for a program. With
$T=\sum_jE_jh_j$, the normal-form coordinates are
\begin{equation}\label{cas:eq-NF}
 \NF_F(G)=p(\id+T)^{-1}G=\sum_{r\ge0}(-1)^r\,pT^rG ,
\end{equation}
and the matrices of multiplication are
\begin{equation}\label{cas:eq-multmatrices}
 (M_H)_{rs}=\NF_F(Hb_s)_r .
\end{equation}
A finite-grid support backend and a constructive ordinary splitting give
effective coefficient access to every entry; matrix multiplication, determinant,
trace and characteristic polynomial then reduce to the corresponding scalar
operations, none of which requires isolating any perturbed root.

The perturbed residue is
\begin{equation}\label{cas:eq-perturbedresidue}
 \Res_F(H)=\sum_{k\in\N^n}(-1)^{|k|}
 \res_{f^{k+1}}\!\bigl(HE_1^{k_1}\cdots E_n^{k_n}\bigr),
\end{equation}
with support in $\supp H+E^*$. At a requested exponent an effective dependency
routine reduces the calculation to finitely many ordinary residues; when the
leading denominators are coordinate powers those ordinary residues are simply
Taylor coefficient extraction, and more general finite-dimensional residues have
established algebraic computation methods \cite{cas:CDS}. The trace identity is
\begin{equation}\label{cas:eq-traceidentity}
 \Res_F(HJ_F)=\Tr(M_H),\qquad
 J_F=\det\!\left(\frac{\partial F_i}{\partial z_j}\right),
\end{equation}
giving both a computation and a consistency check; for $H=1$ the answer is the
ordinary integer $\mu$. A finite-algebra object should record basis,
multiplication matrices, defining relations, a residue vector or pairing matrix,
and enough reduction data to verify the relations --- a representation that
retains the multiplicities and nilpotents a list of distinct points would lose.

\begin{remark}[Verification is not installation]\label{cas:rem-notrewrite}
Verifying the Jacobian trace identity in a worked example does not install it as
a universal rewrite rule for all coherent quotients, and HC distinguishes its
proved multidimensional residue pairing from a general identification of every
localized functional with a chosen Grothendieck-residue theory. A software
implementation should retain that boundary.
\end{remark}

\subsection{Length preservation and the length-drop counterexample}
\label{cas:sub-length}
AG's length-preservation theorem concerns $n$ equations in $n$ variables with an
isolated leading complete intersection. \textbf{An overdetermined system is not
covered.} AG's own example $(z^2,0)\mapsto(z^2,t^\eta z)$ drops the length from
two to one. An implementation must \emph{test or demand} the
isolated-complete-intersection certificate rather than applying a conservation
rule to every infinitesimal perturbation \cite{cas:SourceAG}.

\subsection{The coupled system at two independent scales}
\label{cas:sub-coupled}
Both articles that treat this example imported it from their own supplied
manuscript sets --- one from HC, the other from AG and CT --- and they present
it at different generality. The two-scale form is strictly more informative and
its higher-rank instantiation is the whole point, so we state it as the general
case and record the single-scale specialization separately.

\paragraph{The general two-scale form.}
Consider
\begin{equation}\label{cas:eq-fourexample}
 F_1=x^2-Ay,\qquad F_2=y^2-Bx,\qquad A,B\ne0\ \text{infinitesimal},
\end{equation}
with basis $(1,x,y,xy)$ and $q=AB$. The coordinate multiplication matrices are
\begin{equation}\label{cas:eq-fourmatrices}
 M_x=\begin{pmatrix}0&0&0&0\\1&0&0&q\\0&A&0&0\\0&0&1&0\end{pmatrix},
 \qquad
 M_y=\begin{pmatrix}0&0&0&0\\0&0&B&0\\1&0&0&q\\0&1&0&0\end{pmatrix} .
\end{equation}
They commute and satisfy $M_x^2=AM_y$ and $M_y^2=BM_x$, with characteristic
polynomials
\[
 Q_x(T)=T(T^3-A^2B),\qquad Q_y(T)=T(T^3-AB^2).
\]
All of these are finite polynomial calculations. The four roots are the origin
and
\[
 (\rho\zeta^j,\ \rho^2\zeta^{2j}/A),\qquad j=0,1,2,\qquad \rho^3=A^2B,
\]
for $\zeta$ a primitive cube root of unity, with nonzero coordinate valuations
$(2v(A)+v(B))/3$ and $(v(A)+2v(B))/3$. These roots require an algebraic or
exponent extension \emph{even though the multiplication matrices do not};
keeping the finite algebra avoids forcing that extension before computing traces
or total residues.

The Jacobian is $J_F=4xy-AB$, and the simple residue weights for the constant
numerator are
\[
 -\frac1{AB},\quad \frac1{3AB},\quad \frac1{3AB},\quad \frac1{3AB},
\]
whose sum is \textbf{exactly zero} although each individual weight is infinite,
while $\Res_F(xy)=1$ and $\Res_F(J_F)=4$. A truncation or floating-point
specialization that loses cancellation can therefore fail on a quantity with an
especially simple exact answer. For a monomial numerator the closed formula is
\begin{equation}\label{cas:eq-monomialresidue}
 \Res_F(x^py^r)=
 \begin{cases}
 A^{(2p+r-3)/3}B^{(p+2r-3)/3},&
 \tfrac{2p+r-3}{3},\ \tfrac{p+2r-3}{3}\in\N,\\[0.3em]
 0,&\text{otherwise.}
 \end{cases}
\end{equation}
Taking $A=t$ and $B=u=t^\omega$ preserves the higher-rank distinction without
changing the matrix algorithms.

\paragraph{The single-scale specialization.}
Setting $A=B=t$ gives $F_1=x^2-ty$, $F_2=y^2-tx$, with the same basis, dimension
four, roots $(0,0)$ and $(t\zeta,t\zeta^2)$ for $\zeta^3=1$, matrix entries in
$t$ and $t^2$, and characteristic polynomial $X^4-t^3X$. \emph{This is the form
in which the corresponding symbolic identities were actually executed} in the
layered-representations archive (\cref{cas:sec-executed}). The residue pairing
there has nonzero constant determinant despite infinite individual evaluation
weights.

Both residue-pairing results survive because they are different: the
residue-versus-trace determinant distinction of \cref{cas:sub-pairings}, and the
Gram matrix of determinant $1$ together with the vanishing of the summed
infinite weights above.

More generally, HC treats coupled coherent systems $F_i=w_i^{m_i}+E_i$ with
positive-support errors on one ordinary product domain, constructing a
finite-free algebra of rank $\prod_im_i$ with normal forms supported in $S+E^*$
\cite{cas:SourceResearch}. This motivates a lazy
\texttt{CoherentQuotientAlgebra} constructor whose multiplication matrices are
computed coefficientwise. It is \textbf{not} a generic theorem for arbitrary
analytic systems or arbitrary ordinary leading complete intersections; the
source explicitly delimits those extensions.

\section{Integration and contour operations}
\label{cas:sec-integration}

\subsection{Primitives by category}
\label{cas:sub-primitives}
A formal analytic germ $\sum a_nX^n$ has the normalized formal primitive
$\sum_{n\ge0}a_nX^{n+1}/(n+1)$, representable by a new coefficient generator
with the same local semantic obligations. A Laurent germ with finite principal
part is handled termwise, except that an $X^{-1}$ term requires a logarithm
descriptor or remains a formal logarithmic primitive; the residue is simply the
coefficient of $X^{-1}$. In a formal power-series ring of characteristic zero,
termwise integration with zero constant term is effective when the coefficients
are.

Rational integration in a formal variable $z$ over a characteristic-zero
coefficient field is a well-defined algebraic problem: partial fractions show
that polynomial terms and higher-order poles have rational primitives while a
simple pole contributes a logarithmic obstruction, so a rational differential
has a rational primitive exactly when all its simple-pole residues vanish, after
passing to an algebraic splitting field as necessary. For rational functions
over an effective complexified core, polynomial division and finite partial
fractions give effective residue computations once algebraic roots and
multiplicities are available. An antiderivative need not remain rational;
logarithms and their branch information must be added.

For common-domain Hahn families, SA defines differentiation and integration
coefficientwise, and a \textbf{coherent primitive exists precisely when the
ordinary closed-cycle Hahn periods vanish}. The same manuscript exhibits
fine-locally constant functions with zero derivative but different values on
different fine-open pieces. \emph{There can therefore be no unrestricted
endpoint integral satisfying the fundamental theorem for every fine-locally
analytic function} \cite{cas:SourceAnalysisThree,cas:SourceCT}.

A proposed integration API should distinguish a rational primitive, an
analytic-germ primitive, a coherent primitive on a stated domain, and an
antiderivative for a specified surreal derivation. \textbf{A successful result
in one category is not a result in any other category.} Additive constants can
be normalized locally, but a global primitive also needs compatibility across
the relevant domain. In a broader analytic coefficient domain a definite
integral may be stored exactly without a closed-form reduction: existence,
representability by an integral node, and reduction to elementary functions are
separate outcomes. Closure of the whole surreal field under a chosen
antiderivation is not a universal finite algorithm on arbitrary inputs.

\subsection{The coefficientwise contour functional, and what does not justify it}
\label{cas:sub-coefficientwise}
For an ordinary piecewise smooth path $\gamma$ in $U$ and common-domain Hahn
data the sources define
\begin{equation}\label{cas:eq-contour}
 \HInt_\gamma G(z)\,dz=\sum_\lambda t^\lambda\int_\gamma g_\lambda(z)\,dz ,
\end{equation}
with the result's support contained in that of $G$. Each ordinary coefficient
integral can be evaluated by the underlying CAS when possible, or retained as an
exact integral descriptor; an effective infinite coefficient family again needs
an effective coefficient procedure. On a simply connected ordinary domain,
normalized primitives are constructed coefficientwise; on a multiply connected
domain, periods matter. Nonstandard endpoints add SA's Taylor corrections to
ordinary endpoint integrals.

\begin{remark}[Scope of \eqref{cas:eq-contour}]\label{cas:rem-contourscope}
This is \textbf{the supplied manuscripts' functional}, not a newly asserted
full-class Riemann integral, and it is explicitly \emph{not} justified by taking
Riemann sums in the surreal fine topology: that topology does not make ordinary
nonconstant paths into fine-continuous paths, and the contour in
\eqref{cas:eq-contour} is not a fine-continuous real-parameter path through the
surcomplex class \cite{cas:SourceAnalysisTwo,cas:SourceCT}. Turning
\eqref{cas:eq-contour} into a computable scalar coefficient needs an ordinary
exact or validated integration method for that coefficient class;
\emph{exchanging a general numerical quadrature result for an exact coefficient
would change the contract}. A pole on the contour is not a zero winding number;
it is outside the definition unless a separate regularization has been selected.
The system must not infer a Cauchy theorem merely from the fact that a path is
closed: null-homology or an appropriate residue calculation is necessary, and
the data type of a contour must include orientation and pole-avoidance
conditions.
\end{remark}

Costin and Ehrlich's work on surreal extension and integration is established
background for a large class of resurgent functions and a relevant published
direction \cite{cas:CostinEhrlich}. It is \textbf{explicitly not identified}
here with the local Taylor-lift interface of \cref{cas:sub-taylorlift} or with
the coefficientwise functional \eqref{cas:eq-contour}, and it is not a license
to identify every formal asymptotic series with a unique global surcomplex
function.

\subsection{Algebraic rational periods and winding by certified crossings}
\label{cas:sub-periods}
For a rational differential and a finite semialgebraic cycle avoiding its poles,
CT defines an algebraic period
\begin{equation}\label{cas:eq-algebraicperiod}
 I_C^{\mathrm{alg}}(R\,dz)
 =2\pi i\sum_{a\ \text{a finite pole}}w_F(C,a)\,\res_a(R\,dz),
\end{equation}
where the winding number can be obtained from certified ray crossings or
appropriate real-algebraic index computations \cite{cas:SourceCT}. This is a
practical arbitrary-scale operation when the scalar ordered field is effective.

\begin{remark}[The normalized period may stay algebraic]\label{cas:rem-period}
$I_C^{\mathrm{alg}}/(2\pi i)$ may remain in an algebraic coefficient domain even
when multiplying by $2\pi i$ requires adjoining a named transcendental constant.
CT states agreement with Hahn integration on their common admissible contours.
\end{remark}

\subsection{Microscopic rescaling}
\label{cas:sub-rescaling}
CT rescales a radius-free germ by $z=a+t^sw$, with $a$ infinitesimal and $s>0$.
At each Hahn exponent the resulting coefficient in $w$ is an ordinary
polynomial, because only finitely many Taylor degrees contribute. \textbf{This
supplies a common ordinary domain \emph{after} rescaling where none existed
before} --- exactly the situation of the radius-free germ
\eqref{cas:eq-radiusfree}. The support certificate, rather than a guessed
radius, is what authorizes microscopic Cauchy and residue operations
\cite{cas:SourceCT}.

\subsection{The separating torus, and why the determinant matters}
\label{cas:sub-torus}
For the finite algebra of \cref{cas:sec-quotients}, form the coordinate
characteristic polynomials
\[
 Q_i(T)=\det(TI-M_{z_i})=T^\mu+\sum_{j<\mu}q_{ij}T^j .
\]
Using support-controlled normal forms, CT obtains a matrix $\mathsf C$ with
$Q_i(z_i)=\sum_jC_{ij}F_j$. Choose positive $s_i$ satisfying
\begin{equation}\label{cas:eq-radiusineq}
 (\mu-j)s_i<v(q_{ij})\quad\text{for each nonzero }q_{ij},
\end{equation}
and put $r_i=t^{s_i}$. CT's separating-torus formula is
\begin{equation}\label{cas:eq-torusformula}
 \Res_F(H)=\frac{1}{(2\pi i)^n}\HInt_{\mathbb T_r}
 \frac{H\det\mathsf C}{\prod_iQ_i(z_i)}\,dz_1\wedge\cdots\wedge dz_n ,
\end{equation}
the torus being parameterized by $z_i=r_ie^{i\theta_i}$ with ordinary angular
parameters. Formula \eqref{cas:eq-torusformula} is imported \emph{conditionally}
from CT's stated finite-algebra and transformation hypotheses
\cite{cas:SourceCT}.

\paragraph{Why $\det\mathsf C$ cannot be dropped.}
For the system \eqref{cas:eq-fourexample} take
\[
 \mathsf C=\begin{pmatrix}x^2+Ay&A^2\\B^2&y^2+Bx\end{pmatrix},
 \qquad
 \mathsf C\binom{F_1}{F_2}=\binom{x^4-A^2Bx}{y^4-AB^2y},
\]
so that
\begin{equation}\label{cas:eq-detC}
 \det\mathsf C=x^2y^2+Bx^3+Ay^3+ABxy-A^2B^2 .
\end{equation}
The \emph{separated} system has sixteen simple grid points, while the original
system has four. At a grid point that does not solve the original system,
$\mathsf CF=0$ with $F\ne0$, forcing $\det\mathsf C=0$, so the extraneous simple
residue vanishes. Dropping the determinant would compute the wrong invariant.
At the algebraic level, expanding the reciprocals,
\[
 Q_x(x)^{-1}=x^{-4}\sum_{k\ge0}(A^2B)^kx^{-3k},\qquad
 Q_y(y)^{-1}=y^{-4}\sum_{\ell\ge0}(AB^2)^\ell y^{-3\ell},
\]
which are strongly supported on an admissible scaled torus; extracting
$x^{-1}y^{-1}$ after multiplying by $H\det\mathsf C$ reproduces
\eqref{cas:eq-monomialresidue}. This is a finite coefficient calculation, not a
numerical integration over a fine-topological torus.

The pipeline is therefore concrete: compute multiplication matrices,
characteristic polynomials, an ideal-membership transformation matrix and the
valuation inequalities \eqref{cas:eq-radiusineq}; then reduce the integral to
coefficient extraction. The original denominators $F_i$ need not avoid the
chosen product torus, and a real $n$-torus is not the whole real
$(2n-1)$-dimensional boundary of a polydisk. These are mathematical conditions
that a symbolic simplifier can otherwise accidentally erase.

\subsection{Stokes as a certificate-based rewrite}
\label{cas:sub-stokes}
CT's Hahn differential forms have ordinary smooth coefficient forms with a
common well-ordered support. Exterior derivative, wedge product, admissible
pullback and integration along compact ordinary parameter chains satisfy the
coefficientwise Stokes identity, and positive-support deformations enlarge the
class of chains under stated analyticity conditions \cite{cas:SourceCT}.

A CAS can use this as a certificate-based rewrite: replace the integral of
$d\omega$ over a certified chain by the integral of $\omega$ over its oriented
boundary, or prove invariance under a certified pole-avoiding homotopy. This is
\textbf{a proposed use of an identified theorem, not an algorithm deciding
whether every arbitrary input map is admissible}. Nor does definable Jordan
separation alone furnish a general field-valued integration operation.

\subsection{Three circles that a single real precision parameter would confuse}
\label{cas:sub-threecircles}
In $\Gamma=\Q\omega+\Q$ take
\[
 R(z)=\frac1{z-t}+\frac1{z-u},\qquad u=t^\omega .
\]
A positively oriented microscopic circle $|z|=t^s$ has normalized period
\begin{equation}\label{cas:eq-threecircles}
 \frac1{2\pi i}\HInt_{|z|=t^s}R(z)\,dz
 =\begin{cases}2,&s=1/2,\\ 1,&s=2,\\ 0,&s=2\omega.\end{cases}
\end{equation}
A pole $t^a$ is inside exactly when $a>s$, and both residues equal one, so the
computation reduces to \emph{two exact exponent comparisons}. The example is
from CT and is independently checked at the valuation-vector level
(\cref{cas:sec-executed}). The rational-monomial core already supplies the
essential order operations, while a complete contour wrapper would still need
its chain and pole-avoidance certificates.

\subsection{Coherent primitives, fine-local constants, and a category warning}
\label{cas:sub-categorywarning}
Coherent primitives exist exactly when the ordinary closed-cycle Hahn periods
vanish. SA's fine-locally constant functions have zero derivative yet take
different values on different fine-open pieces.

\begin{remark}[The function category is part of the theorem's type]
\label{cas:rem-category}
A CAS must \textbf{not} combine SA's global fine-analytic logarithm
(\cref{cas:sub-logbranches}) with a coherent fundamental theorem to deduce that
the integral of $dz/z$ around a loop is zero. The global logarithm is a right
inverse of SA's exponentials and has derivative $1/z$, but is \emph{not} a
Hahn-coherent primitive on the punctured ordinary base. The function category is
part of the theorem's type, and the order-geometric route of
\cref{cas:sub-finitealgebra} does not create a fine-continuous contour
integration routine either.
\end{remark}

\section{Exponentials, logarithms, and transserial scale semantics}
\label{cas:sec-transcendental}

\subsection{Gonshor exp and log as exact constructors}
\label{cas:sub-gonshor}
The established surreal exponential $\ExpNo:\No\to\No_{>0}$ and its inverse
$\LogNo$ are canonical structures: an ordered-group isomorphism from $(\No,+)$
onto $(\No_{>0},\cdot)$ extending the ordinary real exponential and logarithm
\cite{cas:Gonshor,cas:vdDE}. Hence finite expressions such as
\[
 \ExpNo(\omega),\qquad \LogNo(\omega+1),\qquad \ExpNo(\omega+\omega^{-1})
\]
have exact meanings, and a finite node \texttt{ExpNo[x]} has an exact meaning
whenever $x$ does. The fact that a chosen implementation cannot yet expand one
of these expressions does not make its denotation inexact. For the last example
the exact group law gives
\begin{equation}\label{cas:eq-expgroup}
 \ExpNo(\omega+\omega^{-1})=\ExpNo(\omega)\sum_{n\ge0}\frac{\omega^{-n}}{n!},
\end{equation}
useful as a structured monomial-scale expression even without normalizing the
first factor in the smallest Hahn backend. \emph{Supporting the constructor is
easier than supporting a complete zero, comparison or normal-form algorithm for
every expression built from it.}

\subsection{The exp and log splits, and why the naive series fails}
\label{cas:sub-expsplit}
Writing $x=p+r+\eta$ with $p$ purely infinite, $r\in\R$ and $\eta$
infinitesimal,
\begin{equation}\label{cas:eq-expsplit}
 \ExpNo(x)=\ExpNo(p)\,e^r\,\Expi(\eta),\qquad
 \Expi(\eta)=\sum_{n\ge0}\frac{\eta^n}{n!}.
\end{equation}
The infinitesimal factor is directly computable by support-controlled expansion.
The purely infinite factor requires Gonshor-compatible monomial or transseries
machinery, or must remain an exact semantic node. \textbf{Defining it by
$\sum p^n/n!$ would be invalid}: for nonzero purely infinite $p$ the term
valuations decrease. Correspondingly, the compatibility identity
$\ExpNo(x+\epsilon)=\ExpNo(x)\sum_n\epsilon^n/n!$ for infinitesimal $\epsilon$
does \emph{not} say that the normal form of $\ExpNo(x)$ can be computed by the
ordinary exponential Taylor series at an arbitrary infinite $x$.

For a positive normal form $a=c\,t^\gamma(1+\eta)$ with $c>0$,
\begin{equation}\label{cas:eq-logsplit}
 \LogNo(a)=\log c+\LogNo(t^\gamma)+\Logi(1+\eta),\qquad
 \Logi(1+\eta)=\sum_{n\ge1}\frac{(-1)^{n+1}\eta^n}{n}.
\end{equation}
The middle term is \emph{not} a licence to conflate the Conway monomial map with
arbitrary exponential powers (\cref{cas:sub-conway}); a backend must implement
the applicable monomial-logarithm identities or retain that term as an atom.

For an effective implementation, the exponential should decompose the argument
into its large, constant and small parts, apply the small-part series where
valid, and introduce or recognize a new growth monomial for the large part.
Normalization, comparison and logarithmic inversion then require a transseries
or surreal-exponential backend, not just a new display head. A fixed Hahn
workspace containing the input may not be closed under these transcendental
operations; closure under algebraic roots in a divisible complex Hahn field is a
separate issue. Bournez and Guilmant study set-sized surreal fields stable under
exponential, logarithmic, derivative and antiderivative operations
\cite{cas:BG}; their structural closure conditions help choose mathematical
domains, but \emph{a bound on birthdays is not by itself a finite executable
representation or a complexity estimate}.

\subsection{An exact exp--log expression language}
\label{cas:sub-explog}
One practical language starts with a decidable algebraic core and permits
$\ExpNo$ and $\LogNo$ on certified positive arguments, plus positive-base real
powers. Its expression graphs are finite and denotationally exact; it can
simplify group identities, cancel inverse functions on their domains, extract
known leading scales, and perform infinitesimal expansions.

It should not claim that every equality between arbitrary such graphs has a
known terminating decision procedure. Broadening the ordinary constant domain
already complicates equality before any surreal exponent is introduced. A
partial simplifier can return a smaller equivalent graph together with a proof
of its rewrites, without producing a canonical normal form for the entire
language.

The backend should distinguish a \textbf{structural constructor} from a
\textbf{normalization request}. The former can often succeed immediately; the
latter may return a normalized object, a conditional result, a residual exact
node, or a resource-limit result. Conflating these operations either rejects
useful exact expressions unnecessarily or encourages unsupported
canonicalization claims. It should also keep positive-real powers
$\ExpNo(b\LogNo a)$ distinct from Conway monomial notation, and have a
documented rewrite theory with a declared transseries grammar when one is used.

\subsection{Lambert \texorpdfstring{$W$}{W} as an exact implicit constructor}
\label{cas:sub-lambert}
There are useful exact constructors beyond a fixed list of elementary functions.
A positive-real Lambert function can be specified by
\[
 W_+(y)\ge0,\qquad W_+(y)\ExpNo\bigl(W_+(y)\bigr)=y,\qquad y\ge0 .
\]
The ordinary real existence statement is first order in the ordered exponential
language, so the elementary-extension result supplies existence for nonnegative
surreal $y$ as well \cite{cas:vdDE}. Uniqueness follows directly: the map
$x\mapsto x\ExpNo(x)$ is strictly increasing for $x\ge0$, by positivity and the
strict increase of the real exponential.

Consequently a selected implicit-root node gives an exact value such as
$W_+(\omega)$ even before a normalizer has computed its logarithmic-scale
expansion; a transseries inversion algorithm can provide that expansion within
its supported class. \emph{The existence and uniqueness argument is a
certificate for this particular function, not a generic transcendental
root-finding algorithm}; general equations without comparable certificates
remain conditional or partial.

\subsection{Transseries backends and their published antecedents}
\label{cas:sub-transseries}
An exponential--logarithmic monomial hierarchy, small-monomial grids,
coefficient algorithms and support-preserving substitutions provide a more
systematic backend than ad hoc rewrite rules. It can expose formal
differentiation, integration within its supported differential class,
composition for admissible large positive arguments, and functional inversion
under the appropriate hypotheses, with extensive mathematical antecedents
\cite{cas:vdH}. The embedding into surreal numbers matters: it identifies the
formal large variable with $\omega$ and specifies the exponential structure. A
set of independent formal labels called $\log\omega$ and $e^\omega$ without
their structural relations is not the same ordered exponential field.

Three pieces should be recorded separately: the abstract transseries expression,
the chosen realization in a surreal field, and the permitted function-composition
or differential structure. Berarducci and Mantova construct surreal derivations
and study transseries as surreal function germs
\cite{cas:BMDerivations,cas:BMFunctions}. Mantova's 2026 preprint, in its May 9
revision, establishes monotonicity of omega-series composition in its second
argument and a Taylor approximation theorem, with an intermediate-value
consequence for the interpreted functions \cite{cas:Mantova2026}.

\begin{remark}[Scope of the transseries results]\label{cas:rem-transscope}
These are \textbf{semantic theorems for specified classes}, not universal
decision algorithms. They do not justify substituting arbitrary surreals into
every transseries expression without checking the chosen framework's
hypotheses, do not identify every arbitrary Hahn series with a globally
composable function on all of $\No$, and are not a claim that arbitrary CAS
expressions or arbitrary coefficient programs can be normalized. \emph{The
entire mathematical field of grid-based transseries is not a finite-input format
when arbitrary coefficient arrays are allowed.} A backend's scope should be
stated by grammar and algorithms, not by saying that it ``supports all
transseries''; terms outside the supported class can remain exact formal names
with reduced capabilities, provided their mathematical meaning has actually been
fixed. A reasonable first target is a finite-depth exponential--logarithmic
tower with explicit dominant-term certificates. Hyperseries
\cite{cas:Hyperseries} are a longer-term representational direction, not a
prerequisite for a first useful implementation.
\end{remark}

\subsection{Asymptotic data do not determine exact values}
\label{cas:sub-asymptotic}
The real functions $f(x)$ and $f(x)+e^{-x}$ have the same asymptotic expansion
in inverse powers of $x$ whenever the first expansion exists, because $e^{-x}$
is smaller than every such power as $x\to+\infty$. Under a selected compatible
surreal extension, however, the additional term at $x=\omega$ is a \emph{nonzero
infinitesimal}, namely $\ExpNo(-\omega)$, not zero. Knowing all algebraic orders
of an asymptotic expansion therefore does not by itself select one exact value.

An exact extension of a special function at infinity must state how
exponentially small data, sectors, branches and any summation choices are fixed.
Similarly the strong Hahn value $\sum n!\,t^n$ is \emph{not} automatically a
Borel sum, a lateral resummation, or the value of a particular analytic function
at $1/t$; such identifications require a summation and extension theorem and can
carry sector and Stokes data. An exact formal Hahn descriptor and an exact
analytic-function descriptor must remain distinct objects.

\subsection{The factorial series and its exponentially small homogeneous solution}
\label{cas:sub-factorial}
The exact formal-germ value $F(t)=\sum_{n\ge0}n!\,t^n$ is meaningful by Hahn
summability and satisfies the coefficient identity
\begin{equation}\label{cas:eq-factorialode}
 t^2D_tF+(t-1)F=-1 .
\end{equation}
Indeed, at degree $n\ge1$ the contribution is $(n-1)(n-1)!+(n-1)!-n!=0$, and the
constant term is $-1$. This provides an efficient recursive coefficient
descriptor. As an ordinary complex power series it diverges, which does not
prevent its strong Hahn value but does prevent identifying it, without further
choices, with an ordinary analytic germ at zero.

The associated homogeneous equation has the solution $e^{-1/t}/t$ on positive
ordinary $t$, as direct differentiation verifies, and it is smaller there than
every power of $t$. An asymptotic power-series solution therefore cannot detect
the addition of a multiple of this homogeneous term. In a Gonshor-compatible
infinite-scale interpretation at $t=\omega^{-1}$, the corresponding positive
term is $\omega\ExpNo(-\omega)$, \textbf{not zero}. Analytic solutions sharing
the same expansion may thus differ invisibly to every power of $t$.

Two further cautions attach to this example. It does not specify a unique
ordinary analytic function at positive real $t$ by convergence of that series.
And if a scalar derivation with $\partial t=-t^2$ is chosen, the displayed
equation is \emph{not} obtained by blindly replacing $D_t$ by $\partial$
(\cref{cas:sec-derivatives}). The example simultaneously tests summation
semantics and derivative semantics.

\subsection{All-scale rigidity as a category boundary}
\label{cas:sub-rigidity}
SA proves, under its explicit all-scale coherence definition, that functions
coherent entire at \emph{every surreal radius} are precisely the polynomials,
and that the analogous all-scale meromorphic functions are precisely the
rational ones. The quantifier ranges over all surreal scales and uses the
manuscript's uniform ordinary-coefficient chart condition
\cite{cas:SourceAnalysisTwo,cas:SourceAnalysisThree}.

\begin{remark}[What rigidity does and does not say]\label{cas:rem-rigidity}
The correct conclusion is a \textbf{separation of categories}, not a prohibition
on transcendental surreal functions and not a universal nonexistence claim. A
package that admits a global exponential should not also label it ``all-scale
Hahn-entire'' in that sense. A global special-function object should instead
advertise local germs, selected coherent charts, a transserial interpretation at
infinity, or a sectorial extension provider, with the compatible charts and
transition rules part of the object. No claim is made here that any one such
global atlas has already been implemented.
\end{remark}

\begin{remark}[Sharpness over a fixed value group]\label{cas:rem-rigiditysharp}
A companion report in this repository --- the surcomplex-foundations merged
report under \path{docs/surcomplex/}, whose analysis volume states the rigidity
theorem for class functions on all of $\No[i]$ --- records a sharpness witness
that should be read alongside the boundary above. Over the \emph{fixed} value
group of $K=\C((t^\Q))$, the function $f(z)=\sum_{n\ge0}t^{n^2}z^n$ is
nonpolynomial yet has a Hahn-coherent chart on the whole ordinary complex plane
at every positive radius in $\R((t^\Q))$. The rigidity proof cannot run there
because it requires a bound $\Delta$ with $|v(a_n)|<\Delta/4$ for all $n$, while
$v(a_n)=n^2$ is unbounded in $\Q$; the proper class of multiplicative scales
available for class functions on the full $\No[i]$ supplies an infinite
$\Delta$, and the fixed value group does not. This is a \textbf{sharpness result
about the scope of the theorem}, not a correction to it: it shows that the
all-scale rigidity statement is specific to the full-class scale system. A CAS
that works inside one fixed set-sized parent is in the second situation, not the
first, and must not invoke rigidity to conclude that its nonpolynomial coherent
objects cannot exist.
\end{remark}

\subsection{Oscillation is a separate layer}
\label{cas:sub-oscillation}
Oscillatory functions deserve their own layer. An expression such as $\sin x$ at
large positive real $x$ does \emph{not} have an eventual sign in the same sense
as a nonzero ordered logarithmic--exponential transseries. Treating oscillation
as just one more positive ordered monomial would lose its essential behaviour.
For infinite surreal arguments the extension convention of
\cref{cas:sub-infinitephases} must also be addressed.

\section{Surcomplex arithmetic, angles, branches, and infinite phases}
\label{cas:sec-surcomplex}

\subsection{Pairs and squared modulus}
\label{cas:sub-pairs}
Represent $z=x+iy$ as a pair of real-surreal objects with a common compatible
domain. Then
\begin{align*}
 (x+iy)(u+iv)&=(xu-yv)+i(xv+yu),\\
 \overline z&=x-iy,\qquad |z|^2=x^2+y^2,\\
 z^{-1}&=\frac{x-iy}{x^2+y^2}\qquad(z\ne0).
\end{align*}
If a real domain is an effective ordered field $F$, pairs implement $F[i]$ with
these formulas. The squared norm lies in the real base field; the norm
$|z|=\sqrt{x^2+y^2}$ may require its real closure, which is why a minimal
rational-monomial package implements \emph{squared} modulus but not unrestricted
modulus. Conjugation, real and imaginary parts, and squared modulus are finite
field operations; integer powers stay in the field, and rational powers
generally require algebraic extension and root selection.

No ordering should be imposed on the complexification as a field: a complex
field has no compatible order, and comparisons must specify real part, modulus,
valuation, or another intended relation. Real comparison gives \texttt{Min},
\texttt{Max} and piecewise evaluation when all required signs are decidable.

\subsection{All directions admit finite surreal angles}
\label{cas:sub-finiteangle}
TR proves that the unit circle
\[
 \mathbb T(\No)=\{u\in\SC:u\overline u=1\}
\]
is the image of canonical $\cis$ on \emph{finite} real surreals, with kernel
$2\pi\Z$. \textbf{Thus every nonzero surcomplex number has a modulus and a
finite angle class, even when its coordinates or its modulus are infinite}
\cite{cas:SourceTrig}.

For computation, let $u=z/|z|$ and $u_0=\st(u)$. Then
\begin{equation}\label{cas:eq-anglesplit}
 u=u_0\Expi(i\delta),\qquad
 \delta=-i\Logi(u/u_0)\in\mm\cap\No .
\end{equation}
The pair $(u_0,\delta)$ avoids any arbitrary angle cut: it is an exact,
\emph{cut-free} phase representation separating the ordinary unit direction from
the infinitesimal correction. Extracting $u_0$ still requires the scalar
backend's standard-part capability.

\subsection{The Cayley coordinate}
\label{cas:sub-cayley}
An entirely algebraic alternative, avoiding angle evaluation, is
\begin{equation}\label{cas:eq-cayley}
 u(q)=\frac{1+iq}{1-iq}=\frac{1-q^2}{1+q^2}+i\frac{2q}{1+q^2},
 \qquad q\in\No .
\end{equation}
It parametrizes the circle except the point $-1$, which must be checked
separately in solving and integration. \emph{Infinite $q$ is allowed and
describes the monad of the omitted direction.} The Cayley coordinate is used
twice below: as a cut-free exact circle parametrization here, and as a reduction
of circle positivity to a semialgebraic question over the real-closed core in
\cref{cas:sub-positivity}.

\subsection{Canonical strips for exponential and trigonometry}
\label{cas:sub-strips}
Write
\[
 \HH=\{x+iy:y\ \text{finite},\ x\in\No\},\qquad
 \mathscr V=\{x+iy:x\ \text{finite},\ y\in\No\}.
\]
TR's canonical strip exponential is
\begin{equation}\label{cas:eq-stripexp}
 E_{\HH}(x+iy)=\ExpNo(x)\cis(y),
\end{equation}
which is onto $\SC^\times$ with kernel $2\pi i\Z$. On $\mathscr V$, canonical
sine and cosine are
\begin{align}
 \Sin(x+iy)&=\sin x\cosh y+i\cos x\sinh y,\label{cas:eq-stripsin}\\
 \Cos(x+iy)&=\cos x\cosh y-i\sin x\sinh y,\label{cas:eq-stripcos}
\end{align}
with the real hyperbolic functions defined via $\ExpNo(\pm y)$. These formulas
admit exact symbolic nodes even at infinite $y$, conditional only on the scalar
backend's capabilities. In particular
\[
 \Sin(i\omega)=i\sinh\omega
\]
\textbf{is canonical in this framework}, whereas $\sin\omega$ and
$E(i\omega)$ are \textbf{not} fixed by its canonical finite-angle data. The
distinction is about the real circular phase, not merely about whether the
argument is large \cite{cas:SourceTrig}.

\subsection{Inverse trigonometry, and why an infinite argument needs no infinite phase}
\label{cas:sub-inversetrig}
The inverse functions have the natural real-surreal domains
$\arctan:\No\to(-\pi/2,\pi/2)$, $\arcsin:[-1,1]\to[-\pi/2,\pi/2]$ and
$\arccos:[-1,1]\to[0,\pi]$. For positive infinite $x$,
\[
 \arctan x=\frac\pi2-\arctan(1/x),
\]
so an exact finite-angle value is available without defining circular sine at
$x$ itself. Concretely,
\[
 \arctan\omega=\frac\pi2-\arctan t
 =\frac\pi2-t+\frac{t^3}{3}-\frac{t^5}{5}+\frac{t^7}{7}+\Ov{t^9},
\]
with $t=\omega^{-1}$: \emph{an infinite argument causes no infinite phase
choice}, unlike a request for $\sin\omega$. Endpoint branches must be explicitly
selected; \eqref{cas:eq-acos} is the endpoint example, requiring the positive
square-root uniformizer.

\subsection{Two logarithm branches must not share one cache key}
\label{cas:sub-logbranches}
An \emph{actual principal-angle} convention chooses a representative in
$(-\pi,\pi]$. SA instead also defines a \emph{standard-direction-anchored}
logarithm
\begin{equation}\label{cas:eq-anchoredlog}
 \mathcal L(z)=\LogNo|z|+i\Arg(u_0)+\Logi(u/u_0),
 \qquad u=z/|z|,\quad u_0=\st(u),
\end{equation}
where $\Arg(u_0)$ is an ordinary principal argument. The ordinary argument
convention is part of the definition. This function is fine-locally analytic on
the full punctured class plane with finite imaginary part, and is a right
inverse of SA's $E_\alpha$ (\cref{cas:sub-infinitephases}); it is \textbf{not} a
coherent primitive of $1/z$ around the ordinary puncture
\cite{cas:SourceAnalysisTwo,cas:SourceAnalysisThree}. See
\cref{cas:rem-category}.

The two choices differ at infinitesimal crossings of the ordinary branch cut.
For $\eta>0$ infinitesimal and $z=-1-i\eta$:
\[
 \text{actual principal angle}=-\pi+\arctan\eta,
 \qquad
 \text{anchored angle}=\pi+\arctan\eta ,
\]
a difference of exactly $2\pi$. Here $|z|=\sqrt{1+\eta^2}$ and
$\LogNo|z|=\tfrac12\Logi(1+\eta^2)$, so the real parts agree; the imaginary
parts differ by $2\pi$; and both exponential images under
\eqref{cas:eq-stripexp} equal $z$. The anchored value is permitted to lie
\emph{infinitesimally above} $\pi$ and \textbf{must not be silently reduced}
back into the principal interval while retaining a claim of the same global
fine-local analytic behaviour. A branch-keyed cache can store both without
inconsistency.

A function descriptor therefore needs a branch policy token such as
\texttt{ActualPrincipal} or \texttt{StandardDirectionAnchored}, and the chosen
policy belongs to exact expression identity, domain checks, differentiation and
serialization.

Two further logarithm-category warnings survive. The local identity
$\log(1+\epsilon)+\log(1+\eta)=\log((1+\epsilon)(1+\eta))$ is valid for the
infinitesimal logarithm at infinitesimal $\epsilon,\eta$; a global identity
$\log(zw)=\log z+\log w$ is \emph{not} justified by that local calculation, and
\texttt{PowerExpand}-style rewrites cannot be applied without recording the
relevant branch and domain assumptions. A safe simplifier may use
$\ExpNo(\LogNo x)=x$ for positive real surreal $x$, and a chosen surcomplex
right-inverse identity where the chosen constructions justify it; it must not
use an unrestricted $\log(\exp z)=z$ or cancel arbitrary complex powers without
branch hypotheses. Returning a set of roots is often preferable to imposing a
hidden principal-root policy on a nonstandard complex plane. A local ordinary
logarithm branch, a global fine-local logarithm, and the real positive logarithm
are different function types even when their printed name is \texttt{Log}.

\subsection{General complex powers and the strip obligation}
\label{cas:sub-powers}
For a selected logarithm $L$, the expression
\[
 z^w=E_{\HH}\bigl(wL(z)\bigr)
\]
is canonical by this formula \textbf{only when} $wL(z)\in\HH$: its imaginary
part need not be finite for arbitrary surcomplex $w$. Outside that strip one
must select a global exponential extension or retain an explicitly unresolved
power node.

Integer powers need none of this and should remain algebraic. Ordinary rational
powers can be implemented by algebraic root selection and finite-angle branch
rules rather than by introducing an unrestricted exponential. For positive real
base and real surreal exponent the real definition $\ExpNo(w\LogNo z)$ is
available. \emph{A single unqualified \texttt{Power} rewrite table cannot safely
cover all these cases.}

\subsection{Infinite phases I: the character classification (from TR)}
\label{cas:sub-characters}
Let $\Pi$ be the additive class of purely infinite real surreals; normal forms
give the additive splitting $\No=\Pi\oplus\OO_\R$, $x=p+\fin(x)$. TR classifies
extensions of the circular group law by characters
\[
 \chi:(\Pi,+)\longrightarrow(\mathbb T(\No),\cdot),
\]
setting
\begin{equation}\label{cas:eq-character}
 U_\chi(p+u)=\chi(p)\cis u,\qquad u\ \text{finite}.
\end{equation}
The real and imaginary parts give global cosine and sine satisfying the addition
laws and TR's fine-local differential identities, and \textbf{conversely every
extension with this group-law property has that form} --- a completeness claim
\emph{within TR's framework} \cite{cas:SourceTrig}.

The convention $\chi=1$ gives $U_0(x)=\cis(\fin(x))$, hence $\sin_0\omega=0$. It
is an \emph{available convention, not an identity forced by the finite-angle
theory}; other characters give different phases at $\omega$. An exact CAS should
default to an \texttt{ExtensionRequired} result for such a circular phase unless
the user selects a provider, and must never silently adopt the trivial
character.

A provider must respect additivity across its domain: independently assigning
phases to arbitrary related inputs can violate the group law. The phase at $2p$
must be the \emph{square} of the phase at $p$, and a phase assigned at $p/2$
must be a compatible square root. \emph{A character descriptor is structural
data, not a mutable table of unrelated evaluations.}

\begin{proposition}[Every global extension acquires an infinite period]
\label{cas:prop-periods}
Every extension \eqref{cas:eq-character} has an infinite period in each nonzero
purely infinite coset $p+\OO_\R$.
\end{proposition}
\begin{proof}
Finite-angle $\cis$ is already onto the entire surreal unit circle, so choose
finite $a$ with $\cis a=\chi(p)^{-1}$. Then $U_\chi(p+a)=1$. If $p\ne0$ the
element $p+a$ is infinite, and the homomorphism law makes it a period. This
reproduces the obstruction established in TR.
\end{proof}

Consequently a global phase provider \textbf{must not also advertise that its
only periods are ordinary integer multiples of $2\pi$}; those two requirements
are incompatible in this framework. Exact results computed with different
providers must retain their provider identifiers, even when they happen to agree
on all ordinary test values.

\subsection{Infinite phases II: the \texorpdfstring{$E_\alpha$}{E-alpha} family (from SA)}
\label{cas:sub-infinitephases}
A separate nonuniqueness phenomenon, for a \emph{different} family from a
\emph{different} supplied manuscript, is SA's explicit one-real-parameter family
of global exponentials. For a real surreal $y$ let $\fin(y)$ retain its
nonnegative $t$-exponent terms and let $\ell(y)$ be its coefficient at
$t^{-1}=\omega$. For an ordinary real $\alpha$ put
$\theta_\alpha(y)=\fin(y)+\alpha\ell(y)$ and
\begin{equation}\label{cas:eq-phaseexp}
 E_\alpha(x+iy)=\ExpNo(x)\bigl(\cos\sh\theta_\alpha(y)
 +i\sin\sh\theta_\alpha(y)\bigr)
 =\ExpNo(x)\exp_{\mathrm{fin}}\bigl(i\theta_\alpha(y)\bigr).
\end{equation}
According to SA these are all global fine-analytic exponential group
homomorphisms extending the ordinary complex exponential and satisfying
$E_\alpha'=E_\alpha$. Nevertheless
\begin{equation}\label{cas:eq-phasecontrast}
 E_0(i\omega)=1,\qquad E_\pi(i\omega)=-1 .
\end{equation}
Agreement on ordinary complex numbers, on the group law, and on fine-local
analyticity therefore do \emph{not} select a unique value at $i\omega$
\cite{cas:SourceAnalysisTwo,cas:SourceAnalysisThree}.

\begin{remark}[Both accounts stand; neither may be extended into the other]
\label{cas:rem-phaseseparate}
These are two separate statements and this report does not merge them.
\begin{itemize}
\item \cref{cas:sub-characters} is a \emph{completeness statement about circular
group-law extensions within TR's framework}. It must not be extended to SA's
exponential family: that would attribute a completeness result to a source that
did not supply it.
\item \cref{cas:eq-phaseexp} is an \emph{explicit family from SA demonstrating
nonuniqueness}, used as a manuscript result and an interface warning. Both
articles carrying it state their non-claims explicitly: one says the statement
``does not assert that no other published convention can be chosen'' --- it says
only that SA's stated local axioms do not make that choice --- and the other
says it uses the contrast ``as a manuscript result and an interface warning, not
as a published uniqueness classification''. \textbf{We therefore do not write
``the classification of global surcomplex exponentials is the $E_\alpha$
family''.}
\end{itemize}
\end{remark}

\begin{remark}[Effectivity caveat, given independently by both carriers]
\label{cas:rem-phaseeffectivity}
Formula \eqref{cas:eq-phaseexp} is not automatically an effective algorithm on
every exact Hahn generator. Extracting $\fin(y)$ and $\ell(y)$ is a
\emph{support-cut operation}: it is mathematically well defined on normal forms,
but need not be computable by scanning an arbitrary coefficient stream. On
finite sparse data and suitably certified structured data it is
straightforward; on an unrestricted program representation its operational
availability must be declared. The representation problem therefore reappears
even after the semantic choice has been settled.
\end{remark}

\begin{remark}[The cache-key rule, three ways]\label{cas:rem-cachekey}
All three sources state a version of the same rule, and all three survive.
\begin{enumerate}
\item The different $E_\alpha$ extensions \textbf{cannot share a cache keyed
only on the printed expression} \texttt{Exp[z]}.
\item \textbf{A coefficient computed under one phase convention cannot be
reused under another}; caches should be keyed by normalized parent, exponent,
branch, derivation and extension policy, and memoized normal forms should
distinguish an exact result from a precision-limited result and never silently
upgrade a guessed leading term to a certificate.
\item Branch, phase, support and derivation identifiers belong in
\emph{persistent} cache keys: reusing a result under a changed interpretation is
\textbf{an exactness bug even when all printed numbers look identical}.
\end{enumerate}
A proposed \texttt{SurcomplexExp[z]} must therefore restrict its input domain,
specify a particular extension policy, or remain unassigned at such arguments,
and must not acquire an unqualified value merely by extending ordinary Taylor
syntax. The same warning applies to unrestricted trigonometric functions and to
expressions such as $E_\alpha(1/z)$ near a microscopic essential singularity.
\end{remark}

\subsection{Finite trigonometric polynomials and the \texorpdfstring{$2n$}{2n} root bound}
\label{cas:sub-trigpoly}
Let $n$ be an ordinary finite integer and consider, on $\mathscr V$,
\[
 T(z)=\sum_{k=-n}^{n}c_kE_{\HH}(ikz),\qquad c_k\in\SC .
\]
Setting $U=E_{\HH}(iz)$ gives
\begin{equation}\label{cas:eq-trigalgebra}
 T(z)=U^{-n}P(U),\qquad P(U)=\sum_{k=-n}^{n}c_kU^{k+n}.
\end{equation}
TR proves that the nonzero roots of $P$ correspond, with multiplicity, to the
zeros of $T$ modulo $2\pi$ in the canonical strip; their number is at most $2n$,
and exactly $2n$ when both extreme coefficients are nonzero
\cite{cas:SourceTrig}.

This gives an implementable exact solver: build $P$ with the scalar polynomial
backend, solve it in the appropriate complex algebraic extension, discard
$U=0$, retain multiplicities, and convert a nonzero $U$ to a strip angle class
through a logarithm. That final analytic reconstruction may leave the algebraic
coefficient core even when solving for $U$ stays inside it. To select finite
real angles, impose $U\overline U=1$; alternatively use the Cayley
parametrization \eqref{cas:eq-cayley} with a real algebraic solver, remembering
the omitted point $-1$.

\begin{remark}[An infinite frequency is not a polynomial degree]
\label{cas:rem-infinitefreq}
An ``infinite integer frequency'' is not an ordinary polynomial degree.
\textbf{Replacing $n$ by $\omega$ in this algorithm is meaningless.}
Infinite-frequency oscillation requires a different function class and, in
general, an infinite-phase provider.
\end{remark}

\subsection{Positivity must see infinitesimal angles}
\label{cas:sub-positivity}
For positive infinitesimal $\varepsilon$, TR's example
\begin{equation}\label{cas:eq-positivity}
 T(\theta)=(\cos\theta-\varepsilon)^2-\varepsilon^4
\end{equation}
is positive at \emph{every ordinary real angle} but is \emph{negative} at the
finite surreal angle $\theta=\arccos\varepsilon$. A positivity test using only
ordinary angular samples or standard parts is therefore \textbf{unsound}
\cite{cas:SourceTrig}.

Two remedies are available for coefficients in the real-closed algebraic core.
One reduces circle positivity to a semialgebraic question using the Cayley
coordinate \eqref{cas:eq-cayley}, inspecting the omitted circle point
separately. The other uses TR's Fej\'er--Riesz theorem for a finite algebraic
factorization route. Both test the actual ordered field, not only its ordinary
shadow.

\section{Differentiation: three operators}
\label{cas:sec-derivatives}

\subsection{The three operators and their worked contrasts}
\label{cas:sub-threederivatives}
A system must distinguish differentiation of a function in a variable from
differentiation of its scalar coefficients.

\paragraph{The coordinate derivative $D_z$.}
For a polynomial or coherent function in a variable $z$,
\[
 F(z)=\sum_\gamma f_\gamma(z)t^\gamma,\qquad
 D_zF=\sum_\gamma f_\gamma'(z)t^\gamma ,
\]
so that every $t^\gamma$ is a \emph{coefficient} and $D_z(t^\gamma)=0$; in
particular $D_z\omega=0$ and $D_z(\omega^{-1})=0$. This is the derivative used
in the supplied analytic manuscripts; it commutes with coefficientwise Hahn
assembly while holding monomials constant, and it is the appropriate operation
for local holomorphic identities, Jacobians, residues and CT's contour calculus.
Concretely,
\[
 D_z(\omega z^2+t\sin z)=2\omega z+t\cos z .
\]

\paragraph{The formal parameter derivative $D_t$ (or $D_T$).}
For a formal power series in an auxiliary indeterminate $T$, the formal
derivative satisfies $D_TT=1$. In a rational or Puiseux field with $t$ an
independent parameter and constant rational exponent $q$, $d(t^q)/dt=qt^{q-1}$;
this changes valuations and may decrease precision. In a multivariate
rational-monomial field, each formal partial derivative treats the other
independent variables as constants. Evaluating a formal series at a fixed
surreal $t$ produces a \emph{number}, not an ordinary variable on which $D_z$
should suddenly act.

\paragraph{A surreal-field derivation $\partial$.}
A surreal derivation is additional global structure on $\No$; the
Berarducci--Mantova construction supplies a distinguished differential-field
framework compatible with the exponential structure
\cite{cas:BMDerivations}. Normalizing by $\partial\omega=1$ and putting
$t=\omega^{-1}$ gives
\[
 \partial t=-t^2,\qquad \partial(\omega^{-1})=-\omega^{-2},
 \qquad \partial(t^q)=-q\,t^{q+1}\ \ (q\in\Q).
\]
Contrast this with $D_t(t^3)=3t^2$ and $\partial(t^3)=-3t^4$, and with
$D_z(\omega^{-1})=0$.

These are \textbf{not contradictory outputs}: they are different operators.
Extending a scalar derivation through a Hahn sum requires its own
strong-additivity hypotheses. For general monomials with nonconstant surreal
exponents one must use the chosen derivation's rules, not an unqualified power
rule copied from ordinary integer exponents: the rational-power identity
$\partial(t^q)=-q\,t^{q+1}$ must not be extended to every surreal exponent of
the omega-map by treating that exponent as an ordinary constant.

\subsection{Three proposed name sets, printed as three proposals}
\label{cas:sub-derivnames}
All three sources propose API names for these operators, and none of the three
is implemented in any delivered file. We print all three as proposals rather
than choosing one.
\begingroup\small
\begin{longtable}{@{}P{0.19\textwidth}P{0.28\textwidth}P{0.24\textwidth}P{0.19\textwidth}@{}}
\toprule
Operator & Capability-tiers proposal & Layered-representations proposal &
Gaussian-prototype proposal\\
\midrule\endhead
Coordinate derivative & \texttt{FunctionDerivative} & \texttt{CoordinateD} &
a separate operator and parent from the surreal derivation\\[0.4em]
Formal series derivative & \texttt{FormalSeriesDerivative} & \texttt{FormalD} &
formal parameter derivative in the declared parent\\[0.4em]
Surreal derivation & \texttt{ScalarDerivation}, with a mandatory derivation
identifier & \texttt{SurrealDerivation[..., derivationID]} & invoked through a
separate operator and parent\\
\bottomrule
\end{longtable}
\endgroup
All three sources agree that converting between the interpretations must be a
deliberate operation, that a derivation identifier is mandatory for the third
operator, and that automatic use of a native \texttt{D} is reasonable only after
the chosen coordinate semantics have been established. \emph{A default
\texttt{D} overload cannot infer all these intentions from a displayed
monomial}, and applying a constant-exponent rule when the exponent is itself
nonconstant for that derivation can simply be wrong. A package should record
whether a symbol is a coordinate, a formal parameter, or a designated surreal
constant.

\subsection{Differential equations must name their category}
\label{cas:sub-decategory}
A fine-local solution of $f'=f$ with a prescribed value at zero need \emph{not}
be globally unique in the entire fine-local category, as the $E_\alpha$ family
of \cref{cas:sub-infinitephases} shows. A local formal initial-value problem can
still be uniquely solved under the usual coefficient-recursion hypotheses.
\textbf{A solver must state the category in which its uniqueness claim is
made.}

\section{Host integration: what Wolfram Language supplies and what an adapter must add}
\label{cas:sec-host}

\subsection{A documentation-based assessment}
\label{cas:sub-docassessment}
The following assessment uses official documentation inspected for the source
articles on September 21, 2026. It identifies reusable facilities. It is
\textbf{not} a claim that every function automatically accepts an arbitrary
user-defined field, and \textbf{not} an exhaustive survey of third-party
packages. \emph{No documented native object with all of the arbitrary-Hahn and
surreal contracts developed here was identified in the inspected material.}

\begingroup\small
\begin{longtable}{@{}P{0.24\textwidth}P{0.33\textwidth}P{0.35\textwidth}@{}}
\toprule
Facility & What it contributes & What still needs a surreal adapter\\
\midrule\endhead
\texttt{Root}, \texttt{AlgebraicNumber}, \texttt{RootReduce} & Exact ordinary
algebraic coefficients and implicit-root precedents. & Ordered monomial domains,
non-Archimedean selectors, custom coefficient semantics; do not borrow ordinary
complex root ordering blindly.\\[0.5em]
\texttt{Together}, \texttt{Cancel}, polynomial routines & Exact fraction and
finite polynomial arithmetic --- the engine all three prototypes actually use. &
Declare the independent monomials and their ordered exponent parent before
interpreting the result.\\[0.5em]
\texttt{SeriesData}, \texttt{SeriesCoefficient} & Finite power and Puiseux
series views, coefficient extraction, order tracking. & Arbitrary ordered
exponent groups, infinite-prefix handling, Hahn support certificates, exact-tail
semantics.\\[0.5em]
\texttt{Series}, \texttt{InverseSeries}, composition & Ordinary local expansions
and reversion, including supported fractional or logarithmic forms. & A proof
that the requested surreal evaluation is admissible and that the chosen branch
is the intended one.\\[0.5em]
\texttt{Asymptotic}, \texttt{AsymptoticSolve} & Ordinary asymptotic
approximations and local equation expansions. & Exact Hahn denotation,
beyond-all-orders data, and selected surreal extension semantics.\\[0.5em]
\texttt{DifferentialRoot} & A finite ordinary-function representation by linear
differential equations plus initial conditions. & Strong evaluation, support and
singular-point certificates, coefficient-domain closure.\\[0.5em]
\texttt{D}, polynomial and matrix routines & Reusable finite algebra, ordinary
coefficient derivatives, symbolic coefficient manipulation. & Exact zero tests;
identify which variables and monomials are constant for the chosen derivation;
not a generic automatic custom-field interface.\\[0.5em]
\texttt{UpValues}, custom heads, \texttt{Inactive} & Extension mechanisms,
controlled symbolic syntax, preserved unevaluated descriptions. & Mathematical
validity, termination policy, branch conventions, deliberate numerical
behaviour; avoid broad rewrites and implicit parent mixing.\\[0.5em]
\texttt{NumericQ}, \texttt{N} & Ordinary numerical evaluation protocol. & A
separate scale-aware numerical interface; see
\cref{cas:sub-numericq}.\\[0.5em]
\texttt{PossibleZeroQ} & A documented \emph{heuristic} facility. & Nothing:
it must not be used as a proof of vanishing or nonvanishing at all.\\[0.5em]
\texttt{Infinity} & Extended-limit arithmetic. & Nothing: it must never be
redefined to mean $\omega$.\\
\bottomrule
\end{longtable}
\endgroup
The representation details are documented in
\cite{cas:WLRoot,cas:WLAlgebraic,cas:WLRootReduce,cas:WLSeriesData,%
cas:WLSeriesCoefficient,cas:WLSeries,cas:WLInverseSeries,cas:WLAsymptotic,%
cas:WLAsymptoticSolve,cas:WLDifferentialRoot,cas:WLUpValues,cas:WLInactive,%
cas:WLNumericQ,cas:WLPossibleZeroQ,cas:WLInfinity,cas:WLSpecial}. The proposed
adapters in the last column are not native-feature claims.

\subsection{Do not understate the current scope of \texttt{Root}}
\label{cas:sub-rootscope}
Two of the three sources give the same self-correction, in nearly identical
words, and it is worth stating plainly because the opposite error is the
tempting one. Current documentation distinguishes indexed polynomial-root
representations from neighbourhood-based representations, and includes exact
roots of certain transcendental equations and systems \cite{cas:WLRoot}.
\textbf{The assertion ``\texttt{Root} can only represent algebraic numbers over
$\Q$'' would be incorrect.} Its ordinary algebraic subset remains particularly
useful as an exact coefficient backend.

That said, none of those documented ordinary numerical neighbourhood mechanisms
by itself chooses between roots whose standard parts are identical but whose
surreal displacements differ, and an ordinary symbolic parameter in
\texttt{Root} does not automatically acquire the ordering of an infinitesimal. A
new \texttt{SurrealAlgebraicRoot} can reuse the same general idea of an implicit
exact object while supplying its own ordered-domain and branch certificates;
native ordinary \texttt{Root} objects can still be used for residual algebraic
constants.

\subsection{\texttt{SeriesData} is a view, not the universal storage format}
\label{cas:sub-seriesdata}
A basic \texttt{SeriesData} object stores a finite coefficient list, integer
bounds, and a common denominator for the exponent progression, together with
omitted-order information, and it supports arithmetic and calculus operations
\cite{cas:WLSeriesData}. It is well suited to a finite rational-power jet and to
a one-parameter adapter. It is \emph{not} a representation of every well-ordered
surreal exponent support; the fact that \texttt{Series} can also produce certain
logarithmic and other expansions does not change that limitation of the basic
exponent-lattice format \cite{cas:WLSeries}. Nested \texttt{SeriesData} objects
can be useful when a fixed iterated expansion convention is supplied; their
existence does not automatically define a general multi-rank Hahn truncation
policy.

\textbf{Converting with \texttt{Normal} discards the omitted-order term}, and it
must not be used internally as a proof that an infinite exact series equals its
displayed polynomial. Round-tripping a jet through plain syntax should retain
its precision certificate in a wrapper or explicitly announce its loss. A
regression suite should check that \texttt{Normal} or a similar jet-to-expression
conversion cannot silently promote a truncation to an exact equality.

\subsection{\texttt{PossibleZeroQ} is heuristic}
\label{cas:sub-possiblezeroq}
\texttt{PossibleZeroQ} is documented as a heuristic facility
\cite{cas:WLPossibleZeroQ}. It must \textbf{never} serve as the proof of a
leading coefficient's nonvanishing or vanishing. Testing a coefficient with an
ordinary numerical heuristic and proceeding as though the result were an exact
algebraic fact is unsound (\cref{cas:sub-finitealgebra}).

\subsection{\texttt{NumericQ}, \texttt{N}, and the Archimedean impossibility}
\label{cas:sub-numericq}
\begin{proposition}[No order-preserving embedding]\label{cas:prop-archimedean}
There is no order-preserving embedding of an ordered field containing a positive
infinitesimal into $\R$ that fixes the rationals.
\end{proposition}
\begin{proof}
If $0<\epsilon<1/n$ for every positive integer $n$, its image would be a
positive real below every $1/n$, contradicting the Archimedean property of $\R$.
Arbitrary precision does not change the argument.
\end{proof}

Three distinct numerical requests are nevertheless useful and should be separate
operations.
\begin{enumerate}
\item \textbf{Coefficient approximation}: approximate the ordinary coefficients
while retaining exact monomials.
\item \textbf{Standard part}: extract the ordinary constant coefficient of a
finite value, discarding infinitesimals by definition, and only with a
certificate of finiteness.
\item \textbf{Scale specialization}: replace finitely many monomials by selected
real parameters for a numerical experiment.
\end{enumerate}

\begin{remark}[A finite-scale specialization is a simulation]
\label{cas:rem-simulation}
The third request is \emph{not} a faithful model of all surreal inequalities: no
fixed positive real $u$ stays below every positive power of a fixed $0<t<1$.
A specialization such as $t_1=10^{-6}$ can illustrate a one-scale expression,
but it is not its surreal value. For several scales, numerical choices can
respect any prescribed \emph{finite} collection of hierarchy inequalities, not
the entire statement $t_2<t_1^{\,n}$ for all ordinary $n$. Results from such a
specialization must be \textbf{labelled as simulations}. They cannot prove exact
non-Archimedean order and cannot resolve an expression whose leading coefficient
is symbolically unknown. A finite algebraic identity survives any defined field
homomorphism, but an arbitrary infinite Hahn expansion need not survive
replacement of its scale by a small positive real number --- the nondiscrete-support
and factorial-series examples give two independent reasons.
\end{remark}

Accordingly, \texttt{NumericQ} is not a declaration of membership in a
non-Archimedean ordered field; its documented role concerns numeric expressions
\cite{cas:WLNumericQ}. Marking an infinite surreal as an ordinary numeric value
can send unrelated numerical algorithms into an impossible real approximation,
or make numerical solvers demand a conventional machine or arbitrary-precision
number the object cannot supply. Use a custom exact-domain predicate and supply
numerical views only through explicitly described coefficient approximation or
finite-scale specialization. A generic \texttt{N[value]} rule should not pretend
to return an ordinary complex approximation to an arbitrary surreal number, and
converting all infinities to native \texttt{Infinity} and all infinitesimals to
zero is neither an injective representation nor a field homomorphism. A written
\texttt{N[surreal,p]} requires a declared meaning: a certified coefficient
approximation, a valuation jet, or a scale-aware structured result.

\subsection{Infinity}
\label{cas:sub-infinity}
\textbf{Never redefine \texttt{System`Infinity} to mean $\omega$.} The latter is
a field element with $\omega-\omega=0$; the former has extended-limit semantics
\cite{cas:WLInfinity}. In all three delivered prototypes \texttt{Infinity}
appears only as the distinguished valuation marker $v(0)=+\infty$.

\subsection{Conservative dispatch, explicit operations before upvalues}
\label{cas:sub-dispatch}
Inert symbolic heads and controlled evaluation keep the host from simplifying
away a domain distinction before the extension sees it, and \texttt{Inactive}
and upvalue mechanisms are documented tools for such an integration
\cite{cas:WLInactive,cas:WLUpValues}. But the first reference implementation
should expose \emph{explicit} operations --- all three delivered prototypes do
exactly this --- because their domains are easier to test than an early
collection of global automatic rewrites. Once the core is stable, narrowly
targeted upvalues can provide natural notation for operations whose semantics
are unconditional in the operands' domains.

\textbf{Do not globally unprotect \texttt{Plus}, \texttt{Times},
\texttt{Power}, \texttt{Exp}, \texttt{Sin}, \texttt{Log} or \texttt{Less} and
redefine them for arbitrary expressions.} Branch-dependent transformations
should remain explicit or return their conditions. The package must avoid
evaluation loops, ambiguous mixed-domain coercions and branch-changing rewrites.
Upvalues and \texttt{Inactive} do not create mathematical correctness on their
own, and they do not remove the responsibility to check mixed parents and
analytic hypotheses. Keep raw monomial symbols internal or require them to be
unassigned; promoting two parents requires an explicit embedding preserving
coefficient arithmetic, exponent relations, order and function policy.

Results should distinguish invalid input, an undefined mathematical operation,
an unsupported algorithm, and an unknown proposition. \textbf{A failed zero test
must not return \texttt{False}}; a negative power at zero is genuinely
undefined; a square root outside a rational parent asks for an extension rather
than proving impossibility. On general symbolic inputs, a conditional result
should preserve its conditions instead of assuming them to simplify the output.

\subsection{Cache keys}
\label{cas:sub-cachekeys}
Cache keys must include the coefficient assumptions, the exponent order, the
embedding, the branch convention, the derivation interpretation and the
extension policy. A normal form computed under one convention must not be reused
after any of those change. See \cref{cas:rem-cachekey} for the three forms in
which the sources state this rule.

\subsection{Serialization and the code-execution risk}
\label{cas:sub-serialization}
Serialization should preserve semantic constructors and parent identifiers
rather than pretty output, and should reference shared expression nodes rather
than copy exponentially growing trees. It should also preserve whether a
displayed term is exact or carries a remainder.

\begin{remark}[Security]\label{cas:rem-security}
\textbf{Coefficient programs and function definitions may contain executable
Wolfram code, so loading untrusted serialized programs is a code-execution
risk.} A production format should use a validated expression grammar and
explicitly trusted extension modules rather than automatically evaluating
arbitrary input strings. Importing an arbitrary coefficient program must not
execute it automatically: mathematical admissibility and safe program execution
are separate checks. Explicit resource budgets and opt-in evaluation of external
code are suitable engineering boundaries. Relatedly, none of the three delivered
prototypes is a hardened parser for arbitrary or untrusted input, and in all
three the internal data heads are private: \emph{manually forging them bypasses
validation and lies outside each package's own contract}.
\end{remark}

\subsection{Lossless coercion before convenient notation}
\label{cas:sub-coercion}
Embedding a smaller ordered monomial field into a larger one is legitimate when
the map is explicit and order preserving, and adjoining algebraic roots is
legitimate. Replacing a named transcendental constant by an independent symbol,
replacing an exact value by a jet, changing a branch, or discarding an
infinitesimal by taking a standard part is \textbf{not} a lossless coercion. The
public interface should make those operations visibly intentional: taking a
standard part requires a certificate of finiteness, and scale specialization
must distinguish ordinary parameter specialization from evaluation at an actual
surreal monomial. Promoting a rational node to a complex pair or an algebraic
extension must preserve the embedding rather than simply discard the older
domain tag.

\subsection{Precedents in other systems}
\label{cas:sub-sage}
The same distinctions are visible in other CAS designs and are useful
architectural precedents. Sage's lazy-series documentation explicitly
distinguishes coefficients computed on demand from equality, which is decidable
only in special cases \cite{cas:SageLazy}. Its asymptotic rings separate
coefficient rings, growth groups, exact terms and order terms
\cite{cas:SageAsymptotic}, working with growth and error terms rather than
purporting to encode every surreal. Existing domain definitions should be reused
where they fit, not relabelled wholesale as the full surreal class. These
systems are sources of algorithms and interfaces, not substitutes for the
semantic choices made here. A CAS other than Wolfram Language can follow the
same architecture: the mathematical requirements are an exact coefficient
engine, programmable expression types, an ordered exponent representation, and
validated adapters for its polynomial and series facilities.

\section{A proposed architecture --- explicitly not delivered}
\label{cas:sec-architecture}

\begin{quote}\small\itshape\raggedright
Every statement in this section is a proposal. Nothing named
\texttt{SurrealDomain}, \texttt{SurrealRational}, \texttt{SurrealAlgebraic},
\texttt{CertifiedHahn}, \texttt{SurrealExpression}, \texttt{SurrealJet},
\texttt{Surcomplex}, \texttt{AnalyticGerm}, \texttt{CoherentFunction},
\texttt{ExactSurreal}, \texttt{SurrealEqual}, \texttt{SurrealExpand},
\texttt{SurrealCoefficient}, \texttt{SurrealRoot}, \texttt{CoherentLift},
\texttt{HahnParent}, \texttt{HahnValue}, \texttt{CoherentFamily} or
\texttt{SurcomplexFunction} exists in any delivered file described in
\cref{cas:sec-prototypes}. No capability named in this section may be reported
as implemented.
\end{quote}

\subsection{Architecture requirements all three sources agree on}
\label{cas:sub-requirements}
Stripped of syntax, the three proposals agree on what a domain, an exact value,
a view, a function and a result must record.

\begin{itemize}
\item A \textbf{domain} (or parent) should record: the coefficient domain; the
ordered exponent domain, including its order and coordinate significance; the
embedding into the surreal numbers; the support presentation or support
grammar; any exponential, branch, extension and derivation conventions; and the
list of available algorithms.
\item An \textbf{exact value} is a finite expression or certified construction
\emph{in} a named domain, together with its certificates.
\item A \textbf{view} is a computed expansion, algebraic root approximation, or
enclosure of an exact value. \emph{Computing another view must not silently
change the value or its domain.} A rational expression, an algebraic root
description and a Hahn expansion can denote the same element while serving
different purposes; keeping all three as related representations is more useful
than forcing every result into one preferred infinite normal form.
\item A \textbf{function} should record its variables, coefficient parent,
ordinary or nonstandard domain, regularity category, branch, and extension
policy.
\item \textbf{Capabilities should be queryable, not guessed from a head.} A
domain should advertise separately whether equality and order are decidable,
whether coefficients and leading terms are computable, and whether requested
cutoff expansions are finite. Crucially, \emph{some capabilities are local to a
value or a request}: a particular series can have a certified leading term even
when its representation class has no general leading-term algorithm, and a
cutoff may admit a finite expansion although a larger cutoff does not.
\item \textbf{Certificates should be checkable more cheaply than the search that
found them.} Useful certificates include an exact coefficient normal form; a
nonzero leading term with a lower-support exclusion; an algebraic root selector;
a positive-grid decomposition bound; a nonvanishing ordinary denominator on a
specified domain; a polynomial residual with a valuation lower bound; a finite
support; a positive-support monoid; an exact polynomial identity; a host
algebraic-number sign certificate; or a branch identity under recorded
inequalities. The checker need only verify the relevant finite statement.
\item \textbf{Unproved is not false, and a resource limit is not a
counterexample.} A missing zero proof is not a mathematical domain error, and an
input outside a chosen sine strip is not a proof that no extension exists. The
constructor should state whether its evidence has been \emph{checked},
\emph{supplied as an assumption}, or \emph{inherited from a proved
construction}. \emph{User assertions must not be relabelled as kernel-verified
theorems.}
\end{itemize}

\subsection{Three concrete syntax sketches, side by side}
\label{cas:sub-sketches}
The three proposals are printed as three alternatives, each labelled with its
origin. No fourth invented syntax is synthesized here, and none of the three is
silently adopted.

\paragraph{Sketch A --- capability-tiers proposal (positional).}
\begin{lstlisting}
SurrealDomain[coefficientDomain, exponentDomain, interpretation]
SurrealRational[domain, numerator, denominator]
SurrealAlgebraic[domain, polynomial, rootSelector]
CertifiedHahn[domain, supportDescriptor, coefficientDescriptor]
SurrealExpression[domain, expressionGraph]
SurrealJet[domain, center, remainderDescriptor]
Surcomplex[realPart, imaginaryPart]
AnalyticGerm[functionDescriptor, center, branch, domain]
CoherentFunction[ordinaryDomain, support, coefficients]
(* explicit operations: SRAdd, SRSin *)
\end{lstlisting}

\paragraph{Sketch B --- layered-representations proposal (option association).}
\begin{lstlisting}
dom = SurrealDomain[
  "Coefficients" -> "RealAlgebraic",
  "Exponents"    -> LexicographicRationals[2],
  "Embedding"    -> ConwayMonomials[{Omega, 1}]];
x = ExactSurreal[representation, dom];
SurrealEqual[x, y, "Evidence" -> True]
SurrealExpand[x, "ValuationCutoff" -> {1, 0}]
SurrealCoefficient[x, {0, 17}]
SurrealRoot[p, "Selector" -> certifiedSelector]
CoherentLift[f, "Domain" -> U, "Branch" -> branchData]
(* explicit operations: SAdd, SMultiply, SInverse, SCompare *)
\end{lstlisting}
The embedding record in Sketch B means that exponent $(a,b)$ denotes
$a\omega+b$. It is not an instruction to identify an unspecified host symbol
with the entire surreal class, and an implementation should store such records
in a validated internal format rather than rely on their printed appearance.

\paragraph{Sketch C --- Gaussian-prototype proposal (parent/value split).}
\begin{lstlisting}
HahnParent[coefficientDomain, exponentDomain,
  "MonomialInterpretation" -> embedding,
  "SupportClass" -> supportGrammar]
HahnValue[parent, representation, certificates]
AnalyticGerm[center, definition, initialData, domain, branch]
CoherentFamily[parent, ordinaryDomain, coefficientProgram]
SurcomplexFunction[definition, domain, extensionPolicy]
(* explicit operations: HAdd, HMul, HCompare *)
\end{lstlisting}

\begin{remark}[Two head collisions among the proposals]\label{cas:rem-headcollisions}
The three sketches are \emph{not} compatible, and two of the collisions are
between heads with the same name and different argument shapes.
\begin{itemize}
\item \texttt{SurrealDomain} appears in Sketch A as a \emph{positional triple}
and in Sketch B as an \emph{option-value association}. A caller who writes one
and reads the other gets nonsense.
\item \texttt{AnalyticGerm} appears in Sketch A with the argument list
\texttt{[functionDescriptor, center, branch, domain]} and in Sketch C with
\texttt{[center, definition, initialData, domain, branch]}. The arities and the
argument orders both differ.
\end{itemize}
Any phrasing implying that \texttt{SurrealDomain}, \texttt{CertifiedHahn},
\texttt{HahnParent}, \texttt{SurrealRoot} or \texttt{CoherentLift} exists in
code would be false three times over.
\end{remark}

\subsection{Three result protocols, side by side}
\label{cas:sub-protocols}
The three sources propose result protocols at three different granularities of
the same invariant. All three are printed; none is picked.

\begingroup\small
\begin{longtable}{@{}P{0.155\textwidth}P{0.40\textwidth}P{0.38\textwidth}@{}}
\toprule
Origin & Result kinds & Notes\\
\midrule\endhead
Capability tiers & \textbf{Five-valued}: \texttt{Exact} /
\texttt{CertifiedApproximation} / \texttt{Conditional} / \texttt{Unknown} /
\texttt{Failure} & Comes with named \emph{standing obligations} that may attach
to a conditional node: a positive argument for a real logarithm; a finite real
part for canonical strip sine; positive support for an infinitesimal series
substitution; a root selector valid in the current ordered field; a common
ordinary domain for coherent composition; a compatible phase character for
infinite circular input. Division by a \emph{certified} zero is an error;
division by an object of \emph{unknown} zero status creates a nonzero
obligation.\\[0.5em]
Layered representations & \textbf{Boolean-plus-witness} /
\texttt{Undetermined[reason, obligations]} / \texttt{Failure} & A successful
equality request returns a Boolean and a checkable witness; an unsupported or
exhausted request returns \texttt{Undetermined}; invalid input returns a
failure.\\[0.5em]
Gaussian prototype & \textbf{Four result kinds}: invalid input / undefined
mathematical operation / unsupported algorithm / unknown proposition & A failed
zero test must not return \texttt{False}; a negative power of zero is genuinely
undefined; a square root outside a rational parent asks for an extension rather
than proving impossibility.\\
\bottomrule
\end{longtable}
\endgroup

The common invariant is the one stated in \cref{cas:sub-requirements}: an
unproved equality is not a proved inequality, and a resource limit is not a
mathematical counterexample. Each operation should additionally report the
capability actually used --- a normalization can succeed in the rational backend
while leaving one transcendental coefficient symbolic; a positivity query can
succeed from a certified leading coefficient without full normalization; a
valuation cutoff can produce a compressed exact block rather than a flat list.
These are useful, distinct \emph{successful} results.

\subsection{Component responsibilities}
\label{cas:sub-components}
\begingroup\small
\begin{longtable}{@{}P{0.24\textwidth}P{0.70\textwidth}@{}}
\toprule
Component & Responsibilities and explicit guarantees\\
\midrule\endhead
Coefficient backend & Exact arithmetic; zero and sign decisions where available;
ordinary analytic identities; coefficient enclosures; conversion to larger
coefficient domains; whether its ordinary constants admit certified numerical
enclosures.\\[0.4em]
Exponent (monomial) backend & Addition, equality, mathematical order, rational
division when supported, rational-root promotion, order-preserving changes of
scale, and the embedding used to interpret it inside the surreal field.\\[0.4em]
Exact algebra & Sparse sums, rational expressions, finite algebraic extensions,
polynomial and matrix operations.\\[0.4em]
Support engine & Certified constructors; finite decomposition bounds;
coefficient production; \emph{separate} membership, coefficient, decomposition,
leading-term and prefix services, with a declaration of which are total; and a
semantic proof of well ordering kept distinct from algorithms.\\[0.4em]
Precision engine & Valuation bounds, generator-index jets, ordinary coefficient
error models, refinement, and dependence tracking.\\[0.4em]
Function layer & Formal germs, ordinary domains, Hahn-coherent data, radius-free
germs, branch selectors, derivative modes, admissible composition rules, and
explicitly selected transserial interpretations.\\[0.4em]
Presentation layer & Readable notation, lossless serialization, evidence
inspection, and safe adapters to native CAS facilities.\\
\bottomrule
\end{longtable}
\endgroup

An operation result should preserve its exact origin and its obligations. If the
answer is a truncation it must include a remainder descriptor and the error
calculus used to obtain it. If the answer is conditional, the condition must
remain visible in further transformations. If an unsupported request has a
well-defined semantic node, the system may retain that node without pretending
to have evaluated or normalized it.

\subsection{Proof-carrying expressions without a proof assistant}
\label{cas:sub-proofcarrying}
A practical certificate can be a small derivation tree built from trusted
constructors, with a checker smaller and more deterministic than the
expression-search engine. A proof-oriented interface could export finite
polynomial identities, support derivations, valuation inequalities and
root-selection certificates to such a checker, which would verify specific
claims without needing to decide every question about every surreal object.

\begin{remark}[The trust boundary]\label{cas:rem-trust}
This does \textbf{not} amount to a machine-checked formalization of the surreal
class. Where certificates rely on trusted mathematical theorems from the
supplied manuscripts, \emph{that trust boundary must be documented}. A later
formal checker could verify the finite algebraic kernel and the soundness of
selected support constructors without requiring every analytic extension to be
formalized at once. No such formalization exists in any of the three delivered
packages (\cref{cas:sec-executed}).
\end{remark}

\section{The three delivered prototypes}
\label{cas:sec-prototypes}

\begin{quote}\small\itshape
There are \textbf{three} shipped Wolfram Language packages here, not one. They
have similar names --- two of them are near-anagrams --- similar API verbs, and
all three are described in their own source articles as implementing ``the
effective rational monomial core''. They are nevertheless three different
artifacts with three different envelopes. This report never writes ``the
prototype''. Each is named: \textbf{RationalHahn}, \textbf{SurrealCASCore},
\textbf{HahnRational}. No capability of one may be attributed to another.
\end{quote}

\subsection{RationalHahn --- the capability-tiers kernel, over \texorpdfstring{$\Q$}{Q}}
\label{cas:sub-rationalhahn}
\textbf{File.} \path{RationalHahn.wl}, 78 lines, approximately 4.3~KB; context
\texttt{RationalHahn`}; shipped with \path{Checks.wl} (34 checks) and a driver
\path{RunChecks.wl}. In this repository the three files are
\begin{sloppypar}\noindent
\path{code/02-capability-tiers-and-rational-kernel-rationalhahn.wl},
\path{code/02-capability-tiers-and-rational-kernel-checks.wl} and
\path{code/02-capability-tiers-and-rational-kernel-runchecks.wl}.
\end{sloppypar}

\textbf{Coefficient field: $\Q$, and only $\Q$.} The predicate is
\texttt{rationalQ[c] := IntegerQ[c] || Head[c] === Rational}, applied through
\texttt{CoefficientRules} to every coefficient of both numerator and
denominator. There is \textbf{no $i$, no Gaussian rationals, no complex
arithmetic and no surcomplex pairs anywhere in this kernel.} Machine reals
(\texttt{1.0}) and symbolic transcendentals (\texttt{Pi}) are rejected with
\texttt{Failure["CoefficientDomain"]}. The article's design target $\RA$ is
explicitly \emph{not} what the kernel implements: ``An initial program may use
$\Q$ instead, as our reference kernel does.'' \emph{Anything this kernel proves
or demonstrates is a statement about a real kernel.}

\textbf{Exponent lattice: integer vectors $\Z^d$.} Elements are quotients of
ordinary multivariate polynomials in $d$ formal symbols, so monomial exponents
come from \texttt{CoefficientRules} in $\N^d$ and valuations lie in $\Z^d$.
There are no rational or fractional exponents, no ramification, and no
$t^{1/2}$ anywhere in this kernel; Puiseux data and root selectors appear only
as later design layers in the prose.

\textbf{Rank and parent identity.} Arbitrary finite $d$, fixed at construction
by the variable list passed to \texttt{RHMake}. Domain identity is
\texttt{SameQ} on that variable list (\texttt{sameDomainQ}), so there is no
automatic embedding between domains: \texttt{RHAdd} of a \texttt{\{t,s\}} object
and a \texttt{\{t\}} object returns \texttt{Failure["DomainMismatch"]}. The
shipped check suite uses $d=2$ with variables \texttt{\{t,s\}}.

\textbf{Monomial interpretation.} The ordered variables $\{t_1,\dots,t_d\}$
denote Conway monomials $t_j=\omega^{-\omega^{j-1}}$, so
$v(t_1)=1,\ v(t_2)=\omega,\ \dots,\ v(t_d)=\omega^{d-1}$ and
$0<t_{j+1}<t_j^{\,n}$ for every ordinary positive integer $n$. They are
explicitly not ordinary numeric variables with small floating-point values.

\textbf{Coordinate significance: REVERSED.} The leading term is the
\texttt{CoefficientRules} entry whose exponent vector is least \emph{after
reversing} its coordinates, so the \emph{highest-index} (highest-rank, most
infinitesimal) variable is the first comparison coordinate. Literally:
\begin{lstlisting}
leadingPolynomial[p_, variables_List] := Module[{r},
  r = First[SortBy[CoefficientRules[p, variables],
                   Reverse[First[#]] &]];
  {First[r], Last[r]}];
\end{lstlisting}
Standard part uses \texttt{vectorSign}, the sign of the first nonzero entry of
the \emph{reversed} vector: positive means infinitesimal, so $\st=0$; the zero
vector returns the leading coefficient; negative returns
\texttt{Failure["NotFinite"]}.

\textbf{Data model.} \texttt{RHData[variables, numerator, denominator]}, a
private head; fractions are reduced by \texttt{Cancel} at construction and the
parts are \texttt{Expand}ed.

\textbf{Exported operations (13, with usage messages).} \texttt{RHData} (private
data head), \texttt{RHMake}, \texttt{RHExpression}, \texttt{RHAdd},
\texttt{RHNegate}, \texttt{RHSubtract}, \texttt{RHMultiply},
\texttt{RHInverse}, \texttt{RHSign} (exact $-1/0/1$), \texttt{RHCompare},
\texttt{RHLeadingData} (returns $\{$valuation vector, leading coefficient$\}$;
\texttt{Missing["ZeroHasNoLeadingTerm"]} for zero), \texttt{RHValuation}
(\texttt{Infinity} only as the distinguished valuation marker for zero),
\texttt{RHStandardPart}.

\textbf{Failure taxonomy.} \texttt{Variables} (empty, duplicate or non-symbol),
\texttt{CoefficientDomain}, \texttt{ZeroDenominator}, \texttt{DomainMismatch},
\texttt{DivisionByZero}, \texttt{NotFinite}.

\textbf{Not implemented} (stated in the package comment, the README and the
article): the effective real closure $R_d$, complex pairs and $C_d$, algebraic
root extensions and root selectors, arbitrary exponent groups, lazy or arbitrary
infinite Hahn supports, transcendental functions, general complex powers, and
proof checking. These are subsequent design layers, \emph{not hidden features}.

\textbf{Session, in this kernel's own convention.}
\begin{lstlisting}
Get["code/RationalHahn.wl"];
Clear[t, s];                 (* t = omega^-1 ; s = omega^(-omega) *)
a = RHMake[1 + t, t, {t, s}];
b = RHMake[s, t^10000, {t, s}];
RHSign[a]                              (* 1 *)
RHValuation[a]                         (* {-1, 0} *)
RHStandardPart[b]                      (* 0 *)
RHCompare[b, RHMake[1, 1, {t, s}]]     (* -1 *)
RHExpression[RHInverse[a]]             (* t/(1 + t) *)
RHValuation[RHMake[(1 + t) s - s, 1, {t, s}]]   (* {1, 1} *)
\end{lstlisting}
The last line is the cancellation trap of \cref{cas:sub-cancellationtrap}: the
leading exponent is $(1,1)$, not $(0,1)$.

\begin{remark}[The $10000$-power comparison is one instance]
\label{cas:rem-10000}
The suite's check \texttt{RHCompare[s, t\^{}10000] === -1} together with
\texttt{RHSign[t\^{}10000 - s] === 1} is a single successful higher-rank order
calculation. \emph{It is not the proof for every $n$}; the universal statement
follows from the exponent-group definition and \cref{cas:thm-core}.
\end{remark}

\subsection{SurrealCASCore --- the layered-representations kernel, over \texorpdfstring{$\Q(i)$ with $\Q^r$}{Q(i) with Q\textasciicircum r} exponents}
\label{cas:sub-surrealcascore}
\textbf{File.} \path{SurrealCASCore.wl}, 82 lines, 4689 bytes; context
\texttt{SurrealCASCore`}; shipped with the harness \path{TestCore.wls} and an
independent SymPy script \path{verify_examples.py}. In this repository:
\path{code/04-layered-representations-and-adapter-surrealcascore.wl},
\path{code/04-layered-representations-and-adapter-testcore.wls},
\path{code/04-layered-representations-and-adapter-verify-examples.py}.

\textbf{Coefficient field: exact Gaussian rationals $\Q(i)$.} The predicate is
\texttt{gaussianQ[x] := ratQ[Re[x]] \&\& ratQ[Im[x]]} with
\texttt{ratQ[x] := MatchQ[x, \_Integer | \_Rational]}. Machine reals are
rejected (\texttt{HCreate[\{\{\{1\},0.5\}\},1]} is a tested \texttt{Failure}).
So: $\Q(i)$, not $\Q$, not $\RA$, not algebraic numbers.

\textbf{Exponent lattice: RATIONAL vectors $\Q^r$.} Each exponent is a list of
length exactly $r$ whose entries each satisfy \texttt{ratQ}. Irrational exponent
coordinates are rejected (\texttt{HCreate[\{\{\{Sqrt[2]\},1\}\},1]} is a tested
\texttt{Failure}), and the suite \emph{exercises} half-integer exponents ---
$1/2$, and \texttt{RandomInteger[\{-4,6\}]/2}. \textbf{This is the only one of
the three kernels with fractional exponents.}

\textbf{Rank and parent identity.} $r$ is a fixed positive ordinary
\emph{integer} field stored inside \texttt{HData[rank, terms]}; the guard is
\texttt{r\_Integer /; r > 0}. Parent identity is equality of that integer.
Operations between data of different rank are rejected, with no automatic
coercion or rank-raising. The tests use $r=1$ and $r=2$ only.

\textbf{Coordinate significance: plain lex, FIRST coordinate most significant.}
Implemented by \texttt{lexLess}, which finds the first index where the vectors
differ and compares there. Terms are stored sorted ascending by exponent, so the
leading (least) exponent sits at position \texttt{[[2,1,1]]}. The surreal
interpretation, stated in the package header, the README and the article, is
$(q_1,\dots,q_r)\mapsto q_1\omega^{r-1}+\cdots+q_r$ with Hahn monomial
$\omega^{-\gamma}$ at exponent $\gamma$, so \emph{a larger exponent vector means
a smaller number}; for $r=2$ the generators are $\omega^{-\omega}$ (exponent
$(1,0)$) and $\omega^{-1}$ (exponent $(0,1)$).

\textbf{Data model: TWO representations.} A canonical sparse map
\texttt{HData[rank, terms]} from exponent vector to coefficient ---
\texttt{HCreate} validates all terms, \emph{gathers} coincident exponents and
\emph{sums} their coefficients, \emph{drops} zero coefficients, and sorts by
\texttt{lexLess}, with \texttt{validQ} re-running \texttt{HCreate} to check
canonicity of any datum handed to an operation. And separately a \emph{fraction
wrapper} \texttt{RData[num, den]} of two such sums. \textbf{The fraction wrapper
is deliberately NOT gcd-normalized}: structurally different \texttt{RData}
wrappers may be equal, and \texttt{REqual}, not \texttt{SameQ}, is the intended
comparison. This is exactly the canonicality/non-canonicality contrast of
\cref{cas:sub-promises}, present inside one package.

\textbf{Exported operations (17 with usage messages).} Sparse layer:
\texttt{HCreate}, \texttt{HAdd}, \texttt{HNegate}, \texttt{HMultiply} (naive
pairwise exponent sums, $O(mn)$ before collection), \texttt{HPower}
(\emph{nonnegative} integer exponent only, by \texttt{Nest} --- deliberately not
an optimized exponentiation), \texttt{HConjugate} (coefficientwise Gaussian
conjugation), \texttt{HValuation}, \texttt{HLeadingCoefficient},
\texttt{HCompareReal} (returns $-1/0/1$ via the sign of the leading coefficient
of $a-b$, and returns \texttt{Failure["NonrealInput"]} if \emph{any} coefficient
of either operand has nonzero imaginary part). Fraction layer:
\texttt{RCreate} (rejects a zero denominator), \texttt{RAdd},
\texttt{RMultiply}, \texttt{RInverse}, \texttt{REqual} (cross multiplication),
\texttt{RValuation} (componentwise difference, so it returns a \emph{vector},
e.g.\ \texttt{\{-2\}}). Plus the two private data heads \texttt{HData} and
\texttt{RData}.

\textbf{Absent from this kernel.} \emph{No standard part.} No \texttt{HRe},
\texttt{HIm} or squared modulus. No inverse on the sparse layer --- only
\texttt{RInverse} on fractions.

\textbf{Not implemented} (stated identically in the package comment, the README,
the test report and the article): arbitrary surreal cuts, algebraic coefficient
extensions, algebraic root isolation and selectors, infinite series generators,
certified analytic domains, general Hahn-coherent functions, the transserial
layer, and standard part. No definitions are attached to \texttt{System`Plus},
\texttt{System`Times}, \texttt{System`Power}, \texttt{System`Infinity},
\texttt{System`N} or \texttt{NumericQ}.

\textbf{Session, in this kernel's own convention.}
\begin{lstlisting}
Get["SurrealCASCore.wl"];
one = HCreate[{{{0}, 1}}, 1];
t   = HCreate[{{{1}, 1}}, 1];
a   = HAdd[one, HNegate[t]];
b   = HAdd[one, t];
HMultiply[a, b]
(* HData[1, {{{0}, 1}, {{2}, -1}}] *)
REqual[RCreate[one, a], RCreate[b, HMultiply[a, b]]]
(* True *)

s = HCreate[{{{0, 1}, 1}}, 2];
u = HCreate[{{{1, 0}, 1}}, 2];
HCompareReal[HPower[s, 100], u]
(* 1 *)
\end{lstlisting}
The first result is the exact identity $(1-t)(1+t)=1-t^2$. The second checks
$1/(1-t)=(1+t)/(1-t^2)$ \emph{without expanding either geometric series}; no
infinite summation routine is needed. In the third, $s=\omega^{-1}$ and
$u=\omega^{-\omega}$, and the output $1$ means $s^{100}>u$.

\begin{remark}[The exponent-$100$ test is one instance]\label{cas:rem-100}
More generally $s^n>u$ for every positive ordinary integer $n$, because
$(0,n)<(1,0)$ lexicographically. \emph{That general assertion follows from the
domain's order definition, not from extrapolating the particular test with
exponent $100$.}
\end{remark}

\subsection{HahnRational --- the Gaussian-prototype kernel, over \texorpdfstring{$\Q(i)$ with $\Z^d$}{Q(i) with Z\textasciicircum d} exponents}
\label{cas:sub-hahnrational}
\textbf{File.} \path{HahnRational.wl}, 97 lines, 5666 bytes; context
\texttt{HahnRational`}; shipped with \path{test_HahnRational.wl} (88 checks), a
short demo \path{examples.wl}, and an independent SymPy script
\path{verify_examples.py}. In this repository:
\path{code/05-gaussian-rational-monomial-prototype-hahnrational.wl} and
companions with the same prefix.

\textbf{Coefficient field: $\Q(i)$.} Enforced by
\texttt{gaussianQ[c] := MatchQ[c, \_Integer | \_Rational | Complex[\_Integer|\_Rational, \_Integer|\_Rational]]},
applied to every coefficient of both parts at construction. \texttt{Pi}, machine
floats and any transcendental coefficient are rejected with
\texttt{Failure["CoefficientDomain", ...]}. Not plain $\Q$; not algebraic
numbers; not \texttt{Root} objects.

\textbf{Exponent lattice: INTEGER vectors $\Z^d$ only}, of rank
$d=\mathrm{Length}$ of the declared variable list. \textbf{Rational exponents
are not supported and are actively rejected}: \texttt{HCreate[\{u,t\},Sqrt[t]]}
is a tested rejection --- fractional powers need a parent extension.
Rational-exponent lattice refinement is discussed in the article as future work
only.

\textbf{Rank and parent identity.} Arbitrary finite $d$, fixed by the declared
variable list; parent identity is \texttt{SameQ} on that list
(\texttt{sameParent}). \textbf{A reversed variable list is a DIFFERENT parent}
and mixing is rejected: \texttt{HAdd[h[t], HCreate[\{t,u\},t]]} fails with
\texttt{ParentMismatch}. Tests and examples use $d=2$.

\textbf{Coordinate significance: plain lex, LEFTMOST coordinate most
significant}, with variables listed from the \emph{most significant} valuation
coordinate to the least. \texttt{polyLeading} is
\texttt{First[SortBy[CoefficientRules[Expand[p], vars], First]]} --- with
\emph{no} reversal. For the parent \texttt{\{u,t\}}: $v(u)=\{1,0\}$,
$v(t)=\{0,1\}$, with the concrete surreal reading $t=\omega^{-1}$ and
$u=t^\omega=\omega^{-\omega}$; \emph{larger valuation means smaller magnitude},
so $u<t^n$ for every ordinary positive integer $n$.

\textbf{Data model.} \texttt{HR[vars, numerator, denominator]}, a private head,
holding normalized expanded polynomials with the denominator's
lexicographically-leading coefficient \emph{normalized to $1$}. Construction
goes through \texttt{Cancel[Together[expr]]} and requires \texttt{PolynomialQ}
of both parts in the declared variables. This is \emph{not} an expanded Hahn
series.

\textbf{Exported operations (19 with usage messages).} \texttt{HCreate} (the
\emph{only} public value constructor), \texttt{HExpression},
\texttt{HVariables}, \texttt{HAdd}, \texttt{HMul}, \texttt{HInv},
\texttt{HPower} (integer exponent; $n=0$ returns $1$; a negative power of zero
is rejected), \texttt{HScale} (by a $\Q(i)$ scalar), \texttt{HEqualQ}
(cross-multiply and \texttt{Expand}; does not require gcd), \texttt{HZeroQ},
\texttt{HValuation} (numerator minus denominator exponent vectors;
\texttt{Infinity} for zero, as metadata only),
\texttt{HLeadingCoefficient}, \texttt{HConjugate}, \texttt{HRe}, \texttt{HIm},
\texttt{HNormSquared} ($a\,\overline a$, no square root), \texttt{HSign}
(rejects input not equal to its own conjugate), \texttt{HCompare} (returns
$-1/0/1$; \emph{rejects nonreal inputs even when their difference is real}),
\texttt{HStandardPart} ($0$ for infinitesimals; the leading coefficient ---
possibly a Gaussian rational such as $2+3i$ --- at valuation-sign zero; and
\texttt{Failure["InfiniteValue"]} for infinite values).

\textbf{This is the only one of the three with the full surcomplex surface}:
conjugation, real and imaginary parts, squared modulus, sign and comparison with
nonreal rejection, and standard part.

\textbf{Not implemented} (stated by both the README and the article): general
Hahn-series coefficient programs, rational-exponent parent refinement, algebraic
closure, algebraic root selection, arbitrary analytic germs, transseries, global
exponential extensions, unrestricted modulus (only \emph{squared} modulus),
contour integration, AG's germ division and preparation, and CT's integration
theory.

\textbf{Session, in this kernel's own convention.}
\begin{lstlisting}
Get["code/HahnRational.wl"];
Clear[u, t];
a = HCreate[{u, t}, 1/(1-t) + u];
b = HCreate[{u, t}, 1/(1-t)];
HValuation[HAdd[a, HScale[b, -1]]]
(* {1, 0} *)
HCompare[HCreate[{u, t}, u], HCreate[{u, t}, t^100]]
(* -1 *)

z = HCreate[{u, t}, t + I u];
HExpression[HNormSquared[z]]
(* t^2 + u^2 *)
HStandardPart[HCreate[{u, t}, (2 + 3 I + t)/(1 - u)]]
(* 2 + 3 I *)
HStandardPart[HCreate[{u, t}, t/u]]
(* Failure["InfiniteValue", ...] *)
\end{lstlisting}
The first operation \emph{cancels an entire infinite geometric block exactly}
and reveals a term at a new valuation rank; a coefficient stream with no way to
recognize the equal geometric tails could fail to reach it. The second
comparison is exact and does not depend on choosing a large real value for
$\omega$. The quotient $t/u$ is infinite and has no finite standard part in this
convention, so the package returns a \texttt{Failure} rather than a
floating-point infinity.

\subsection{Comparison, symbol collisions, and four cross-attribution warnings}
\label{cas:sub-comparison}

\subsubsection{The full property table}
\label{cas:subsub-table}
\begingroup\footnotesize
\begin{longtable}{@{}P{0.155\textwidth}P{0.245\textwidth}P{0.245\textwidth}P{0.245\textwidth}@{}}
\toprule
Property & \textbf{RationalHahn} (02) & \textbf{SurrealCASCore} (04) &
\textbf{HahnRational} (05)\\
\midrule\endhead
File and size & \texttt{RationalHahn.wl}, 78 lines, $\approx$4.3 KB &
\texttt{SurrealCASCore.wl}, 82 lines, 4689 B &
\texttt{HahnRational.wl}, 97 lines, 5666 B\\[0.4em]
Context & \texttt{RationalHahn`} & \texttt{SurrealCASCore`} &
\texttt{HahnRational`}\\[0.4em]
Data model & Reduced quotient \texttt{RHData[vars,num,den]}; \texttt{Cancel} at
construction & \textbf{Two}: canonical sparse \texttt{HData[rank,terms]} plus a
deliberately \emph{unreduced} fraction \texttt{RData[num,den]} &
Reduced quotient \texttt{HR[vars,num,den]};
\texttt{Cancel[Together[...]]}, denominator leading coefficient scaled to
$1$\\[0.4em]
Coefficient field & \textbf{$\Q$ only} --- no $i$ anywhere & $\Q(i)$ & $\Q(i)$\\[0.4em]
Exponent lattice & $\Z^d$ & \textbf{$\Q^r$ --- rational exponents, actually
exercised} & $\Z^d$\\[0.4em]
Rank / parent identity & Arbitrary $d$; \texttt{SameQ} on the variable list &
An integer field $r$; equality of that integer & Arbitrary $d$; \texttt{SameQ}
on the variable list\\[0.4em]
Coordinate significance & \textbf{REVERSED}: highest-index coordinate compared
first & Plain lex, first coordinate most significant &
Plain lex, first coordinate most significant\\[0.4em]
Monomial reading & $t_j=\omega^{-\omega^{j-1}}$; $v(t_1)=1$, $v(t_2)=\omega$ &
$(q_1,\dots,q_r)\mapsto q_1\omega^{r-1}+\cdots+q_r$ &
$(a_1,\dots,a_d)\mapsto a_1\omega^{d-1}+\cdots+a_d$; $u=\omega^{-\omega}$,
$t=\omega^{-1}$\\[0.4em]
Standard part & \textbf{Yes} (\texttt{RHStandardPart}) & \textbf{No} &
\textbf{Yes} (\texttt{HStandardPart})\\[0.4em]
Conjugation / Re / Im / $|z|^2$ & \textbf{None at all} &
\texttt{HConjugate} only & \texttt{HConjugate}, \texttt{HRe}, \texttt{HIm},
\texttt{HNormSquared}\\[0.4em]
Real comparison & \texttt{RHSign}, \texttt{RHCompare} &
\texttt{HCompareReal}, rejecting nonreal input &
\texttt{HSign}, \texttt{HCompare}, both rejecting nonreal input\\[0.4em]
Inverse & \texttt{RHInverse} & \textbf{Only} \texttt{RInverse}, on fractions &
\texttt{HInv}\\[0.4em]
gcd-reduced fractions & \textbf{Yes} & \textbf{No} --- use \texttt{REqual}, not
\texttt{SameQ} & \textbf{Yes}\\[0.4em]
Exported symbols & 13 & 17 & 19\\[0.4em]
Rejections tested & machine reals, \texttt{Pi}, \texttt{Sin[t]}, zero
denominator, zero inverse, domain mismatch, duplicate/empty variables &
machine reals, \texttt{Sqrt[2]} exponents, rank mismatch, negative power, zero
denominator, zero inverse, nonreal ordering &
\texttt{Pi}, machine reals, \texttt{Exp[t]}, \textbf{\texttt{Sqrt[t]}} (needs a
parent extension), parent mismatch, zero inverse, negative power of zero,
nonreal ordering\\[0.4em]
Not implemented & real closure, complex pairs, algebraic roots, arbitrary
exponent groups, lazy supports, transcendental functions, general complex
powers, proof checking &
algebraic coefficient fields, algebraic root selectors, infinite series
generators, certified analytic domains, arbitrary Conway cuts, transserial
layer, \textbf{standard part} &
general Hahn coefficient programs, \textbf{rational-exponent refinement},
algebraic closure, arbitrary analytic germs, transseries, global exponential
extensions, contour integration\\[0.4em]
What was executed & \textbf{Definitions submitted as code}; 31/34 then three
corrected test expressions; the shipped file does \emph{not} parse &
\textbf{Real files \texttt{Get}-loaded}; 110/110; plus SymPy 69/69 &
\textbf{\texttt{Get[StringToStream[...]]}}; 88/88; plus SymPy 323/323\\
\bottomrule
\end{longtable}
\endgroup

\subsubsection{The symbol-name collision}
\label{cas:subsub-collision}
\begin{remark}[Six colliding public symbols]\label{cas:rem-symbolcollision}
\textbf{SurrealCASCore and HahnRational both declare, with \texttt{::usage},
the six public symbols}
\[
 \texttt{HAdd},\quad \texttt{HConjugate},\quad \texttt{HCreate},\quad
 \texttt{HLeadingCoefficient},\quad \texttt{HPower},\quad \texttt{HValuation},
\]
\textbf{with incompatible signatures.} The clearest instance is the
constructor:
\begin{itemize}
\item SurrealCASCore:
\texttt{HCreate[\{\{exponentVector, coefficient\},\dots\}, rankInteger]}
\item HahnRational:
\texttt{HCreate[\{u,t,\dots\}, rationalExpression]}
\end{itemize}
They also
diverge on the names of operations that do the same job ---
\texttt{HMultiply} versus \texttt{HMul}, \texttt{HCompareReal} versus
\texttt{HCompare}, \texttt{RInverse} versus \texttt{HInv}.
\textbf{Loading both packages into one session silently breaks both}, since one
set of definitions shadows the other.

Additionally, note the \emph{name} clash at the file level:
\path{RationalHahn.wl} (02) and \path{HahnRational.wl} (05) are near-anagrams of
each other and are two entirely different artifacts.
\end{remark}

\subsubsection{Four cross-attribution warnings}
\label{cas:subsub-crossattribution}
\begin{enumerate}
\item[(a)] \textbf{$\Q(i)$, conjugation and surcomplex operations belong to
SurrealCASCore and HahnRational, and NOT to RationalHahn}, whose kernel is
$\Q$-only and has no conjugation, no real or imaginary part, and no modulus of
any kind.
\item[(b)] \textbf{Rational and fractional exponents belong to SurrealCASCore
ONLY.} RationalHahn and HahnRational both have integer exponent lattices, and
HahnRational's suite \emph{actively rejects} \texttt{Sqrt[t]} as needing a
parent extension. So a construction needing a fractional exponent --- a ramified
root, a Puiseux-type uniformizer --- is representable in SurrealCASCore and not
in the other two.
\item[(c)] \textbf{Standard part belongs to RationalHahn and HahnRational, and
NOT to SurrealCASCore}, which has no standard-part operation at all.
\item[(d)] \textbf{gcd-reduced fractions belong to RationalHahn and
HahnRational, and NOT to SurrealCASCore}, whose fraction wrapper is deliberately
unreduced, so that structurally different wrappers may be equal and
\texttt{REqual} rather than structural equality is the intended comparison.
Correspondingly, claims about ``not expanding an infinite tail'' by means of a
separate \emph{fraction object} layered over a canonical sparse sum belong to
SurrealCASCore; the other two carry a rational function directly.
\end{enumerate}

\begin{remark}[The sentence never to write]\label{cas:rem-neverwrite}
Any sentence of the form ``the prototype implements the rational monomial core
over $\Q(i)$ with rational lexicographic exponents, conjugation and standard
part, and passes its checks'' is \textbf{true of no delivered artifact}. It
silently claims a union of capabilities that none of the three has.
\end{remark}

\subsubsection{The reversed variable-order convention}
\label{cas:sub-orderconventions}
\begin{remark}[The conventions are reversed between two of the three kernels]
\label{cas:rem-reversal}
\textbf{RationalHahn lists variables least-significant-valuation first;
HahnRational lists them most-significant first.} Verified directly against the
sources:
\begin{itemize}
\item RationalHahn: $v(t_1)=1$, $v(t_2)=\omega$, \dots, $v(t_d)=\omega^{d-1}$,
so $0<t_{j+1}<t_j^{\,n}$. \emph{The first-listed variable is the largest element
and has the least valuation.} Exponent vectors are compared after
\texttt{Reverse}.
\item HahnRational: ``Variables are listed from the most significant valuation
coordinate to the least. For \texttt{\{u,t\}}, \texttt{u} is smaller than every
positive power of \texttt{t}.'' \emph{The first-listed variable is the smallest
element and has the greatest valuation.} Checked by hand: $v(u)=\{1,0\}$,
$v(t^n)=\{0,n\}$, and lexicographically $\{1,0\}>\{0,n\}$, hence $u<t^n$.
\end{itemize}
HahnRational's own README warns that \emph{a reversed variable list is a
different parent}. \textbf{A reader who carries one convention into the other
gets every inequality backwards.} SurrealCASCore agrees with HahnRational on
coordinate significance, and the two agree on the interpretation
$(q_1,\dots,q_r)\mapsto q_1\omega^{r-1}+\cdots+q_r$; they differ only on whether
the $q$ are rational or integral.
\end{remark}

All three kernels express the same mathematical fact,
$\omega^{-\omega}<\omega^{-n}$ for every ordinary $n$, under two different
coordinate layouts and --- in SurrealCASCore's case --- with reversed argument
order in the comparison call. The three displayed comparisons are:
\begin{itemize}
\item \textbf{RationalHahn}: \texttt{RHCompare[s, t\^{}10000] === -1}, together
with \texttt{RHSign[t\^{}10000 - s] === 1}.
\item \textbf{SurrealCASCore}: \texttt{HCompareReal[HPower[s,100], u] === 1},
meaning $s^{100}>u$ --- note the reversed argument order relative to the other
two.
\item \textbf{HahnRational}: \texttt{HCompare[h[u], h[t\^{}100]] === -1}.
\end{itemize}
\textbf{None of these has been rewritten into a common form}, and this report
states no single global coordinate convention that these quoted outputs would
then contradict. The global rule stated once in \cref{cas:rem-convention} is
that coordinate significance is \emph{declared as part of the domain
descriptor}; each kernel's own declaration is given in its own subsection above.

\subsection{Shared contract boundaries of all three prototypes}
\label{cas:sub-sharedboundaries}
Six statements hold of all three delivered packages and are recorded here once.

\begin{enumerate}
\item \textbf{None of the three is the architecture.} Each source article states
this in its own text, and \cref{cas:sec-architecture} repeats it: the proposed
object layers, capability queries and result protocols are \emph{proposed
interfaces}, and most of that architecture is exported by no package at all. A
small, separately delimited prototype demonstrating exact sparse or rational
arithmetic does not purport to implement it.
\item \textbf{Only the public constructors' outputs may be passed to the
arithmetic routines.} The internal heads --- \texttt{RHData} in RationalHahn,
\texttt{HData} and \texttt{RData} in SurrealCASCore, \texttt{HR} in
HahnRational --- are private. \emph{Manually forging them bypasses validation
and lies outside each package's own contract.}
\item \textbf{None of the three is a hardened parser} for arbitrary or untrusted
Wolfram expressions, and arbitrary malformed inputs are not handled as a
production-grade general symbolic language would handle them.
\item \textbf{Rejection is a statement about the core, not about mathematics.}
When a kernel rejects an input --- \texttt{Pi}, a machine float,
\texttt{Exp[t]}, \texttt{Sqrt[t]}, an irrational exponent coordinate, a
mismatched parent --- that means \emph{this core does not implement that
representation}, not that the corresponding exact surreal value cannot be
defined. The first belongs in a larger coefficient parent, the second in an
exponent or algebraic extension, the third in a functional or lazy-series
extension.
\item \textbf{The three non-implementation lists stand separately and must not
be pooled.} RationalHahn omits the real closure, complex pairs, algebraic root
extensions, arbitrary exponent groups, lazy infinite supports, transcendental
functions, general complex powers and proof checking. SurrealCASCore omits
algebraic coefficient fields, algebraic root selectors, infinite series
generators, certified analytic domains, arbitrary Conway cuts, the transserial
layer and \emph{standard part}. HahnRational omits general Hahn coefficient
programs, \emph{rational-exponent parent refinement}, algebraic closure,
arbitrary analytic germs, transseries, global exponential extensions and contour
integration. Merging these into one list would understate each package's gaps
and overstate the group's coverage.
\item \textbf{No canonical infinite circular phase is claimed}, by any of the
three, and none of the three implements any phase provider. The trivial
character is a convention and never a default
(\cref{cas:sub-characters,cas:sub-infinitephases}).
\end{enumerate}

\section{What was actually executed}
\label{cas:sec-executed}

\begin{quote}\small\itshape
Three separate provenance records follow, never one aggregate. \textbf{The
Wolfram check counts $34$, $110$ and $88$ must not be added, and the Python
counts $69$ and $323$ must not be added.} One of the source articles explicitly
forbids that arithmetic for its own two suites, and the same reasoning applies
across articles: different suites check different things, and their counts are
not a proof of that many general theorems.
\end{quote}

\subsection{What the three archives recorded}
\label{cas:sub-recorded}

\paragraph{RationalHahn (capability tiers).}
The arithmetic definitions corresponding to \path{RationalHahn.wl} and the
checks corresponding to \path{Checks.wl} were \textbf{submitted as code} to a
stateless connected Wolfram Language evaluator. \emph{The local package path was
explicitly not loaded into a separately installed kernel.} The first evaluation
reported \textbf{34 checks, 31 passed}. Three \emph{test expressions} were then
corrected --- the inverse-cosine endpoint expansion needed the uniformizer
$2\arcsin(u/\sqrt2)/(\sqrt2\,u)$ rather than \texttt{ArcCos[1-u\^{}2]} used
without an explicit positive square root, and the Cayley and
trigonometric-algebraization identities needed \texttt{Together} rather than
\texttt{Cancel} alone on a sum of rational terms --- and all three passed in a
subsequent evaluator call. The distributed suite contains the corrections.

Three disclosures of this record survive verbatim.
\begin{itemize}
\item ``This is not reported as one original 34/34 run.''
\item The first connector evaluation emitted warnings about global-symbol
preprocessing and messages during the inverse-cosine series request, and
\textbf{no claim of a warning-free initial evaluation is made}.
\item An additional request to simplify $\arccos(1-u^2)=2\arcsin(u/\sqrt2)$ for
$0<u<1$ \textbf{remained unevaluated} and \emph{is explicitly not counted as a
passed computer check}; the positive-branch identity is justified by the
ordinary half-angle identity and inverse-function ranges, and a failure to
simplify it is not evidence against that mathematical proof.
\end{itemize}
The promised $34/34$ local run is described only as an \emph{expected}
reproducibility result, not a performed one. Of the 34 checks, only 22 touch the
kernel (eight of those are negative or rejection tests); the remaining 12 are
ordinary \texttt{Series}, \texttt{Together} and \texttt{Simplify} identities
that never call the kernel at all.

\paragraph{SurrealCASCore (layered representations).}
The \emph{actual} \path{SurrealCASCore.wl} and \path{TestCore.wls} files were
transferred as literal text, written into one temporary directory, and
\path{TestCore.wls} was \texttt{Get}-loaded --- it in turn loads the package
from the same directory. The recorded result was
\begin{lstlisting}
<|"Kernel" -> "15.0.1 for Linux x86 (64-bit) (July 2, 2026)",
  "Checks" -> 110, "Passed" -> 110, "Failures" -> {}|>
\end{lstlisting}
with no warning messages on the final run: \textbf{20 deterministic checks plus
90 seeded property checks}, using \texttt{SeedRandom[20260921]}. The record
discloses that \emph{an earlier inline-definition exploratory run passed its 105
checks but produced a static undefined-global-symbol warning, and is explicitly
not the reported result}; five additional deterministic tests were included in
the final run. A separate exact-symbolic script \path{verify_examples.py} was
recorded as running under \textbf{Python 3.13.5 and SymPy 1.14.0} with
\textbf{69 of 69 checks passing}, using \texttt{random.Random(20260921)} and
pinning \texttt{sympy==1.14.0}.

\paragraph{HahnRational (Gaussian prototype).}
The package's executable definitions and the test body were loaded via
\texttt{Get[StringToStream[...]]} in the connected kernel, preserving sequential
package-context parsing. The recorded result was
\begin{lstlisting}
<|"Kernel" -> "15.0.1 for Linux x86 (64-bit) (July 2, 2026)",
  "Tests" -> 88, "Passed" -> 88, "Failures" -> {}|>
\end{lstlisting}
with \texttt{SeedRandom[20260921]} for the randomized polynomial examples. The
record discloses that \emph{a preliminary direct evaluation also passed the same
tests but emitted namespace and documentation setup messages}, and that the
clean package-loading run is the recorded final validation. An independent
Python script was recorded as passing \textbf{323 exact checks} under
\textbf{Python 3.13.5 and SymPy 1.14.0}, covering quotient-algebra matrices,
trace and residue identities, separating denominators, radius-free coefficients
through exponent sixteen, exact geometric remainders, local power-series
identities, and valuation-based contour counts.

\paragraph{Common to all three records.}
The same kernel --- \textbf{Wolfram Language 15.0.1 for Linux x86 (64-bit),
July 2, 2026} --- reached through a \emph{connected evaluator}, on the same date
(September 21, 2026), with \textbf{no local installation exercised by any of the
three}. Two of the three additionally used Python 3.13.5 and SymPy 1.14.0.
\textbf{None of the three contains any Lean file, formalization blueprint or
formalization audit, and no git SHA, Lean toolchain or mathlib revision is
pinned anywhere.}

\subsection{What was independently reproduced for this report}
\label{cas:sub-reproduced}
The following was run directly on this machine on September 21, 2026, and is
\emph{not} carried across from the archives.

\textbf{Environment.} Wolfram Engine at
\path{C:\Program Files\Wolfram Research\Wolfram\15.0\wolframscript.exe},
reporting \textbf{15.0.1 for Microsoft Windows (64-bit) (July 2, 2026)} --- the
same release and build date as the archives' Linux kernel, on a different
platform. \textbf{Python 3.14.4} with \textbf{SymPy 1.14.0}; since the archives
recorded their Python runs under 3.13.5, the Python results reproduce across a
minor Python version as well.

\begingroup\small
\begin{longtable}{@{}P{0.19\textwidth}P{0.17\textwidth}P{0.24\textwidth}P{0.13\textwidth}P{0.14\textwidth}@{}}
\toprule
Package & Wolfram recorded & Wolfram reproduced & Python recorded &
Python reproduced\\
\midrule\endhead
\textbf{RationalHahn} & 34 checks, corrected suite &
\textbf{Does not parse as shipped} (12/34: exactly the 12
checks that never call the kernel); \textbf{34/34} after a
one-character fix &
--- & ---\\[0.45em]
\textbf{SurrealCASCore} & 110/110 & \textbf{110/110} & 69/69 &
\textbf{69/69}\\[0.45em]
\textbf{HahnRational} & 88/88 & \textbf{88/88} & 323/323 &
\textbf{323/323}\\
\bottomrule
\end{longtable}
\endgroup

Every recorded claim reproduces. The quoted sessions of
\cref{cas:sec-prototypes} were also re-executed and produce the outputs shown,
including the cancellation trap $\{1,1\}$, the acceptance of the half-integer
exponent $1/2$ by SurrealCASCore, its rejection of $\sqrt2$, and all three
displayed order comparisons in their own conventions.

\subsection{The one real defect: RationalHahn cannot be loaded as distributed}
\label{cas:sub-defect}
\begin{remark}[A verified packaging defect, found by this merge]
\label{cas:rem-defect}
\path{RationalHahn.wl} \textbf{does not parse}. The kernel's own diagnosis is
\begin{lstlisting}
Syntax::sntue: Unexpected end of file (probably unfinished expression).
(line 79 of .../RationalHahn.wl)
\end{lstlisting}
The cause is in \texttt{RHInverse}, at lines 48--49 as shipped:
\begin{lstlisting}
RHInverse[a_RHData] := If[a[[2]] === 0,
  Failure["DivisionByZero",<||>],RHMake[a[[3]],a[[2]],a[[1]]];
\end{lstlisting}
The \texttt{If[} opened on the first line is \emph{never closed}. A brace census
gives exactly one unmatched \texttt{[} in the whole file (179 against 178,
outside comments and strings), whereas the other two packages balance exactly;
the whole file parses into 38 top-level expressions once a single \texttt{]} is
inserted before the \texttt{;}. \textbf{The fix is one character}:
\texttt{...a[[1]]]];} for \texttt{...a[[1]]];}.

\textbf{Consequence.} \texttt{Get["code/RationalHahn.wl"]}, the driver
\path{RunChecks.wl}, and the README's
\texttt{wolframscript -file code/RunChecks.wl} instruction cannot load the
package as delivered. Running the driver unrepaired on this machine returned
\texttt{Checks -> 34, Passed -> 12}, the twelve passing checks being precisely
the ordinary \texttt{Series}/\texttt{Together}/\texttt{Simplify} identities that
never call the kernel; all 22 kernel checks failed. After the one-character
repair the same driver returned \texttt{Checks -> 34, Passed -> 34, Failures ->
\{\}} with no warnings.

\textbf{Therefore, as delivered, the archive's ``expected 34/34 local run'' is
not merely unverified but currently impossible.} A merged report that repeated
that \texttt{wolframscript} instruction without this note would be printing a
false reproducibility claim; the instruction is repeated in
\cref{cas:app-repro} \emph{only} together with the repair.

\textbf{This is consistent with, and explains, the archive's own verification
record}, which states that the definitions were submitted as code to the
stateless evaluator and that ``the local package path itself was not loaded into
a separately installed kernel''. The package was never \texttt{Get}-loaded
because it \emph{cannot} be. The archive does not claim otherwise; but neither
does it report the defect --- it is not among its list of three corrected
checks, which were all harness-level. Its README also lists
\path{SHA256SUMS.txt} among the distribution files. That is not a defect: the
archive did ship one, and this repository does not carry it, because the
collection ships no checksum manifests. The README is preserved verbatim and so
still names it.

\textbf{SurrealCASCore and HahnRational both \texttt{Get}-load cleanly} (both
bracket-balanced at depth zero), and their reported numbers are the numbers one
gets from running their own harnesses.
\end{remark}

\subsection{What the executed checks do and do not establish}
\label{cas:sub-checkscope}
The executed checks are exact finite algebra and formal-series checks. They are
\textbf{not} a formal verification of: Hahn summability; the surreal
foundations; the effective real-closure construction or real-closed-field
quantifier elimination; global analytic structures; the phase classification or
the $E_\alpha$ family; AG's germ division or preparation; the ISSAC 2026
transseries zero-test; CT's integration theory; or any infinite or transfinite
recursion. A finite series check is evidence about \emph{those coefficients
only}, and a successful driver run does not upgrade finite coefficient
calculations into proofs of the infinite constructions. None of the three
prototypes executes real-closed-field quantifier elimination, constructs all
algebraic roots, solves arbitrary support-recursive equations, or certifies the
supplied manuscripts' analytic results. \textbf{Compatibility with Wolfram
versions other than 15.0.1, and with Mathics, was not tested.} And, once more:
no Lean file, formalization blueprint or formalization audit exists in any of
the three packages, and nothing here is machine-verified.

\section{Implementation priorities, costs, and regression boundaries}
\label{cas:sec-roadmap}

\subsection{A union staged plan}
\label{cas:sub-stages}
All three sources propose staged plans; the union below is their common shape
with each source's distinctive emphasis retained. \emph{The ordering is not a
claim that later stages are easy}; it prevents their difficulties from blocking
a valuable exact algebraic system.

\paragraph{Stage 1 --- exact ordered rational algebra and surcomplex pairs.}
Validated coefficient and exponent domains; rational fractions; arithmetic,
equality, sign, valuation and finite standard parts; surcomplex pairs and exact
squared norms; efficient multivariate rational normalization; exponent-domain
embeddings; expression sharing; exact matrix and polynomial adapters; a stable
serialization format; and a precise public data contract whose \emph{errors} are
tested as thoroughly as its successful identities. This stage already supports
many all-scale geometric calculations, and many polynomial-manuscript examples
including resultants, discriminants, finite quotient algebras and initial
polynomials for coefficients whose valuations are available.

\paragraph{Stage 2 --- algebraic extensions and root selectors.}
Real-root selectors, ordered extension towers, exact comparison, complex
polynomial roots, and their surcomplex pairs; finite-dimensional linear algebra
and basic polynomial tools over those coefficients; rank-one or nested Puiseux
computations. \textbf{Keep root selection independent of ordinary numerical
approximation}: use exact Newton data as a \emph{view} of a root, never as a
substitute for its selector, and preserve branch identity under refinement. The
central deliverable is exact algebraic closure with robust branch selection, not
a universal root expansion.

\paragraph{Stage 3 --- support-certified series, grids, precision and compressed cutoffs.}
Positive-grid constructors (finite unions), finite coefficient-dependency
procedures, reciprocal and binomial expansions, precision propagation,
rank-one finite truncations, rational-grid coefficient extraction, and
compressed cutoff objects. Initially restrict support grammars enough that every
advertised coefficient request terminates, and preserve exact origin graphs
behind truncations. \textbf{Implement \cref{cas:thm-grid}'s dependency guarantee
before offering unrestricted-looking sum and composition syntax.} High-rank
requests should remain compressed or return an informative
unsupported-finite-view result.

\paragraph{Stage 4 --- analytic lifts, coherent data, quotient algebras, residues and periods.}
Ordinary-centre Taylor lifts; finite-angle trigonometry; inverse functions with
endpoint ramification; logarithm branch policies; finite trigonometric
polynomial algebraization; regular implicit germs; Hahn-coherent positive
perturbations with executable coefficient rules; coherent ordinary-contour
functionals for a controlled coefficient-function grammar; finite quotient
algebras, complete-intersection deformation, residue pairings, and rational
contour periods. Then add preparation, division and residue methods from the
supplied sources \emph{under their stated hypotheses}. The arbitrary radius-free
analytic-germ generalization should follow only when an effective ordinary-germ
division backend exists for the chosen language. Reuse host coefficient
functions only with domain checks.

\paragraph{Stage 5 --- exp--log and transserial scales, with phases separate.}
Develop the exp--log hierarchy and its compatible surreal interpretation; admit
forms of transserial composition and differential operations, each with its own
support and leading-term algorithms; keep general semantic exp--log nodes
available even when normalization is incomplete. \textbf{Surcomplex phase
policies, sectorial summation, resurgent special-function extensions,
hyperseries and global chart compatibility are separate projects.} They should
not delay a usable exact algebraic implementation, and must not be marked
implemented merely because the parser accepts their names.

\subsection{Complexity must include representation and output size}
\label{cas:sub-complexity}
Multiplying sparse supports of sizes $m$ and $n$ initially produces at most $mn$
pairs before collection. This is an arithmetic-operation count, not a
bit-complexity estimate: the bit lengths of rational coordinates and algebraic
coefficients matter as much as the number of terms, and exponent comparison and
coefficient growth can dominate the actual cost. Repeated cross multiplication
causes fraction growth, so gcd reduction, common-subexpression sharing and
delayed expansion are important improvements; a straightforward power routine is
intentionally not an optimized exponentiation algorithm. A compact rational
expression can have a very large series expansion, and an algebraic root can
have complicated nested coefficients, so store shared exact graphs, cache
verified leading data, and make expansion demand-driven. Recursive exponent
comparison may dominate an otherwise small calculation.

For grid coefficient extraction the finite box of \cref{cas:thm-grid} proves
termination but can be large; integer linear-equation methods, dynamic
programming and memoized convolution can improve it while preserving exact
treatment of coincident sums. Finite-algebra computations scale with the
quotient dimension and coefficient complexity. Polynomial root descriptions can
have large algebraic degrees, and \emph{quantifier elimination is a sound
foundation, not a guarantee of low computational cost}; specialized Newton and
Hensel methods can provide far better useful views without replacing the exact
fallback.

\begin{remark}[Output-representation limits are not complexity]
\label{cas:rem-outputlimit}
\textbf{No algorithm can print an infinite cutoff prefix as a finite explicit
term list.} This is an output-representation limitation, not merely poor
complexity. Likewise, \textbf{no time budget transforms the halting reductions
of \cref{cas:sec-limits} into decision procedures}. A useful system reports
whether a request is expensive, outside its implemented domain, impossible in
the requested output format, or undetermined by its current evidence. Outside a
decidable domain, every budget exit must remain distinguishable from
mathematical falsehood. A total decision procedure need not be the first
procedure tried: fast domain-specific simplifiers can precede an expensive
complete algebraic fallback.
\end{remark}

Resource accounting should include coefficient bit sizes, polynomial degrees,
extension degrees, exponent-tree sizes, cached coefficient counts and
support-dependency counts. \emph{A single ``precision'' integer is inadequate
for a hierarchy of infinitesimal scales.} Interactive output should show the
leading scales and permit inspection of the exact underlying descriptor.

\subsection{A union regression and acceptance list}
\label{cas:sub-regression}
A regression suite should include much more than successful expansions. The
union of the three sources' lists:

\begin{itemize}
\item Cancellation to an exact finite value; zeros hidden behind long
coefficient prefixes; exact identity tracking proving $A-A=0$.
\item A mathematically valid but \emph{nonenumerable} finite-cutoff request;
preservation of the rank-crossing counterexample
\eqref{cas:eq-geometricremainder}, so that increasing an ordinary integer
precision parameter never claims accuracy beyond $u$; and prevention of a
high-rank cutoff being handled by an endless rank-one loop.
\item Rejection of $\sum_{n\ge1}t^{1/n}$ as an unordered-support Hahn sum;
rejection of positive but non-well-ordered supports and of well-ordered unions
with infinitely many contributions at one exponent.
\item Rejection of the geometric expansion of $\sum t^nz^n$ at $z=t^{-2}$ while
\emph{permitting} the separate rational continuation \eqref{cas:eq-blindsub};
distinguishing $\exp(t)$ from an invalid strong expansion of $\exp(1/t)$.
\item Retention of the difference between an exact factorial Hahn series and an
ordinary divergent series after real specialization.
\item Division by \emph{exact} zero versus division by an object of
\emph{unknown} zero status; standard part of an infinite input; incompatible
exponent domains and parent mismatches.
\item Root branches separated only infinitesimally, which ordinary decimal
neighbourhoods cannot separate; algebraic root modules verifying defining
equations, branch separation, parent embeddings and multiplicities including
repeated roots.
\item Quotient-algebra tests including repeated roots, so that a nondegenerate
\emph{residue} pairing is not confused with a nondegenerate \emph{trace}
pairing; the four-point example; the overdetermined length-drop counterexample;
and the sharp Hensel threshold example \eqref{cas:eq-collision}.
\item Contour modules detecting poles on the contour, using the transformation
determinant $\det\mathsf C$, and distinguishing ordinary-parameter Hahn chains
from fine-continuous paths.
\item Actual versus anchored logarithms; preservation of branch conditions in
logarithmic rewrites; unsupported infinite circular phases returning an
obligation; and a global surcomplex exponential module making the two values of
\eqref{cas:eq-phasecontrast} explicit under the two named policies and
\emph{never caching them as if they belonged to one function}.
\item Differentiation under each named convention; refusal to decide coefficient
equality from insufficient numerical precision; and a check that
\texttt{Normal}-style conversion cannot silently promote a truncation to an
exact equality.
\end{itemize}

For each successful approximate calculation, verify the defining residual and
the claimed error certificate whenever that check is available. A displayed root
expansion without a branch certificate and residual bound is only a
\emph{candidate}; with those certificates it is a dependable computed view of an
exact root object. \textbf{Negative tests are essential}: rejecting an
inadmissible sum or returning a phase obligation is often the correct
mathematical answer, not missing functionality. These acceptance tests target
errors that ordinary complex arithmetic tests would miss, and they \emph{do not
replace proofs of the general algorithms}.

\section{Capability and operation matrix}
\label{cas:sec-matrix}

The following is the union of the two capability matrices in the sources, with
an added column naming which of the three kernels, if any, implements each row.
``Effective'' always includes the stated coefficient and exponent operations.
``Conditional'' means that support, branch, coefficient-algorithm or category
hypotheses must be supplied and checked; it is not a claim of a general
terminating algorithm.

\begingroup\footnotesize
\begin{longtable}{@{}P{0.175\textwidth}P{0.235\textwidth}P{0.295\textwidth}P{0.205\textwidth}@{}}
\toprule
Layer or operation & Exact representation available &
Algorithmic promise and restriction & Implemented in\\
\midrule\endhead
Rational monomial core: $+,\times$ & Polynomial fractions with exact
coefficients & Terminating arithmetic and decidable equality &
\textbf{All three}\\[0.35em]
Inverse & Rational node & Exact for a proved nonzero value &
RationalHahn; HahnRational; SurrealCASCore \emph{on fractions only}\\[0.35em]
Equality and real order & Rational or algebraic data & Total for the rational
core; partial in general; total in supported structured classes including
appropriate transseries zero-test domains & \textbf{All three} (real
comparison)\\[0.35em]
Valuation and leading coefficient & Rational node & Exact finite polynomial
calculation; cannot be total for arbitrary coefficient programs &
\textbf{All three}\\[0.35em]
Standard part & Finite value with known leading data & Exact on finite inputs;
may need a coefficient query plus a finiteness certificate &
RationalHahn; HahnRational. \textbf{Not} SurrealCASCore\\[0.35em]
Conjugation, $\mathrm{Re}$, $\mathrm{Im}$, $|z|^2$ & Complex rational core &
Exact, coefficientwise & HahnRational (all four); SurrealCASCore (conjugation
only); \textbf{not} RationalHahn\\[0.35em]
Rational / fractional exponents & Refined exponent lattice & Exact after
common-denominator refinement & \textbf{SurrealCASCore only}\\[0.35em]
Real closure and complex roots & Finite equations plus ordered root selectors &
Implementable terminating algebraic operations; cost not bounded &
\textbf{None}\\[0.35em]
Modulus and algebraic roots & Selected real/algebraic extension & Exact root
descriptor and branch selection; finite normal form not guaranteed &
\textbf{None}\\[0.35em]
Infinitesimal reciprocal and powers & Rational/root node or positive-support
series & Coefficients effective with finite dependency procedures;
exponent/coefficient promotion allowed & \textbf{None} (series layer)\\[0.35em]
Coefficient at one exponent & Finite expansion or grid & Effective positive
grids and nested supported descriptions; arbitrary supports need extra data &
\textbf{None}\\[0.35em]
All terms below a cut & Compressed exact rational object & A finite list may not
exist even for a valid Hahn series & \textbf{None}\\[0.35em]
Ordinary analytic lifts & Germ descriptor at $c+\eta$ & Exact under branch and
regular-centre hypotheses; coefficients may leave the algebraic core &
\textbf{None}\\[0.35em]
Radius-free Hahn germs & AG's $\Acal_n$ descriptor & Strictly weaker than
common-domain coherence; needs separate support and germ programs &
\textbf{None}\\[0.35em]
Real exp and log & Gonshor-semantic graph & All real inputs for exp, positive
for log; normalization only where the backend supports the scales &
\textbf{None}\\[0.35em]
Canonical complex exp & Strip object on $\HH$ & Imaginary part finite; real
growth may be infinite & \textbf{None}\\[0.35em]
Canonical sine and cosine & Finite-angle or $\mathscr V$ strip object & Real
part finite; arbitrary imaginary growth allowed & \textbf{None}\\[0.35em]
Inverse real trigonometry & Finite-angle or ramified germ data & Natural
real-surreal domains; endpoint branches explicitly selected &
\textbf{None}\\[0.35em]
Infinite circular phases & Named character extension or named $E_\alpha$ policy
& Exact only after a coherent provider is fixed; extraction of $\fin$ and $\ell$
may not be effective & \textbf{None}\\[0.35em]
Polynomial / trig solving & Algebraic or finite cluster objects & Complete over
the effective algebraic core; broader coefficients need effective zero and root
capabilities & \textbf{None}\\[0.35em]
Function differentiation & Germ or coherent derivative descriptor & Scalars
fixed; not a scalar-field derivation & \textbf{None}\\[0.35em]
Scalar derivation & Named differential-field structure & Only supported
identities and summable families; different from formal $D_t$ &
\textbf{None}\\[0.35em]
Root counts and finite schemes & Polynomial algebra, matrices, selected roots &
Initial polynomials and profiles are exact; deformation methods need effective
ordinary division and support & \textbf{None}\\[0.35em]
Integration, residues, periods & Coefficientwise integral descriptor; algebraic
period & Common support, ordinary contour and analytic-domain conditions
retained; no universal fine-local fundamental theorem & \textbf{None}\\[0.35em]
General lazy Hahn equality & Finite program plus semantic certificate & No total
universal procedure, even for $\{0,1\}$ coefficients on $\N$ &
\textbf{None}\\[0.35em]
Numerical display & Coefficient approximation or explicit scale specialization &
Not an ordinary numeric value of an infinite surreal; a specialization is a
simulation & \textbf{None}\\
\bottomrule
\end{longtable}
\endgroup

\section{Conclusions}
\label{cas:sec-conclusions}

\subsection{The representable subsets and their operations}
\label{cas:sub-conclrep}
A substantial exact symbolic surreal CAS is possible without a universal
representation of all surreals. A large subset is available immediately: finite
normal forms over an effective coefficient field, rational expressions in
ordered monomials, and all algebraic extensions of a specified effective ordered
monomial field. The real-closed fields $F$ and their algebraically closed
complexifications $C$ of \cref{cas:thm-realclosure} make this precise --- every
value there has a finite algebraic description even when its Hahn normal form is
infinite. This already gives exact arithmetic, comparisons, polynomial roots,
linear algebra and a large part of finite geometry at infinite and infinitesimal
scales.

Beyond that centre, exact representation remains possible without complete
normalization: certified recurrence-defined series, positive-grid Hahn
generators, regular implicit solutions, ordinary analytic lifts, Hahn-coherent
functions on specified ordinary domains, radius-free germs, and selected
transserial expressions. Exact field arithmetic, real comparison, conjugation,
algebraic root specification, finite polynomial algebra and finite linear
algebra are realistic terminating operations in the effective core. Leading
terms, standard parts, Newton profiles, root-cluster counts and certified series
refinement are realistic when the representation supplies the required valuation
and support services. Finite quotient algebras can often compute residues,
traces and multiplicities without expanding individual roots. Formal and
coherent differentiation, regular implicit solving, positive-support Taylor
substitution, local inverses, preparation and coefficientwise contour operations
can be implemented for suitably described functions. Global real exponentiation
can be admitted with its established surreal semantics; global surcomplex
transcendental operations must additionally state their extension and branch
policies; and integration must name its derivative and function category.

Their exactness means that their \emph{semantics are specified}. It does not
imply that all their identities are decidable or that every requested expansion
is finite. \textbf{There is no largest finite-description class that
simultaneously solves all these problems}; enlarging the expression language
usually changes its closure, equality, order and evaluation guarantees. The
right product is a family of compatible exact domains with explicit
capabilities.

\subsection{Four independent sources of limitation}
\label{cas:sub-fourlimits}
The boundaries have four different origins, and conflating them obscures what
can be fixed and what cannot.
\begin{enumerate}
\item \textbf{Cardinality.} Finite syntax cannot name every real, a fortiori not
every surreal (\cref{cas:prop-countable}). No language extension removes this.
\item \textbf{Undecidability.} Arbitrary computable coefficient streams admit no
universal terminating equality test, no total support-admissibility validator,
and no universal leading-term or sign algorithm in higher rank, even under a
nonzero promise (\cref{cas:thm-zerotest,cas:prop-validation,cas:prop-leading}).
The representation, not the value set, is what carries the guarantee.
\item \textbf{Higher-rank support geometry.} A valid finite-valuation cutoff
request need not have a finite term-list answer, and rank one alone does not
suffice either (\cref{cas:sub-finiteprefix}). This is an output-representation
limitation.
\item \textbf{Analytic category and branch choice.} Fine-local germs,
common-domain coherent data, radius-free germs and contour-integrable data are
different categories; global logarithms, infinite circular phases and global
surcomplex exponentials require declared policies
(\cref{cas:sec-functions,cas:sec-integration,cas:sec-surcomplex}).
\end{enumerate}
None of these makes the implementable subsets uninteresting. Together they
determine a sound software architecture.

\subsection{Attach the necessary information to the object}
\label{cas:sub-attach}
The constructive lesson is uniform across all three sources. A \emph{parent}
identifies arithmetic and order, including coordinate significance. A
\emph{support presentation} identifies coefficient dependencies. A \emph{root
descriptor} identifies a branch. A \emph{derivative} identifies its constants. A
\emph{contour} identifies a chain category and pole-avoidance conditions. A
\emph{global surcomplex function} identifies its extension policy. With those
interfaces the strong formulas of the supplied manuscripts become either genuine
algorithms on effective inputs or clearly labelled conditional constructions;
without them, elegant symbolic notation can conceal undefined operations or
nonterminating searches.

The most useful meaning of ``exact surreal support'' is therefore not a single
unchecked promise. It is a visible contract: an exact denotation, a known
representation class, a declared set of effective operations, certified
remainder information when approximating, and explicit analytic conventions. The
central design principle is not to equate an infinite object with an
approximation, or a mathematical existence theorem with a program: preserve
exact objects, expose certified finite information about them, and make every
change of representation, scale, domain or analytic interpretation explicit.
That principle turns the supplied surcomplex theory into a credible
computer-algebra program rather than a collection of unsafe rewrite rules. The
three delivered prototypes demonstrate the central algebraic point ---
\emph{exact multi-rank arithmetic does not require truncating infinitesimals or
replacing infinity by a numerical limit} --- within three narrow and
separately stated envelopes; the more ambitious layers are specified here as
extensions, with their mathematical and computational obligations made explicit.

\appendix

\section{Source map and boundaries of reliance}
\label{cas:app-audit}

This appendix merges the two source-audit tables from the sources into one, with
an \emph{origin} column naming the manuscript and, where relevant, the revision.
\textbf{References to theorem names are references to the supplied source files,
not claims that those manuscripts implemented the algorithms.} The final column
preserves both tables' ``what is NOT inferred'' discipline.

\begingroup\footnotesize
\begin{longtable}{@{}P{0.135\textwidth}P{0.265\textwidth}P{0.265\textwidth}P{0.25\textwidth}@{}}
\toprule
Origin & Source result or convention & Role in this report &
What is \emph{not} inferred\\
\midrule\endhead
SA (2) and (3) & Strong summability and the Hahn--Neumann positive-monoid lemma
& Semantic certificates for products, inverses and infinitesimal evaluation &
No automatic zero test; no effective support decomposer.\\[0.4em]
SA (2) and (3) & ``Evaluation of arbitrary formal series'' after a sufficiently
small rescaling & Existence of sufficiently small analytic scales
(\cref{cas:sub-smallscale}) & No algorithm extracting a global bound from an
arbitrary program; the scale may leave the workspace.\\[0.4em]
SA (2) and (3) & Hahn-coherent data on a fixed domain, and the composition
theorem & \cref{cas:def-coherent}, \eqref{cas:eq-composition} & No gluing of
arbitrary unrelated fine-local germs; no unconditional global-support
discovery.\\[0.4em]
SA (2) and (3) & Preparation and exact root lifting; implicit simple roots &
\eqref{cas:eq-implicit}; finite clusters; coefficient recursions & No completion
of an arbitrary transfinite recursion by a finite loop.\\[0.4em]
SA (2) and (3) & Contour and moving-pole residue theorems &
\eqref{cas:eq-contour}, \eqref{cas:eq-residue} & No unrestricted
fine-topological Riemann integration; not a full-class Riemann
integral.\\[0.4em]
SA (2) and (3) & Global fine-analytic logarithm; all-scale rigidity &
\cref{cas:sub-logbranches}, \cref{cas:sub-rigidity} & No universal nonexistence
of transcendental functions; not a coherent primitive of $1/z$; no implemented
global atlas.\\[0.4em]
SA (2) and (3) & The $E_\alpha$ phase family with
$E_0(i\omega)=1$, $E_\pi(i\omega)=-1$ & \cref{cas:sub-infinitephases} as an
interface requirement & \textbf{Not a published uniqueness classification}; no
claim that no other convention can be chosen; not automatically
effective.\\[0.4em]
SA (2) and (3) & The fine topology, in which set-indexed convergent nets are
eventually constant & \cref{cas:sub-finetopology} & Truncation sequences are not
fine-convergent; a rank-one valuation topology is a different
structure.\\[0.4em]
TR & Finite-angle circle, polar decomposition, inverse functions &
\cref{cas:sub-finiteangle,cas:sub-inversetrig} & No necessity for infinite
circular angles; extracting $u_0$ still needs a standard-part
capability.\\[0.4em]
TR & Canonical strip functions $E_{\HH}$, $\Sin$, $\Cos$ &
\cref{cas:sub-strips} & No uniquely specified evaluation beyond the strips;
$\sin\omega$ and $E(i\omega)$ are not fixed.\\[0.4em]
TR & Character classification of circular group-law extensions; infinite periods
& \cref{cas:sub-characters}, \cref{cas:prop-periods} & No default value of
$\sin\omega$ forced by the local axioms; \textbf{completeness is claimed only
within TR's framework} and is not extended to SA's $E_\alpha$.\\[0.4em]
TR & Tangency and valuation loss \eqref{cas:eq-acosprecision}; the $2n$ root
bound; the positivity example & \cref{cas:sub-tangency,cas:sub-trigpoly,%
cas:sub-positivity} & Replacing $n$ by $\omega$ is meaningless; ordinary-angle
sampling is unsound.\\[0.4em]
PA & Initial polynomial and its ball counts; Newton profiles; Hensel cluster
lifting & \cref{cas:sub-initial,cas:sub-newton,cas:sub-hensel} & The valuations
must be supplied for general coefficients; no promise to list a
positive-dimensional family; classical Newton--Puiseux is a rank-one
specialization.\\[0.4em]
PA & Root-perturbation bounds and the squarefree/discriminant conditions &
\cref{cas:subsub-univariate} & Not a multivariate statement; the polynomial must
be normalized to the finite-coefficient setting first.\\[0.4em]
PA & Residue functional, trace identity, residue-vs-trace determinants &
\cref{cas:sub-remainders,cas:sub-pairings} & Not a universal rewrite rule; not
an identification with a chosen Grothendieck-residue theory.\\[0.4em]
HC & Fixed-domain coefficient rings; positive-operator inverses; coupled
finite-free systems; $S+E^*$ bounds &
\cref{cas:sub-inverses,cas:sub-coupled} & Do not synthesize finite certificates
or decide identities of arbitrary holomorphic coefficients; the extensions are
explicitly delimited.\\[0.4em]
AG & Radius-free germ algebra $\Acal_n$; evaluation and translation; finite jets
& \cref{cas:sub-radiusfree,cas:sub-jetsgerms} & Not every convergent germ is
already a CAS object; \emph{finite mathematical dependencies need not be
effectively enumerable} --- AG's own warning.\\[0.4em]
AG & Weierstrass preparation and division operators \eqref{cas:eq-weierstrass} &
\cref{cas:sub-weierstrass} & \textbf{Exact formulas, not an implemented
routine}; an arbitrary choice-based splitting is not executable; semidecidable
germ zero-detection may not terminate.\\[0.4em]
AG & Normal forms, multiplication matrices, perturbed residue, trace identity;
length preservation & \cref{cas:sub-normalforms,cas:sub-length} & Requires an
isolated leading complete intersection with $n$ equations in $n$ variables; an
overdetermined system is not covered.\\[0.4em]
AG & Root separation $\kappa_a$ and certified continuation &
\cref{cas:subsub-multivariate} & Conditional on obtaining the exact bounds; the
equality case $\delta=2\kappa_a$ cannot be accepted.\\[0.4em]
CT & Coefficientwise Hahn integration and Stokes; microscopic rescaling &
\cref{cas:sub-coefficientwise,cas:sub-rescaling,cas:sub-stokes} & Not a limit of
fine-topological Riemann sums; not an algorithm deciding admissibility of every
input map; Jordan separation alone gives no integration operation.\\[0.4em]
CT & Separating torus with $\det\mathsf C$; algebraic rational periods &
\cref{cas:sub-torus,cas:sub-periods} & Imported \emph{conditionally} under CT's
finite-algebra and transformation hypotheses; a pole on the contour is outside
the definition.\\[0.4em]
Published literature & Conway, Gonshor, van den Dries--Ehrlich, Poonen,
Perrucci--Roy, van der Hoeven, Berarducci--Mantova, Costin--Ehrlich,
Bagayoko--van der Hoeven, Mantova 2026, Richardson, Chen--Fang--van der Hoeven,
Bournez--Guilmant, Cattani--Dickenstein--Sturmfels & Foundations, exponential
structure, algebraic closure, quantifier elimination, transseries, derivations,
integration, zero tests & See \cref{cas:sec-nonclaims-note}. No claim that one
cited framework equals any other.\\[0.4em]
Official CAS documentation & Wolfram and Sage reference pages, consulted
September 21, 2026 & \cref{cas:sec-host} & A statement about the interfaces
inspected, not an exhaustive survey; these motivate reuse, not a claim that the
facilities already implement surreal semantics.\\[0.4em]
The three source articles & The countability bound, both halting reductions, the
core theorem, the grid lemma in two constructions, the precision rules, the
cutoff obstruction, the proposed architectures, and the executed checks &
Throughout & \textbf{No priority claim} is made for any of these elementary
deductions. The executed checks are not proofs of the general theorems.\\
\bottomrule
\end{longtable}
\endgroup

\subsection{A note on what is not claimed}
\label{cas:sec-nonclaims-note}
Collected for the published literature specifically: hyperseries and the
Bagayoko--van der Hoeven representation theorem are cited for the mathematical
result only, not for a finite universal encoding or an implementation claim, and
the countability argument still applies to them. The Hahn-field
algebraic-closure input (Poonen) is a theorem about the complete mathematical
field, not a terminating algorithm for arbitrary coded series, and the full
abstract Puiseux field still contains noncomputable coefficient sequences.
Mantova's 2026 monotonicity and Taylor-approximation results,
Berarducci--Mantova's derivations and germs, Bournez--Guilmant's stable
set-sized fields, and the elementary-extension theorem for a surreal exponential
structure are all \emph{semantic theorems for specified classes}, not universal
decision algorithms, and a birthday bound is not a finite executable
representation or a complexity estimate. Costin--Ehrlich integration is
established background for suitable analytic and resurgent classes and is
explicitly not identified with the local Taylor-lift interface or with the
manuscripts' coefficientwise contour functional. Richardson's undecidability
results have restrictions that matter and must not be inflated into a blanket
claim about every restricted exponential-constant identity problem. The ISSAC
2026 D-algebraic transseries zero-test holds under its effective-field,
transbasis and solution-selection hypotheses and is not a zero-test for
arbitrary series-producing programs.

\section{Minimal backend contracts and component responsibilities}
\label{cas:app-contracts}

A \textbf{coefficient backend} should specify exact construction; zero and sign
capabilities; arithmetic; conversion to larger coefficient domains; and whether
its ordinary constants admit certified numerical enclosures.

A \textbf{monomial (exponent) backend} should specify equality, multiplication
(addition of exponents), order, rational-root promotion where supported, and the
embedding used to interpret it inside the surreal field --- including its
coordinate significance, which is part of the descriptor and never inferred from
variable names.

A \textbf{support backend} should distinguish a \emph{semantic proof of well
ordering} from \emph{algorithms} for membership, finite coefficient
dependencies, leading-term discovery and cutoff materialization, and it should
declare which of those operations are total.

An \textbf{analytic backend} should specify the ordinary or surreal domain;
branch and extension identifiers; Taylor coefficients; the derivative mode; the
regularity category (fine-local, common-domain coherent, radius-free germ, or
contour-integrable); and the admissible composition rules.

An \textbf{operation result} should preserve its exact origin and its
obligations. If the answer is a truncation it must include a remainder
descriptor and the error calculus used to obtain it. If the answer is
conditional, the condition must remain visible in further transformations. If an
unsupported request has a well-defined semantic node, the system may retain that
node without pretending to have evaluated or normalized it. The component
responsibility table of \cref{cas:sub-components} applies to this appendix in
full.

\section{Notation and semantic conventions}
\label{cas:app-notation}

\begingroup\small
\begin{longtable}{@{}P{0.25\textwidth}P{0.68\textwidth}@{}}
\toprule
Notation & Meaning\\
\midrule\endhead
$\No$, $\SC$ & The surreal class field and its complexification $\No[i]$.\\[0.3em]
$t^\gamma$ & The Conway monomial $\omega^{-\gamma}$, \emph{not} a generic power
operation.\\[0.3em]
$\RA$ & The real algebraic numbers.\\[0.3em]
$\Q^r_{\mathrm{lex}}$, $\Z^d_{\mathrm{lex}}$ & Rational or integral exponent
vectors ordered by their first differing coordinate --- with the coordinate
significance declared by the domain descriptor
(\cref{cas:rem-convention}).\\[0.3em]
$E_r$, $F_r$, $C_r$ & An ordered rational monomial field, its relative real
closure, and the latter's complexification.\\[0.3em]
$v(A)$, $\lc(A)$ & The least Hahn exponent and its coefficient;
$v(0)=+\infty$, used only as a distinguished marker.\\[0.3em]
$|z|$ & The surreal modulus $\sqrt{z\bar z}$, distinct from the
valuation.\\[0.3em]
$\mm$, $\OO$, $\st$ & Infinitesimals, finite elements, and the standard part on
finite elements.\\[0.3em]
$\fin$ & TR's finite-part projection, removing all negative-exponent terms;
additive but \emph{not} multiplicative.\\[0.3em]
$\Pi$, $\Pi\oplus\Z$ & The purely infinite additive class, and the ordered
subring used for the floor \eqref{cas:eq-floor}.\\[0.3em]
$E^*$ & The monoid of finite sums from a positive well-ordered exponent set
$E$.\\[0.3em]
$\Ov{\beta}$ & A remainder with valuation at least $\beta$, including zero; the
\emph{weak} inequality unless a strict one is displayed.\\[0.3em]
$U\sh$, $f\sh$ & Finite surcomplex points with ordinary standard part in $U$,
and the ordinary analytic Taylor lift.\\[0.3em]
$\HH(U)$, $\Acal_n$ & Common-domain Hahn-coherent coefficient data, and AG's
radius-free Hahn-germ algebra (strictly weaker).\\[0.3em]
$D_z$, $D_t$ (or $D_T$), $\partial$ & Coordinate derivative, formal
scale-variable derivative, and a separately selected surreal-field derivation;
three different operators.\\[0.3em]
$E_\alpha$ & SA's explicitly phase-selected global surcomplex exponential;
\emph{not} an unqualified canonical complex exponential.\\[0.3em]
$U_\chi$ & TR's circular group-law extension by a character
$\chi:(\Pi,+)\to(\mathbb T(\No),\cdot)$.\\[0.3em]
$\HH$ (strip), $\mathscr V$ & TR's canonical strips: finite imaginary part, and
finite real part respectively.\\[0.3em]
$I_{a,\rho}(P)$ & The initial polynomial at a centre and a valuation
scale.\\[0.3em]
$\lambda_P(H)$, $\Res_F$ & The top coefficient of $H$ modulo monic $P$, and AG's
perturbed residue functional.\\[0.3em]
$\kappa_a$ & AG's root-separation constant \eqref{cas:eq-kappa}.\\
\bottomrule
\end{longtable}
\endgroup

\paragraph{Per-kernel coordinate conventions, restated.}
\begingroup\small
\begin{longtable}{@{}P{0.20\textwidth}P{0.37\textwidth}P{0.37\textwidth}@{}}
\toprule
Kernel & Variable listing & Comparison\\
\midrule\endhead
\textbf{RationalHahn} & \emph{Least}-significant-valuation first:
$v(t_1)=1$, $v(t_2)=\omega$, \dots & Exponent vectors compared after
\texttt{Reverse}; highest-index coordinate first.\\[0.4em]
\textbf{SurrealCASCore} & \emph{Most}-significant first;
$(q_1,\dots,q_r)\mapsto q_1\omega^{r-1}+\cdots+q_r$ & Plain lex on $\Q^r$; first
coordinate most significant.\\[0.4em]
\textbf{HahnRational} & \emph{Most}-significant first; for \texttt{\{u,t\}},
$v(u)=\{1,0\}$, $v(t)=\{0,1\}$ & Plain lex on $\Z^d$; leftmost coordinate most
significant; a reversed list is a \emph{different parent}.\\
\bottomrule
\end{longtable}
\endgroup

\section{Reproducibility and archive contents}
\label{cas:app-repro}

\subsection{Building this report}
\label{cas:app-build}
This article is standalone LaTeX with an internal bibliography. It uses no
external figures, data files, or bibliography database. From its own directory:
\begin{lstlisting}
latexmk -pdf -interaction=nonstopmode article.tex
\end{lstlisting}
Alternatively run \texttt{pdflatex} until the references and table of contents
stabilize. Standard TeX Live packages and Latin Modern are sufficient.

\subsection{Archive contents}
\label{cas:app-contents}
Beside this article, the directory carries the three source manuscripts and
their READMEs under \path{sources/}, the executable files of all three
prototypes under \path{code/}, and the recorded verification reports under
\path{data/}. \textbf{The files in \path{sources/}, \path{code/} and
\path{data/} ship verbatim as delivered and are not modified here} --- including
the defective \path{RationalHahn.wl} (see \cref{cas:app-defectnote}). The two
manuscripts bundled by one archive are not re-shipped: they are already
committed at
\begin{sloppypar}\noindent
\path{docs/surcomplex/trigonometry/sources/18-canonical-angles-oscillation.tex}
and
\path{docs/surcomplex/analysis/sources/07-infinitesimal-calculus-contour-theory.tex}
(see \cref{cas:sub-provenance}).
\end{sloppypar}

\subsection{Three separate run instructions}
\label{cas:app-runs}

\paragraph{RationalHahn --- requires a one-character repair first.}
\label{cas:app-defectnote}
The archive's own instruction is
\begin{lstlisting}
wolframscript -file code/RunChecks.wl
\end{lstlisting}
\textbf{As delivered this does not work}: the package does not parse, and the
driver reports \texttt{Checks -> 34, Passed -> 12} --- the twelve
kernel-independent identities --- with the kernel's own
\texttt{Syntax::sntue} error. Apply the one-character repair of
\cref{cas:rem-defect} to a \emph{copy} of the package (close the \texttt{If[} in
\texttt{RHInverse}: write \texttt{...a[[1]]]];} for \texttt{...a[[1]]];}), place
that copy beside \path{Checks.wl} and \path{RunChecks.wl} under the names the
driver expects, and the driver then returns \texttt{Checks -> 34, Passed -> 34,
Failures -> \{\}} with no warnings. Run it in a fresh kernel or a disposable
context: \path{Checks.wl} uses named global variables for readability, and the
package rejects variables that have already evaluated to nonsymbolic values.

\paragraph{SurrealCASCore.}
\begin{lstlisting}
wolframscript -file TestCore.wls
python -m pip install -r requirements.txt
python verify_examples.py
\end{lstlisting}
The harness locates the package relative to its own path, so the working
directory need not be the package directory; the Python script regenerates its
report beside itself and fails if an assertion fails. Expect $110/110$ and
$69/69$.

\paragraph{HahnRational.}
\begin{lstlisting}
wolframscript -file code/test_HahnRational.wl
python code/verify_examples.py
\end{lstlisting}
The test file loads the package relative to its own path when the package is not
already loaded. Expect $88/88$ and $323$ exact Python checks. The short demo
\path{examples.wl} prints four of the displayed session results.

\subsection{Combined statement}
\label{cas:app-combined}
One kernel release (15.0.1 of July 2, 2026, recorded on Linux x86-64 by the
archives and reproduced on Windows x86-64 here); one date (September 21, 2026);
\textbf{no local installation exercised by any of the three original runs};
\textbf{counts not additive}, either within an article or across articles; no
formal verification of any general theorem; and \textbf{no formalization
anywhere} --- no Lean file, no blueprint, no formalization audit, no pinned git
commit, Lean toolchain or mathlib revision in any of the three packages. The
executed checks are exact finite algebra and formal-series checks, and a
successful driver run does not upgrade them into proofs of the infinite
constructions.

\clearpage
\begin{thebibliography}{99}
\addcontentsline{toc}{section}{References}
\small
\setlength{\itemsep}{0.35em}

\bibitem{cas:SourceAnalysisTwo}
\emph{Surcomplex Analysis: Infinitesimal Calculus, Hahn-Coherent Holomorphy,
and Contour Theory}. Supplied research manuscript, September 21, 2026; file
\path{surcomplex_analysis(2).tex}. \textbf{Revision (2)}, supplied to the
capability-tiers article, which bundled it unchanged. Identical, after CRLF
normalization, to
\path{docs/surcomplex/analysis/sources/07-infinitesimal-calculus-contour-theory.tex}
in this repository. Source labels relied upon include \texttt{def:strong},
\texttt{thm:smallscale}, \texttt{def:coherent}, \texttt{thm:preparation},
\texttt{thm:globallog} and \texttt{thm:rigidity}. No refereeing or proof
verification is inferred.

\bibitem{cas:SourceAnalysisThree}
\emph{Surcomplex Analysis: Infinitesimal Calculus, Hahn-Coherent Holomorphy,
and Contour Theory}. Supplied research manuscript, September 21, 2026; file
\path{surcomplex_analysis(3).tex}. \textbf{Revision (3)}, supplied to the
layered-representations and Gaussian-prototype articles, which cited rather than
bundled it. Its separate notions of local germs, coherent data, coefficientwise
contours and all-scale functions are retained here. No refereeing or proof
verification is inferred.

\bibitem{cas:SourceTrig}
\emph{Trigonometry on the Surcomplex Plane: Canonical Angles, Infinite
Triangles, Infinitesimal Contact, and Hahn-Analytic Oscillation}. Supplied
research manuscript, \path{surcomplex_trigonometry.tex}, September 21, 2026.
Byte-identical to
\path{docs/surcomplex/trigonometry/sources/18-canonical-angles-oscillation.tex}
in this repository. Primary basis for the finite-angle, strip, branch,
root-count and infinite-phase conventions; source labels relied upon include
\texttt{thm:polar}, \texttt{thm:tangency}, \texttt{thm:stripexp},
\texttt{thm:polyroots}, \texttt{thm:characters} and
\texttt{thm:infiniteperiods}. Two further archives mentioned inside this
manuscript were not supplied and are not treated as reviewed sources.

\bibitem{cas:SourcePolynomial}
\emph{Surcomplex Polynomial Algebra: Order Geometry, Newton Profiles, and
Scale-Resolved Stability}. Supplied research manuscript, September 21, 2026;
file \path{surcomplex_polynomial_algebra(2).tex}. The workspaces,
initial-polynomial formulas, Newton profiles, Hensel lifting, perturbation
bounds and finite residue algebra are used with their stated hypotheses. A
fourth, global-divisor manuscript cited inside this file was not supplied and is
not used here as an independently inspected source.

\bibitem{cas:SourceResearch}
\emph{Hartogs Coherence, Finite Maps, and Residue Duality over the Surcomplex
Numbers}. Supplied research manuscript, September 21, 2026; file
\path{surcomplex_research.tex}. In particular the positive-operator calculus,
fixed-domain division, coordinate-power perturbations, coupled finite-free
systems and finite residue pairings.

\bibitem{cas:SourceAG}
\emph{Infinitesimal Analytic Geometry over the Surcomplex Numbers: Radius-Free
Hahn Germs, a Nullstellensatz, and Conservation of Singular Zeros}. Supplied
research manuscript, September 21, 2026; file
\path{surcomplex_analytic_geometry.tex}. Especially its Sections 2--4, 6--12 and
Appendix B. Its general theorems are not independently machine-verified here.

\bibitem{cas:SourceCT}
\emph{Jordan Separation, Cauchy Theory, and Stokes Formulae on the Surcomplex
Plane: A Category-Sensitive Contour Calculus}. Supplied research continuation,
September 21, 2026; file \path{surcomplex_contours(2).tex}. Especially its
Sections 2, 4--10 and Appendices A--C. The stated category hypotheses are
retained.

\bibitem{cas:Conway}
John H. Conway. \emph{On Numbers and Games}. Academic Press, 1976; second
edition, A K Peters, 2001.

\bibitem{cas:Gonshor}
Harry Gonshor. \emph{An Introduction to the Theory of Surreal Numbers}. London
Mathematical Society Lecture Note Series \textbf{110}, Cambridge University
Press, 1986.
\href{https://doi.org/10.1017/CBO9780511629143}{doi:10.1017/CBO9780511629143}.

\bibitem{cas:vdDE}
Lou van den Dries and Philip Ehrlich. ``Fields of surreal numbers and
exponentiation.'' \emph{Fundamenta Mathematicae} \textbf{167} (2001), 173--188.
\href{https://doi.org/10.4064/fm167-2-3}{doi:10.4064/fm167-2-3}. Erratum,
\textbf{168} (2001), 295--297. Used for the normal-form and
exponential-structure background, not for ordinal-length estimates affected by
the erratum.

\bibitem{cas:Poonen}
Bjorn Poonen. ``Maximally complete fields.'' \emph{L'Enseignement
Math\'ematique} (2) \textbf{39} (1993), 87--106.
\href{https://math.mit.edu/~poonen/papers/amsval.pdf}{Author-hosted text}.
Corollary 4 supplies the algebraic-closure input for divisible Hahn workspaces;
Corollary 7 the characteristic-zero Puiseux extension in rank one.

\bibitem{cas:PerrucciRoy}
Daniel Perrucci and Marie-Fran\c{c}oise Roy. ``Elementary recursive quantifier
elimination based on Thom encoding and sign determination.'' \emph{Annals of
Pure and Applied Logic} \textbf{168} (2017), no.~8, 1588--1604.
\href{https://doi.org/10.1016/j.apal.2017.03.001}{doi:10.1016/j.apal.2017.03.001};
\href{https://arxiv.org/abs/1609.02879}{arXiv:1609.02879v2}.

\bibitem{cas:Richardson}
Daniel Richardson. ``Some undecidable problems involving elementary functions of
a real variable.'' \emph{The Journal of Symbolic Logic} \textbf{33} (1968),
no.~4, 514--520. \href{https://doi.org/10.2307/2271358}{doi:10.2307/2271358}.
The restrictions of this theorem matter; it is not used to assert undecidability
of every restricted exponential language.

\bibitem{cas:vdH}
Joris van der Hoeven. \emph{Transseries and Real Differential Algebra}. Lecture
Notes in Mathematics \textbf{1888}, Springer, 2006.
\href{https://doi.org/10.1007/3-540-35590-1}{doi:10.1007/3-540-35590-1};
\href{https://www.texmacs.org/joris/ln/ln.html}{author-hosted text}. Especially
the chapters on grid-based series, Newton polygons and operations on
transseries. None of the delivered prototypes implements that theory.

\bibitem{cas:BMDerivations}
Alessandro Berarducci and Vincenzo Mantova. ``Surreal numbers, derivations and
transseries.'' \emph{Journal of the European Mathematical Society} \textbf{20}
(2018), 339--390.
\href{https://doi.org/10.4171/JEMS/769}{doi:10.4171/JEMS/769};
\href{https://arxiv.org/abs/1503.00315}{arXiv:1503.00315}.

\bibitem{cas:BMFunctions}
Alessandro Berarducci and Vincenzo Mantova. ``Transseries as germs of surreal
functions.'' \emph{Transactions of the American Mathematical Society}
\textbf{371} (2019), 3549--3592.
\href{https://doi.org/10.1090/tran/7428}{doi:10.1090/tran/7428};
\href{https://arxiv.org/abs/1703.01995}{arXiv:1703.01995v3}.

\bibitem{cas:CostinEhrlich}
Ovidiu Costin and Philip Ehrlich. ``Integration on the surreals.''
\emph{Advances in Mathematics} \textbf{452} (2024), article 109823.
\href{https://doi.org/10.1016/j.aim.2024.109823}{doi:10.1016/j.aim.2024.109823};
\href{https://arxiv.org/abs/2208.14331}{arXiv:2208.14331v5}. A distinct
constructive extension and integration framework with its own specified function
classes; \emph{not} identified here with the supplied manuscripts'
coefficientwise contour functional or with the local Taylor-lift interface.

\bibitem{cas:Hyperseries}
Vincent Bagayoko and Joris van der Hoeven. ``Surreal numbers as hyperseries.''
\href{https://arxiv.org/abs/2310.14879}{arXiv:2310.14879v1}, October 23, 2023.
Cited for its mathematical representation theorem, not for a finite universal
encoding or an implementation claim.

\bibitem{cas:Mantova2026}
Vincenzo Mantova. ``Monotonicity and a Taylor approximation theorem for
transseries.'' \href{https://arxiv.org/abs/2601.07747}{arXiv:2601.07747v2},
revised May 9, 2026. Preprint record consulted September 21, 2026. A semantic
theorem for specified transserial classes.

\bibitem{cas:CFH}
Shaoshi Chen, Hanqian Fang, and Joris van der Hoeven. ``A Zero-Test for
D-Algebraic Transseries.'' In \emph{Proceedings of ISSAC 2026}, ACM, 2026,
pp.~105--113. Published July 12, 2026.
\href{https://doi.org/10.1145/3815436.3815443}{doi:10.1145/3815436.3815443};
\href{https://arxiv.org/abs/2602.01188}{arXiv:2602.01188};
\href{https://www.texmacs.org/joris/zttr/zttr.html}{author-hosted text}.
Especially Sections 2.4--2.5 and 7, including Theorem 18. \textbf{The only
positive general decidability result carried here beyond the algebraic core},
and only under its stated hypotheses.

\bibitem{cas:BG}
Olivier Bournez and Quentin Guilmant. ``Surreal fields stable under exponential,
logarithmic, derivative and anti-derivative functions.'' Preprint, 2022.
\href{https://arxiv.org/abs/2211.08396}{arXiv:2211.08396}.

\bibitem{cas:CDS}
Eduardo Cattani, Alicia Dickenstein, and Bernd Sturmfels. ``Computing
multidimensional residues.'' In L.~Gonz\'alez-Vega and T.~Recio (eds.),
\emph{Algorithms in Algebraic Geometry and Applications}, Progress in
Mathematics \textbf{143}, Birkh\"auser, 1996, pp.~135--164.
\href{https://doi.org/10.1007/978-3-0348-9104-2_8}{doi:10.1007/978-3-0348-9104-2\_8};
\href{https://arxiv.org/abs/alg-geom/9404011}{arXiv:alg-geom/9404011}.

\bibitem{cas:WLRoot}
Wolfram Research. \emph{Root}. Wolfram Language documentation.
\url{https://reference.wolfram.com/language/ref/Root.html}. Consulted
September 21, 2026. The consulted documentation includes general-equation
neighbourhood representations as well as polynomial indexing.

\bibitem{cas:WLAlgebraic}
Wolfram Research. \emph{AlgebraicNumber}; \emph{Algebraic Number Fields}.
Wolfram Language documentation.
\url{https://reference.wolfram.com/language/ref/AlgebraicNumber.html};
\url{https://reference.wolfram.com/language/tutorial/AlgebraicNumberFields.html}.
Consulted September 21, 2026.

\bibitem{cas:WLRootReduce}
Wolfram Research. \emph{RootReduce}. Wolfram Language documentation.
\url{https://reference.wolfram.com/language/ref/RootReduce.html}. Consulted
September 21, 2026.

\bibitem{cas:WLSeriesData}
Wolfram Research. \emph{SeriesData}. Wolfram Language documentation.
\url{https://reference.wolfram.com/language/ref/SeriesData.html}. Consulted
September 21, 2026.

\bibitem{cas:WLSeriesCoefficient}
Wolfram Research. \emph{SeriesCoefficient}. Wolfram Language documentation.
\url{https://reference.wolfram.com/language/ref/SeriesCoefficient.html}.
Consulted September 21, 2026.

\bibitem{cas:WLSeries}
Wolfram Research. \emph{Series}. Wolfram Language documentation.
\url{https://reference.wolfram.com/language/ref/Series.html}. Consulted
September 21, 2026.

\bibitem{cas:WLInverseSeries}
Wolfram Research. \emph{InverseSeries}. Wolfram Language documentation.
\url{https://reference.wolfram.com/language/ref/InverseSeries.html}. Consulted
September 21, 2026.

\bibitem{cas:WLAsymptotic}
Wolfram Research. \emph{Asymptotic}. Wolfram Language documentation.
\url{https://reference.wolfram.com/language/ref/Asymptotic.html}. Consulted
September 21, 2026.

\bibitem{cas:WLAsymptoticSolve}
Wolfram Research. \emph{AsymptoticSolve}. Wolfram Language documentation.
\url{https://reference.wolfram.com/language/ref/AsymptoticSolve.html}.
Consulted September 21, 2026.

\bibitem{cas:WLDifferentialRoot}
Wolfram Research. \emph{DifferentialRoot}. Wolfram Language documentation.
\url{https://reference.wolfram.com/language/ref/DifferentialRoot.html}.
Consulted September 21, 2026.

\bibitem{cas:WLUpValues}
Wolfram Research. \emph{UpValues}. Wolfram Language documentation.
\url{https://reference.wolfram.com/language/ref/UpValues.html}. Consulted
September 21, 2026.

\bibitem{cas:WLInactive}
Wolfram Research. \emph{Inactive}. Wolfram Language documentation.
\url{https://reference.wolfram.com/language/ref/Inactive.html}. Consulted
September 21, 2026.

\bibitem{cas:WLNumericQ}
Wolfram Research. \emph{NumericQ}. Wolfram Language documentation.
\url{https://reference.wolfram.com/language/ref/NumericQ.html}. Consulted
September 21, 2026.

\bibitem{cas:WLPossibleZeroQ}
Wolfram Research. \emph{PossibleZeroQ}. Wolfram Language documentation.
\url{https://reference.wolfram.com/language/ref/PossibleZeroQ.html}. Consulted
September 21, 2026. Documented as a heuristic facility.

\bibitem{cas:WLInfinity}
Wolfram Research. \emph{Infinity}. Wolfram Language documentation.
\url{https://reference.wolfram.com/language/ref/Infinity.html}. Consulted
September 21, 2026.

\bibitem{cas:WLSpecial}
Wolfram Research. \emph{Mathematical Functions: Working with Special
Functions}. Wolfram Language documentation.
\url{https://reference.wolfram.com/language/tutorial/MathematicalFunctions}.
Consulted September 21, 2026.

\bibitem{cas:SageLazy}
The Sage Developers. \emph{Lazy Series}. SageMath Reference Manual.
\url{https://doc.sagemath.org/html/en/reference/power_series/sage/rings/lazy_series.html}.
Consulted September 21, 2026.

\bibitem{cas:SageAsymptotic}
The Sage Developers. \emph{Asymptotic Ring}. SageMath Reference Manual.
\url{https://doc.sagemath.org/html/en/reference/asymptotic/sage/rings/asymptotic/asymptotic_ring.html}.
Consulted September 21, 2026.

\end{thebibliography}
\end{document}
